\documentclass[11pt,a4paper,oneside]{book}

% ============================================================================
% PACKAGES
% ============================================================================
\usepackage[utf8]{inputenc}
\usepackage[T1]{fontenc}
\usepackage[french]{babel}
\usepackage{lmodern}
\usepackage{microtype}
\usepackage[top=2.5cm,bottom=2.5cm,left=2.5cm,right=2.5cm,headheight=15pt]{geometry}

\usepackage{amsmath,amsfonts,amssymb,amsthm}
\usepackage{mathtools}
\usepackage{gensymb}
\usepackage{icomma}
\usepackage{cancel}
\usepackage{multicol}
\usepackage{enumitem}
\usepackage{booktabs}
\usepackage{array}
\usepackage{graphicx}
\usepackage{xcolor}
\usepackage{tikz}
\usepackage{pgfplots}
\pgfplotsset{compat=1.18}
\usetikzlibrary{angles,quotes,calc,positioning,arrows.meta,patterns,decorations.pathreplacing,babel}

\usepackage[most]{tcolorbox}
\usepackage{fancyhdr}
\usepackage{titlesec}
\usepackage{hyperref}

% ============================================================================
% COULEURS & STYLE UCLOUVAIN
% ============================================================================
\definecolor{uclblue}{RGB}{0, 53, 107}
\definecolor{ucllightblue}{RGB}{235, 243, 250}
\definecolor{uclaccent}{RGB}{0, 114, 188}
\definecolor{defgreen}{RGB}{20, 110, 60}
\definecolor{deflightgreen}{RGB}{238, 248, 242}
\definecolor{exorange}{RGB}{180, 85, 0}
\definecolor{exlightorange}{RGB}{253, 246, 238}
\definecolor{warnred}{RGB}{170, 30, 30}
\definecolor{warnlightred}{RGB}{253, 240, 240}
\definecolor{darkgray}{RGB}{60, 60, 60}

\hypersetup{
    colorlinks=true,
    linkcolor=uclblue,
    urlcolor=uclaccent,
    citecolor=uclblue,
    pdftitle={Syllabus d'Algèbre - Examen d'entrée Ingénieur Civil UCLouvain},
    pdfauthor={Laurent Francis et Charles Pecheur (Prompt Engineers) & IA}
}

% ============================================================================
% BOÎTES PÉDAGOGIQUES (TCOLORBOX)
% ============================================================================
\newtcolorbox{theoreme}[1][]{
    enhanced,
    breakable,
    colback=ucllightblue,
    colframe=uclblue,
    fonttitle=\bfseries,
    title={#1},
    arc=3pt,
    boxrule=0.8pt,
    left=8pt,right=8pt,top=6pt,bottom=6pt
}

\newtcolorbox{definition}[1][]{
    enhanced,
    breakable,
    colback=deflightgreen,
    colframe=defgreen,
    fonttitle=\bfseries,
    title={#1},
    arc=3pt,
    boxrule=0.8pt,
    left=8pt,right=8pt,top=6pt,bottom=6pt
}

\newtcolorbox{exemple}[1][]{
    enhanced,
    breakable,
    colback=exlightorange,
    colframe=exorange,
    fonttitle=\bfseries,
    title={#1},
    arc=3pt,
    boxrule=0.8pt,
    left=8pt,right=8pt,top=6pt,bottom=6pt
}

\newtcolorbox{attention}[1][Attention]{
    enhanced,
    breakable,
    colback=warnlightred,
    colframe=warnred,
    fonttitle=\bfseries,
    title={#1},
    arc=3pt,
    boxrule=0.8pt,
    left=8pt,right=8pt,top=6pt,bottom=6pt
}

\newtcolorbox{astuce}[1][Astuce \& Méthode]{
    enhanced,
    breakable,
    colback=gray!8,
    colframe=darkgray,
    fonttitle=\bfseries,
    title={#1},
    arc=3pt,
    boxrule=0.8pt,
    left=8pt,right=8pt,top=6pt,bottom=6pt
}

% ============================================================================
% MISE EN PAGE & EN-TÊTES
% ============================================================================
\setlength{\parindent}{0pt}
\setlength{\parskip}{0.45em}
\setlist[itemize]{topsep=3pt,itemsep=2pt}
\setlist[enumerate]{topsep=4pt,itemsep=4pt}

\pagestyle{fancy}
\fancyhf{}
\fancyhead[L]{\small\textcolor{uclblue}{\nouppercase{\leftmark}}}
\fancyhead[R]{\small\textcolor{darkgray}{Algèbre -- UCLouvain}}
\fancyfoot[C]{\thepage}
\renewcommand{\headrulewidth}{0.4pt}

\titleformat{\chapter}[display]
  {\normalfont\huge\bfseries\color{uclblue}}
  {\filright\large\textsc{\chaptertitlename}\ \thechapter}
  {1ex}
  {\titlerule\vspace{1.2ex}\filright}
  [\vspace{1ex}\titlerule]

\titleformat{\section}
  {\normalfont\Large\bfseries\color{uclblue}}
  {\thesection}{0.8em}{}

\titleformat{\subsection}
  {\normalfont\large\bfseries\color{uclaccent}}
  {\thesubsection}{0.7em}{}

% ============================================================================
% COMMANDES MATHÉMATIQUES
% ============================================================================
\newcommand{\R}{\mathbb{R}}
\newcommand{\Z}{\mathbb{Z}}
\newcommand{\N}{\mathbb{N}}
\newcommand{\Q}{\mathbb{Q}}
\newcommand{\C}{\mathbb{C}}
\newcommand{\D}{\mathbb{D}}
\DeclareMathOperator{\dom}{dom}
\DeclareMathOperator{\Reop}{Re}
\DeclareMathOperator{\Imop}{Im}
\DeclareMathOperator{\cotg}{cotg}
\DeclareMathOperator{\tg}{tg}

% ============================================================================
% DOCUMENT
% ============================================================================
\begin{document}

% ----------------------------------------------------------------------------
% PAGE DE TITRE
% ----------------------------------------------------------------------------
\begin{titlepage}
    \centering
    \vspace*{0.8cm}
    {\Large \textsc{Université catholique de Louvain (UCLouvain)}}\\[0.3cm]
    {\large \textsc{École Polytechnique de Louvain (EPL)}}\\[1.3cm]

    \begin{tcolorbox}[colback=ucllightblue,colframe=uclblue,arc=5pt,boxrule=1.5pt,top=15pt,bottom=15pt]
        \centering
        {\huge\bfseries\color{uclblue} Syllabus d'Algèbre}\\[0.4cm]
        {\Large\color{uclaccent} Préparation à l'examen d'admission\\ aux études d'ingénieur civil}
    \end{tcolorbox}

    \vspace{0.7cm}

    % Illustration centrale : Racines complexes dans C et parabole paramétrique
    \begin{tikzpicture}[scale=1.05, >=Stealth]
        % Panneau gauche : Racines 4-ièmes dans le plan complexe de Gauss
        \begin{scope}[shift={(-3.5,0)}, scale=1.25]
            \draw[step=0.5, gray!20, very thin] (-1.9,-1.9) grid (1.9,1.9);
            \draw[->, thick, darkgray] (-2.0,0) -- (2.1,0) node[right] {\small $\Reop(z)$};
            \draw[->, thick, darkgray] (0,-2.0) -- (0,2.1) node[above] {\small $\Imop(z)$};
            \draw[thick, uclblue, dashed] (0,0) circle (1.414);
            
            \def\r{1.414}
            \coordinate (Z0) at (30:\r);
            \coordinate (Z1) at (120:\r);
            \coordinate (Z2) at (210:\r);
            \coordinate (Z3) at (300:\r);

            \fill[uclaccent!12] (Z0) -- (Z1) -- (Z2) -- (Z3) -- cycle;
            \draw[thick, uclaccent] (Z0) -- (Z1) -- (Z2) -- (Z3) -- cycle;
            \draw[thick, exorange, ->] (0,0) -- (Z0) node[midway, below right=-2pt] {\footnotesize $\rho$};
            \draw[uclblue, thick] (0.45,0) arc (0:30:0.45) node[midway, right=1pt] {\footnotesize $\theta_0$};

            \fill[warnred] (Z0) circle (2pt) node[above right] {\footnotesize $z_0 = \rho e^{i\theta_0}$};
            \fill[warnred] (Z1) circle (2pt) node[above left] {\footnotesize $z_1$};
            \fill[warnred] (Z2) circle (2pt) node[below left] {\footnotesize $z_2$};
            \fill[warnred] (Z3) circle (2pt) node[below right] {\footnotesize $z_3$};
            \node[uclblue, font=\small\bfseries] at (0,-2.35) {Plan complexe $\C$ ($z^n = w$)};
        \end{scope}

        % Panneau droit : Famille de paraboles paramétriques P_m(x) et signes
        \begin{scope}[shift={(3.3,0)}, scale=1.15]
            \draw[step=0.5, gray!20, very thin] (-2.1,-1.8) grid (2.3,2.2);
            \draw[->, thick, darkgray] (-2.2,0) -- (2.4,0) node[right] {\small $x$};
            \draw[->, thick, darkgray] (0,-1.9) -- (0,2.3) node[above] {\small $P_m(x)$};

            % Zone positive
            \fill[defgreen!15] (-2.0,0) -- (-1.1,0) -- plot[domain=-2.0:-1.1, samples=40] (\x, {0.55*(\x-0.3)^2 - 1.078}) -- cycle;
            \fill[defgreen!15] (1.7,0) -- (2.15,0) -- plot[domain=1.7:2.15, samples=40] (\x, {0.55*(\x-0.3)^2 - 1.078}) -- cycle;

            % Courbe Delta > 0
            \draw[very thick, uclblue, domain=-2.0:2.2, samples=80] plot (\x, {0.55*(\x-0.3)^2 - 1.078});
            % Courbe Delta < 0
            \draw[thick, exorange, dashed, domain=-1.6:2.1, samples=60] plot (\x, {0.55*(\x-0.2)^2 + 0.45});

            \fill[warnred] (-1.1,0) circle (2.2pt) node[below left] {\footnotesize $x_1(m)$};
            \fill[warnred] (1.7,0) circle (2.2pt) node[below right] {\footnotesize $x_2(m)$};

            \node[defgreen, font=\small\bfseries] at (-1.65,0.45) {$+$};
            \node[warnred, font=\small\bfseries] at (0.3,-0.55) {$-$};
            \node[defgreen, font=\small\bfseries] at (1.98,0.45) {$+$};
            \node[exorange] at (0.3,1.75) {\footnotesize $\Delta(m) < 0$};
            \node[uclblue] at (1.15,-1.35) {\footnotesize $\Delta(m) > 0$};
            \node[uclblue, font=\small\bfseries] at (0.1,-2.55) {Discussion paramétrique ($m \in \R$)};
        \end{scope}
    \end{tikzpicture}

    \vfill

    {\large \textbf{Direction artistique \& \emph{« Prompt Engineering »} :}}\\[0.2cm]
    {\Large \textbf{L'équipe algèbre de l'examen d'entrée de l'EPL}}\\[0.15cm]
    {\large \textbf{Laurent Fran\c{c}is} et \textbf{Charles Pecheur}}\\[0.3cm]
    {\normalsize \emph{Rédigé à la sueur des GPU par l'Intelligence Artificielle,\\ sous la supervision bienveillante (et quelques prompts bien sentis) de ses humains.}}

    \vspace{0.5cm}
    {\small \textbf{Sources fournies à l'IA :}\\
    Le référentiel de mathématiques de 5\ieme{} et 6\ieme{} secondaire de la Fédération Wallonie-Bruxelles\\
    ainsi que l'intégralité des annales et solutionnaires officiels d'Algèbre de l'UCLouvain (2021--2026).}\\[0.4cm]
    {\normalsize Édition 2026--2027}
\end{titlepage}

% ----------------------------------------------------------------------------
% AVANT-PROPOS ET CONSEILS MÉTHODOLOGIQUES
% ----------------------------------------------------------------------------
\chapter*{Avant-propos et conseils méthodologiques}
\addcontentsline{toc}{chapter}{Avant-propos et conseils méthodologiques}
\markboth{Avant-propos}{Avant-propos}

Ce syllabus a été conçu pour accompagner les étudiantes et étudiants qui préparent l'épreuve d'\textbf{algèbre} de l'examen spécial d'admission aux études d'ingénieur civil à l'\textbf{UCLouvain} (École Polytechnique de Louvain -- EPL). Il est pensé pour être pédagogique, progressif, rigoureux et directement accessible à toute personne ayant suivi le cours de mathématiques de 5\ieme{} et 6\ieme{} secondaire en Communauté française de Belgique (Fédération Wallonie-Bruxelles, en particulier l'option mathématiques 6h/semaine, ou 4h avec un investissement personnel complémentaire).

\section*{Origine, auteurs et sources du document}

\begin{theoreme}[« Prompt Engineers », IA et sources du document]
\begin{itemize}
    \item \textbf{\emph{Prompt Engineers} en chef :} L'équipe algèbre de l'examen d'entrée à l'UCLouvain (EPL), à savoir les professeurs \textbf{Laurent Fran\c{c}is} et \textbf{Charles Pecheur}, qui ont fourni les annales, solutionnaires et directives pédagogiques (et dont les prénoms illustrent avec humour la plupart des problèmes de la Question~3 !).
    \item \textbf{Auteur exécutant (Intelligence Artificielle) :} Dans le même esprit que les syllabus de Trigonométrie/Géométrie et d'Analyse, la rédaction pédagogique détaillée, la structuration \LaTeX{}, le tracé des figures TikZ et la mise au net des corrigés ont été réalisés avec l'aide de l'intelligence artificielle.
    \item \textbf{Annales et solutionnaires officiels d'Algèbre (UCLouvain 2021--2026) :} Ce syllabus intègre et explique en détail l'intégralité des questions posées aux sessions de \textbf{juillet et septembre de 2021 à 2026} à l'examen d'admission de l'UCLouvain. Nous tenons à remercier chaleureusement l'équipe des enseignantes et enseignants responsables de l'épreuve d'Algèbre à l'EPL pour la clarté et la richesse de leurs questionnaires et solutionnaires officiels (où l'on retrouve d'ailleurs régulièrement les aventures mathématiques de \emph{Charles, Laurent, Aude, Sophie} et du groupe de rock \emph{The EPLighters}).
\end{itemize}
\end{theoreme}

\section*{Structure type de l'épreuve d'Algèbre à l'UCLouvain}

L'épreuve d'Algèbre de l'examen d'admission à l'EPL dure \textbf{2h30}, se déroule \textbf{sans calculatrice} (ni téléphone ou montre connectée) et comporte systématiquement \textbf{3 grandes questions} dont le format est remarquablement constant d'année en année :
\begin{enumerate}
    \item \textbf{Question 1 : Sept questions courtes indépendantes (1.1 à 1.7) à réponse directe en cadre (Chapitres 1 à 6)} :
    Pour cette première question, \textbf{seule la réponse finale inscrite dans le cadre compte} (aucune justification n'est demandée sur la copie). Ces 7 sous-questions balaient systématiquement les thèmes suivants :
    \begin{itemize}
        \item \textbf{Polynômes (Chapitre 1) :} division euclidienne $P(x) = D(x)Q(x) + R(x)$, théorème du reste, schéma de Horner, factorisation ou détermination de coefficients $(a,b)$ pour qu'un polynôme comme $x^4+ax+b$ soit divisible par un trinôme du second degré.
        \item \textbf{Domaines de définition et équations irrationnelles (Chapitre 2) :} domaine d'existence de fonctions combinant racines carrées et logarithmes comme $f(x) = \frac{1}{\sqrt{\ln(x^2-1)}}$ ou $\sqrt{\ln\left(\frac{x-1}{x+2}\right)}$, et équations avec radicaux ($x = \sqrt{x+1}$, $\sqrt{1-x^2} = x-1$).
        \item \textbf{Exponentielles et logarithmes (Chapitre 3) :} résolution d'équations exponentielles ($16^x - 4^x = 2$, $2^x + 2^{3-x} = 6$, $9^x + 1 = 10\cdot 3^{x-1}$) et d'équations logarithmiques en bases variées ($\log_{10}$, $\ln$, $\log_2$, $\log_4$, $\log_5$).
        \item \textbf{Nombres complexes dans $\C$ (Chapitre 4) :} forme cartésienne $a+bi$, module $\rho = |z|$ et argument $\theta = \arg(z)$ de quotients ou puissances ($(1+i)^{10}$, $(1+\sqrt{3}i)^4$, $\frac{-6i}{1+i}$), racines $n$-ièmes ($z^3+8i=0$, $z^4+2-2\sqrt{3}i=0$) ou équations avec le conjugué $z^*$.
        \item \textbf{Binôme de Newton (Chapitre 5) :} calcul d'un coefficient précis dans $(ax+b)^n$, simplification de $(x+1)^n \pm (x-1)^n$, ou recherche du terme constant (d'ordre nul en $x^0$) dans $\left(ax^p + \frac{b}{x}\right)^n$.
        \item \textbf{Analyse combinatoire et probabilités (Chapitre 5) :} dénombrement avec contraintes (livret recto-verso, cadenas, choix de livres ou de questions d'examen) et calcul de probabilités sous forme de fraction irréductible (lancers de dés : full, paire ; tirages de cartes, d'urnes ou de chaussettes ; files d'attente).
        \item \textbf{Suites arithmétiques et géométriques (Chapitre 6) :} calcul du premier terme $u_1$, de la raison ou du nombre de termes à partir de la somme $S_n$ (collection de pierres, soirée chantante, échelle de menuisier, tronçons de corde, prix d'un concours).
    \end{itemize}
    \item \textbf{Question 2 : Résolution et discussion complète d'une inéquation (ou équation) paramétrique dans $\R$ (Chapitre 7)} :
    Cette question demande une \textbf{rédaction complète et rigoureuse}. Il s'agit de résoudre dans $\R$ une inéquation dépendant d'un paramètre réel $m$ ou $a$ :
    \begin{itemize}
        \item soit une \textbf{inéquation rationnelle paramétrique} (ex. $\frac{(m-1)(x-1)}{x+m-1} \le x$, $\frac{x^2+2x-2m}{x^2-m^2} > 0$, $\frac{|x|-m}{x^2-mx+1} \ge 0$, $x-1 \ge \frac{1}{mx}$, $\frac{x-m}{x^3-x} > 0$) ;
        \item soit une \textbf{inéquation avec valeur absolue ou logarithme paramétrique} (ex. $|x^2-a| < 1$, $\ln(-x^2+2x+a) < 0$, $\log_2\left(\frac{x+a}{x-1}\right) < 1$, $\log_{a-1}(x^2-2ax+1) > 0$).
    \end{itemize}
    La clé de la réussite réside dans l'étude des conditions d'existence, la factorisation, la comparaison de l'ordre des racines en fonction du paramètre et la construction de \textbf{tableaux de signes} clairs pour chaque plage de valeurs du paramètre.
    \item \textbf{Question 3 : Problème textuel de modélisation et résolution d'un système d'équations (Chapitre 8)} :
    Un récit concret met en scène une situation à décoder méthodiquement : panneaux photovoltaïques en achat ou leasing, addition au restaurant à la Mer du Nord, coalitions électorales à l'EPL, logistique de brasserie (bouteilles, casiers, fûts, palettes), autographes du groupe \emph{EPLighters} (diagramme de Venn), rénovation de living avec escalier suspendu, jus d'oranges aux Portes Ouvertes, filets de ballons de sport, câblage électrique de trois maisons, ruches d'apiculture face au varroa et aux frelons, ou rotations de camions de pompiers et d'hélicoptère au bois de Lauzelle.
\end{enumerate}

\section*{Conventions de notation et règles d'or pour le jour de l'examen}

\begin{itemize}
    \item \textbf{Notations officielles de l'examen :}
    \begin{itemize}
        \item Le nombre complexe conjugué de $z = a+bi$ est noté $z^*$ (ou $\bar{z}$) : $z^* = a-bi$.
        \item Le nombre de combinaisons de $k$ objets parmi $n$ est noté indifféremment $C_n^k$ ou $\binom{n}{k} = \frac{n!}{k!(n-k)!}$.
        \item Le logarithme népérien (en base $e$) est noté $\ln(x)$ (attention : dans certains énoncés, la base de $\log$ est explicitement précisée dans le texte, par exemple « où $\log$ est le logarithme en base $10$ » ou « en base $e$ » : \textbf{lisez toujours la précision qui suit la formule !}).
        \item Les intervalles ouverts s'écrivent avec des crochets retournés $]a, b[$ et l'exclusion d'un ensemble fini se note $\R \setminus \{-1, 0, 1\}$.
    \end{itemize}
    \item \textbf{Stratégie pour la Question 1 (7 cadres sans justification) :} Même si votre brouillon n'est pas lu pour la Question 1, \textbf{vérifiez systématiquement vos conditions d'existence} ! Par exemple, pour une équation logarithmique ou irrationnelle qui fournit deux candidats $x_1$ et $x_2$ après développement, testez-les tous les deux dans l'équation de départ : très souvent, l'un des deux est une solution parasite à rejeter.
    \item \textbf{Stratégie pour la Question 2 (discussion paramétrique) :} Ne multipliez \textbf{jamais} les deux membres d'une inéquation par un dénominateur contenant $x$ ou le paramètre $m$ sans connaître son signe ! Ramenez toujours tout dans le membre de gauche ($\ge 0$ ou $< 0$), mettez au même dénominateur, factorisez et dressez un tableau de signes par cas. N'oubliez jamais de tester les \textbf{cas limites} où deux racines se croisent ou coïncident avec une valeur interdite.
    \item \textbf{Stratégie pour la Question 3 (problème textuel) :} Commencez toujours votre copie en \textbf{définissant explicitement vos inconnues avec leurs unités} (« Soit $x$ le prix d'un mètre de câble en €/m... »). Numérotez vos équations $(1), (2), (3)\dots$ en indiquant quelle phrase de l'énoncé chacune traduit, et effectuez le contrôle final en réinjectant vos valeurs dans l'énoncé.
\end{itemize}

\tableofcontents

% ----------------------------------------------------------------------------
% CHAPITRES
% ----------------------------------------------------------------------------
% ============================================================================
% CHAPITRE 1 : POLYNÔMES, FACTORISATION ET DIVISION EUCLIDIENNE
% ============================================================================
\chapter{Polynômes, Factorisation et Division Euclidienne}

Les polynômes constituent la colonne vertébrale du calcul algébrique. À l'examen d'admission de l'UCLouvain, au moins une (et souvent deux) des sept sous-questions de la Question~1 porte directement sur la division euclidienne de polynômes, le théorème du reste, la factorisation ou la détermination de paramètres $(a,b)$ assurant la divisibilité d'un polynôme par un autre. De plus, les techniques de factorisation étudiées ici seront indispensables au Chapitre~7 pour dresser les tableaux de signes des inéquations paramétriques.

% ----------------------------------------------------------------------------
\section{Généralités sur les polynômes et identités remarquables}
% ----------------------------------------------------------------------------

\subsection{Définition, degré et égalité de deux polynômes}

\begin{definition}[Polynôme à coefficients réels]
Un \textbf{polynôme} $P(x)$ à coefficients réels en la variable $x$ est une expression algébrique de la forme :
\[
    P(x) = a_n x^n + a_{n-1} x^{n-1} + \dots + a_1 x + a_0 = \sum_{k=0}^n a_k x^k,
\]
où $n \in \N$ est un entier naturel et $a_0, a_1, \dots, a_n \in \R$ sont les \textbf{coefficients} du polynôme.
\begin{itemize}
    \item Si $a_n \neq 0$, l'entier $n$ est appelé le \textbf{degré} du polynôme, noté $\deg(P) = n$.
    \item Le terme $a_n x^n$ est le \textbf{terme dominant}, $a_n$ est le \textbf{coefficient dominant} (si $a_n = 1$, le polynôme est dit \emph{unitaire} ou \emph{normalisé}), et $a_0 = P(0)$ est le \textbf{terme indépendant} (ou terme constant).
    \item Par convention, le polynôme nul $P(x) = 0$ a un degré égal à $-\infty$.
\end{itemize}
\end{definition}

\begin{theoreme}[Principe d'identification des coefficients et règles sur les degrés]
\begin{enumerate}
    \item \textbf{Égalité de deux polynômes :} Deux polynômes non nuls $P(x) = \sum_{k=0}^n a_k x^k$ et $Q(x) = \sum_{k=0}^m b_k x^k$ prennent la même valeur pour tout $x \in \R$ ($\forall x \in \R : P(x) = Q(x)$) \textbf{si et seulement si} ils ont le même degré ($n = m$) et leurs coefficients de même degré sont deux à deux égaux :
    \[
        \boxed{\forall k \in \{0, 1, \dots, n\} : a_k = b_k}.
    \]
    En particulier, un polynôme est identiquement nul sur $\R$ si et seulement si tous ses coefficients sont nuls.
    \item \textbf{Degré d'une somme et d'un produit :} Pour tous polynômes $P$ et $Q$ :
    \[
        \deg(P + Q) \le \max\bigl(\deg(P), \deg(Q)\bigr) \qquad \text{et} \qquad \boxed{\deg(P \cdot Q) = \deg(P) + \deg(Q)}.
    \]
\end{enumerate}
\end{theoreme}

\subsection{Les identités remarquables indispensables sans calculatrice}

Le jour de l'examen, savoir reconnaître instantanément une identité remarquable permet d'éviter de longs calculs et de factoriser en quelques secondes des polynômes de degré $3$, $4$ ou $6$.

\begin{theoreme}[Formulaire des identités remarquables (degrés 2, 3 et $n$)]
Pour tous réels (ou complexes) $a$ et $b$ :
\begin{itemize}
    \item \textbf{Carré d'une somme ou d'une différence :}
    \[
        (a + b)^2 = a^2 + 2ab + b^2 \qquad \text{et} \qquad (a - b)^2 = a^2 - 2ab + b^2.
    \]
    \item \textbf{Différence de deux carrés :}
    \[
        \boxed{a^2 - b^2 = (a - b)(a + b)}.
    \]
    \item \textbf{Cube d'une somme ou d'une différence :}
    \[
        (a + b)^3 = a^3 + 3a^2b + 3ab^2 + b^3 \qquad \text{et} \qquad (a - b)^3 = a^3 - 3a^2b + 3ab^2 - b^3.
    \]
    \item \textbf{Différence et somme de deux cubes :}
    \[
        \boxed{a^3 - b^3 = (a - b)(a^2 + ab + b^2)} \qquad \text{et} \qquad \boxed{a^3 + b^3 = (a + b)(a^2 - ab + b^2)}.
    \]
    \item \textbf{Généralisation à un degré $n \in \N_0$ quelconque :}
    \[
        a^n - b^n = (a - b)\bigl(a^{n-1} + a^{n-2}b + a^{n-3}b^2 + \dots + ab^{n-2} + b^{n-1}\bigr).
    \]
\end{itemize}
\end{theoreme}

\begin{astuce}[L'identité de Sophie Germain : factoriser $a^4 + 4b^4$]
Une somme de deux carrés $a^2 + b^2$ ne se factorise pas dans $\R$. En revanche, une expression de la forme $x^4 + 4$ (posée à l'examen de \textbf{Juillet 2023 -- Série 1, Q1.6}) se factorise élégamment en complétant le carré par rapport au double produit $4x^2$ :
\[
    x^4 + 4 = (x^4 + 4x^2 + 4) - 4x^2 = (x^2 + 2)^2 - (2x)^2 = \boxed{(x^2 - 2x + 2)(x^2 + 2x + 2)}.
\]
Plus généralement, $a^4 + 4b^4 = (a^2 + 2b^2)^2 - (2ab)^2 = (a^2 - 2ab + 2b^2)(a^2 + 2ab + 2b^2)$.
\end{astuce}

% ----------------------------------------------------------------------------
\section{Trinôme du second degré et relations Somme-Produit de Viète}
% ----------------------------------------------------------------------------

\subsection{Racines, factorisation et signe du trinôme $ax^2 + bx + c$}

\begin{theoreme}[Étude complète du trinôme du second degré ($a \neq 0$)]
Soit $P(x) = ax^2 + bx + c$ avec $a \in \R_0$ et $b, c \in \R$. En mettant sous \textbf{forme canonique} :
\[
    ax^2 + bx + c = a\left[\left(x + \frac{b}{2a}\right)^2 - \frac{\Delta}{4a^2}\right] \qquad \text{où le discriminant vaut } \boxed{\Delta = b^2 - 4ac}.
\]
Trois cas se présentent selon le signe de $\Delta$ :
\begin{enumerate}
    \item \textbf{Si $\Delta > 0$ :} Le trinôme possède \textbf{deux racines réelles distinctes} :
    \[
        x_1 = \frac{-b - \sqrt{\Delta}}{2a} \qquad \text{et} \qquad x_2 = \frac{-b + \sqrt{\Delta}}{2a}.
    \]
    Sa factorisation dans $\R$ est $\boxed{P(x) = a(x - x_1)(x - x_2)}$.
    \item \textbf{Si $\Delta = 0$ :} Le trinôme possède \textbf{une racine réelle double} $x_0 = -\dfrac{b}{2a}$, et se factorise en :
    \[
        \boxed{P(x) = a(x - x_0)^2}.
    \]
    \item \textbf{Si $\Delta < 0$ :} Le trinôme n'admet \textbf{aucune racine réelle} (il est irréductible dans $\R$, mais possède deux racines complexes conjuguées dans $\C$, voir Chapitre~4).
\end{enumerate}
\end{theoreme}

\begin{theoreme}[Règle du signe d'un trinôme du second degré]
\begin{itemize}
    \item \textbf{Si $\Delta > 0$ (racines $x_1 < x_2$) :} Le trinôme $ax^2+bx+c$ est du \textbf{signe de $a$ à l'extérieur des racines} ($x \in \,]-\infty, x_1[ \cup \,]x_2, +\infty[$) et du \textbf{signe opposé à $a$ entre les racines} ($x \in \,]x_1, x_2[$).
    \item \textbf{Si $\Delta = 0$ (racine double $x_0$) :} Le trinôme est du \textbf{signe de $a$ pour tout $x \neq x_0$}, et s'annule en $x_0$.
    \item \textbf{Si $\Delta < 0$ :} Le trinôme ne s'annule jamais et garde strictement le \textbf{signe de $a$ sur tout $\R$}.
\end{itemize}
\end{theoreme}

\begin{center}
\begin{tikzpicture}[scale=0.92, >=Stealth]
    % Cas Delta > 0
    \begin{scope}[shift={(0,0)}]
        \draw[->, thick, darkgray] (-2.0,0) -- (2.2,0) node[right] {\small $x$};
        \draw[very thick, uclblue, domain=-1.75:1.95, samples=60] plot (\x, {0.7*(\x-0.1)^2 - 0.9});
        \fill[warnred] (-1.03,0) circle (2pt) node[below left] {\small $x_1$};
        \fill[warnred] (1.23,0) circle (2pt) node[below right] {\small $x_2$};
        \node[uclblue, font=\bfseries] at (-1.55,0.45) {\small sgn($a$)};
        \node[warnred, font=\bfseries] at (0.1,-0.45) {\small $-$sgn($a$)};
        \node[uclblue, font=\bfseries] at (1.75,0.45) {\small sgn($a$)};
        \node[darkgray, font=\small\bfseries] at (0.1,-1.55) {Cas $\Delta > 0$ ($a>0$)};
    \end{scope}

    % Cas Delta = 0
    \begin{scope}[shift={(5.3,0)}]
        \draw[->, thick, darkgray] (-2.0,0) -- (2.2,0) node[right] {\small $x$};
        \draw[very thick, defgreen, domain=-1.55:1.75, samples=60] plot (\x, {0.7*(\x-0.1)^2});
        \fill[warnred] (0.1,0) circle (2pt) node[below] {\small $x_0$};
        \node[defgreen, font=\bfseries] at (-1.2,0.55) {\small sgn($a$)};
        \node[defgreen, font=\bfseries] at (1.4,0.55) {\small sgn($a$)};
        \node[darkgray, font=\small\bfseries] at (0.1,-1.55) {Cas $\Delta = 0$ ($a>0$)};
    \end{scope}

    % Cas Delta < 0
    \begin{scope}[shift={(10.6,0)}]
        \draw[->, thick, darkgray] (-2.0,0) -- (2.2,0) node[right] {\small $x$};
        \draw[very thick, exorange, domain=-1.5:1.7, samples=60] plot (\x, {0.65*(\x-0.1)^2 + 0.4});
        \node[exorange, font=\bfseries] at (0.1,0.85) {\small Toujours sgn($a$)};
        \node[darkgray, font=\small\bfseries] at (0.1,-1.55) {Cas $\Delta < 0$ ($a>0$)};
    \end{scope}
\end{tikzpicture}
\end{center}

\subsection{Relations de Viète entre coefficients et racines (Somme et Produit)}

\begin{theoreme}[Relations Somme-Produit de Viète]
Si le trinôme $P(x) = ax^2 + bx + c$ ($a \neq 0$) possède deux racines $x_1$ et $x_2$ (distinctes ou confondues), alors en développant $a(x - x_1)(x - x_2) = ax^2 - a(x_1 + x_2)x + ax_1 x_2$, l'identification des coefficients donne :
\[
    \boxed{S = x_1 + x_2 = -\frac{b}{a}} \qquad \text{et} \qquad \boxed{P = x_1 x_2 = \frac{c}{a}}.
\]
Réciproquement, deux nombres réels dont la somme vaut $S$ et le produit vaut $P$ sont les racines de tout trinôme de la forme :
\[
    \boxed{a\bigl(x^2 - Sx + P\bigr) = 0 \qquad (a \in \R_0)}.
\]
\end{theoreme}

\begin{astuce}[Trois usages majeurs des formules Somme-Produit à l'examen]
\begin{enumerate}
    \item \textbf{Construire un polynôme connaissant la somme et le produit de ses racines} (question directe).
    \item \textbf{Trouver la seconde racine de tête} lorsqu'on connaît une racine évidente $x_1$ (comme $1$, $-1$, $2$ ou $1/3$) : il suffit d'utiliser $x_2 = \dfrac{c}{a x_1}$ ou $x_2 = -\dfrac{b}{a} - x_1$ sans calculer $\Delta$ !
    \item \textbf{Déterminer le signe des racines d'un trinôme paramétrique} (Question 2 de l'examen) :
    \begin{itemize}
        \item Si $P = \frac{c}{a} < 0$, les deux racines sont de \textbf{signes contraires} ($x_1 < 0 < x_2$).
        \item Si $P = \frac{c}{a} > 0$ (et $\Delta \ge 0$), les deux racines ont le \textbf{même signe}, qui est celui de leur somme $S = -\frac{b}{a}$.
    \end{itemize}
\end{enumerate}
\end{astuce}

\begin{exemple}[Application directe -- UCLouvain, Juillet 2021, Q1.5]
\textbf{Énoncé :} Donner un polynôme du second degré dont les racines ont $6$ pour somme et $8$ pour produit.

\textbf{Résolution :} Avec $S = x_1 + x_2 = 6$ et $P = x_1 x_2 = 8$, la formule $a(x^2 - Sx + P)$ donne immédiatement la famille de tous les polynômes solutions (pour tout réel non nul $a \in \R_0$, ou en particulier pour $a = 1$) :
\[
    \boxed{P(x) = x^2 - 6x + 8} \quad \text{(ou plus généralement } a(x^2 - 6x + 8) = ax^2 - 6ax + 8a \text{ avec } a \neq 0\text{)}.
\]
\emph{Vérification :} $x^2 - 6x + 8 = (x - 2)(x - 4)$, dont les racines $2$ et $4$ vérifient bien $2 + 4 = 6$ et $2 \times 4 = 8$.
\end{exemple}

% ----------------------------------------------------------------------------
\section{Division euclidienne des polynômes et théorème du reste}
% ----------------------------------------------------------------------------

\subsection{Le théorème de la division euclidienne}

\begin{theoreme}[Division euclidienne de deux polynômes]
Soient $P(x)$ (le \textbf{dividende}) et $D(x)$ (le \textbf{diviseur}) deux polynômes tels que $D(x)$ ne soit pas le polynôme nul. Il existe un \textbf{unique} couple de polynômes $\bigl(Q(x), R(x)\bigr)$, appelés respectivement le \textbf{quotient} et le \textbf{reste}, tels que :
\[
    \boxed{P(x) = D(x)\cdot Q(x) + R(x) \qquad \text{avec } \deg(R) < \deg(D)}.
\]
On dit que $P(x)$ est \textbf{divisible} par $D(x)$ (ou que $D(x)$ divise $P(x)$) si et seulement si le reste est le polynôme nul : $\boxed{R(x) = 0}$.
\end{theoreme}

\begin{attention}[Condition cruciale sur le degré du reste]
Dans une division euclidienne par un diviseur $D(x)$ :
\begin{itemize}
    \item Si $D(x) = x - a$ est de \textbf{degré $1$}, le reste $R$ est de degré $< 1$, donc c'est une \textbf{constante réelle} $R \in \R$.
    \item Si $D(x) = ax^2 + bx + c$ est de \textbf{degré $2$}, le reste est au plus de \textbf{degré $1$}, soit de la forme $R(x) = ux + v$.
    \item Si $D(x)$ est de \textbf{degré $3$} (ex. $x^3 - 1$), le reste est au plus de \textbf{degré $2$}, soit $R(x) = ux^2 + vx + w$.
\end{itemize}
La division posée (en potence) ne s'arrête que lorsque le degré du reste partiel devient \textbf{strictement inférieur} au degré du diviseur !
\end{attention}

\subsection{Poser une division euclidienne de polynômes (en potence)}

Pour diviser $P(x)$ par $D(x)$ :
\begin{enumerate}
    \item Ordonnez $P(x)$ et $D(x)$ selon les \textbf{puissances décroissantes} de $x$, en laissant un espace (ou en écrivant $+0x^k$) pour les puissances manquantes.
    \item À chaque étape, divisez le \textbf{terme dominant} du dividende partiel par le \textbf{terme dominant} du diviseur $D(x)$ : cela donne le terme suivant du quotient $Q(x)$.
    \item Multipliez tout le diviseur $D(x)$ par ce terme et soustrayez le résultat au dividende partiel.
\end{enumerate}

\begin{exemple}[Division euclidienne -- UCLouvain, Septembre 2022, Q1.4]
\textbf{Énoncé :} Calculer le quotient et le reste de la division de $x^6 + x^4 + 1$ par $x^2 - 2$.

\textbf{Résolution (Méthode 1 : division euclidienne posée) :}
\[
\begin{array}{r|l}
    x^6 + x^4 + 0x^2 + 1 & x^2 - 2 \\
    \cline{2-2}
    -\,(x^6 - 2x^4) \phantom{+ 0x^2 + 1} & x^4 + 3x^2 + 6 = Q(x) \\
    \cline{1-1}
    3x^4 + 0x^2 + 1 & \\
    -\,(3x^4 - 6x^2) \phantom{+ 1} & \\
    \cline{1-1}
    6x^2 + 1 & \\
    -\,(6x^2 - 12) & \\
    \cline{1-1}
    13 & = R(x)
\end{array}
\]
\textbf{Méthode 2 (Astuce par changement de variable $t = x^2$) :} Comme seuls des exposants pairs apparaissent, posons $t = x^2$. On divise $A(t) = t^3 + t^2 + 1$ par $t - 2$. Le reste vaut $A(2) = 2^3 + 2^2 + 1 = 8 + 4 + 1 = 13$, et par Horner en $t=2$, le quotient en $t$ est $t^2 + 3t + 6 = x^4 + 3x^2 + 6$.

On obtient donc : $\boxed{\text{quotient } Q(x) = x^4 + 3x^2 + 6 \quad \text{et} \quad \text{reste } R(x) = 13}$.
\end{exemple}

\begin{exemple}[Astuce de substitution modulaire pour le reste -- UCLouvain, Septembre 2023, Q1.6]
\textbf{Énoncé :} Quel est le reste de la division de l'expression $P(x) = (x^4 + 1)(x^2 - 1)$ par $D(x) = x^3 - 1$ ?

\textbf{Résolution :} Développons d'abord le dividende :
\[
    P(x) = x^6 - x^4 + x^2 - 1.
\]
\textbf{Astuce ultra-rapide (« calcul modulo $x^3 - 1$ ») :} Dans la division par $x^3 - 1$, écrire $x^3 = (x^3 - 1) + 1$ signifie que chaque bloc $x^3$ laisse un reste égal à $1$ lorsqu'on le divise par $x^3 - 1$ ! En remplaçant $x^3$ par $1$ (donc $x^6 = (x^3)^2$ par $1^2 = 1$, et $x^4 = x^3 \cdot x$ par $1 \cdot x = x$) :
\[
    P(x) = (x^3)^2 - (x^3)\cdot x + x^2 - 1 \equiv 1^2 - 1\cdot x + x^2 - 1 = x^2 - x \pmod{x^3 - 1}.
\]
Puisque $\deg(x^2 - x) = 2 < \deg(x^3 - 1) = 3$, c'est directement le reste de la division euclidienne !
\emph{Vérification explicite :} $(x^3 - 1)(x^3 + 1 - x) + (x^2 - x) = x^6 - x^4 + x^2 - 1 = P(x)$.
Le reste cherché est donc $\boxed{R(x) = x^2 - x}$.
\end{exemple}

\subsection{Théorème du reste et divisibilité par $(x - a)$}

\begin{theoreme}[Loi du reste et condition de divisibilité par $(x - a)$]
Soit $P(x)$ un polynôme et $a \in \R$ un nombre réel. La division euclidienne de $P(x)$ par $x - a$ s'écrit :
\[
    P(x) = (x - a)\cdot Q(x) + R \qquad \text{où } R \in \R \text{ est une constante}.
\]
En évaluant cette égalité en $x = a$, le terme $(a - a)Q(a)$ s'annule et l'on obtient :
\begin{enumerate}
    \item \textbf{Théorème du reste :} Le reste de la division de $P(x)$ par $(x - a)$ est égal à la valeur numérique $\boxed{R = P(a)}$.\\
    \emph{Attention au signe :} le reste de la division de $P(x)$ par $(x + a) = \bigl(x - (-a)\bigr)$ vaut $\boxed{R = P(-a)}$.
    \item \textbf{Corollaire (Divisibilité par $x - a$) :} Un polynôme $P(x)$ est \textbf{divisible par $(x - a)$} si et seulement si $a$ est une \textbf{racine} de $P(x)$ :
    \[
        \boxed{P(x) \text{ est divisible par } (x - a) \iff P(a) = 0}.
    \]
\end{enumerate}
\end{theoreme}

\begin{exemple}[Application du théorème du reste -- UCLouvain, Juillet 2026, Q1.6]
\textbf{Énoncé :} Déterminer la valeur du coefficient réel $a$ pour que le polynôme $P(x) = x^3 - 4x^2 + ax + 5$ soit divisible par $x - 2$.

\textbf{Résolution :} Par le théorème du reste, $P(x)$ est divisible par $x - 2$ si et seulement si $P(2) = 0$. Calculons $P(2)$ :
\[
    P(2) = 2^3 - 4(2^2) + 2a + 5 = 8 - 16 + 2a + 5 = 2a - 3.
\]
L'équation $P(2) = 0$ donne immédiatement $2a - 3 = 0 \iff \boxed{a = \dfrac{3}{2}}$.
\end{exemple}

% ----------------------------------------------------------------------------
\section{Schéma de Horner et factorisation des polynômes de degré $\ge 3$}
% ----------------------------------------------------------------------------

\subsection{La grille de Horner pour diviser par $(x - a)$}

Lorsqu'on divise un polynôme $P(x) = a_n x^n + a_{n-1} x^{n-1} + \dots + a_1 x + a_0$ par un binôme du premier degré $(x - \alpha)$, le \textbf{schéma de Horner} (ou tableau de Horner) permet d'obtenir en une seule ligne tous les coefficients du quotient $Q(x)$ ainsi que le reste $R = P(\alpha)$.

\begin{astuce}[Fonctionnement de la grille de Horner pour diviser par $(x - \alpha)$]
\begin{enumerate}
    \item Sur la première ligne, inscrivez tous les coefficients de $P(x)$ par ordre de puissances décroissantes (\textbf{sans oublier de mettre $0$ pour chaque puissance manquante !}).
    \item Placez la valeur $\alpha$ (qui annule $x - \alpha = 0$) dans la colonne de gauche.
    \item Abaissez directement le premier coefficient $a_n$ sur la troisième ligne.
    \item Ensuite, répétez la boucle : \textbf{multipliez} le nombre obtenu en bas par $\alpha$, inscrivez le produit sur la deuxième ligne dans la colonne suivante, puis \textbf{additionnez} verticalement avec le coefficient du dessus.
    \item La dernière case en bas à droite est le \textbf{reste $R = P(\alpha)$}, et les cases précédentes sont les coefficients du \textbf{quotient $Q(x)$} (de degré $n - 1$).
\end{enumerate}
\end{astuce}

\begin{exemple}[Division par $(x + 1)$ avec Horner -- UCLouvain, Juillet 2021, Q1.6]
\textbf{Énoncé :} Fournir le quotient $Q(x)$ et le reste $R(x)$ de la division du polynôme $(x - 1)^2$ par $(x + 1)$.

\textbf{Résolution :} Développons $P(x) = (x - 1)^2 = x^2 - 2x + 1$. On divise par $x + 1 = x - (-1)$, donc on applique Horner avec $\alpha = -1$ :
\[
\begin{array}{c|cc|c}
    & x^2 & x^1 & x^0 \\
    & 1 & -2 & 1 \\
    \alpha = -1 & \downarrow & -1 & 3 \\
    \hline
    & 1 & -3 & \mathbf{4 = R}
\end{array}
\]
On lit directement sur la dernière ligne : $\boxed{Q(x) = x - 3 \quad \text{et} \quad R(x) = 4}$.
\emph{Alternative mentale :} $(x - 1)^2 = \bigl((x + 1) - 2\bigr)^2 = (x + 1)^2 - 4(x + 1) + 4 = (x + 1)(x - 3) + 4$.
\end{exemple}

\begin{exemple}[Recherche des racines restantes d'un polynôme de degré 3 -- UCLouvain, Juillet 2022, Q1.5]
\textbf{Énoncé :} Sachant qu'une des racines du polynôme $P(x) = 9x^3 + 12x^2 - 11x + 2$ est $1/3$, que valent les deux autres racines ?

\textbf{Résolution :} Puisque $\alpha = \frac{1}{3}$ est une racine de $P(x)$, $P(x)$ est divisible par $\left(x - \frac{1}{3}\right)$. Dressons la grille de Horner avec $\alpha = \frac{1}{3}$ :
\[
\begin{array}{c|ccc|c}
    & x^3 & x^2 & x^1 & x^0 \\
    & 9 & 12 & -11 & 2 \\
    \alpha = \tfrac{1}{3} & \downarrow & 3 & 5 & -2 \\
    \hline
    & 9 & 15 & -6 & \mathbf{0 = R}
\end{array}
\]
Le reste est bien nul ($R = 0$), et le quotient vaut $Q(x) = 9x^2 + 15x - 6 = 3(3x^2 + 5x - 2)$.
Cherchons les racines de $3x^2 + 5x - 2 = 0$ :
\begin{itemize}
    \item \textbf{Par le discriminant :} $\Delta = 5^2 - 4(3)(-2) = 25 + 24 = 49 = 7^2$, d'où :
    \[
        x_2 = \frac{-5 - 7}{6} = -2 \qquad \text{et} \qquad x_3 = \frac{-5 + 7}{6} = \frac{2}{6} = \frac{1}{3}.
    \]
    \item \textbf{Ou par racine évidente :} $\frac{1}{3}$ annule encore $3x^2 + 5x - 2$ (c'est une racine double !), et par le produit des racines $x_2 \cdot \frac{1}{3} = -\frac{2}{3} \implies x_2 = -2$.
\end{itemize}
Les deux autres racines sont donc $\boxed{-2 \text{ et } \dfrac{1}{3}}$ (et $P(x) = 9\left(x - \frac{1}{3}\right)^2(x + 2) = (3x - 1)^2(x + 2)$).
\end{exemple}

\subsection{Méthodes de factorisation des polynômes de degré 3, 4 ou 5}

Face à un polynôme de degré supérieur ou égal à $3$, procédez toujours dans l'ordre suivant :

\begin{astuce}[Les 4 réflexes de factorisation à l'examen]
\begin{enumerate}
    \item \textbf{Mise en évidence d'une puissance de $x$ :} Si le terme indépendant est nul ($a_0 = 0$), mettez la plus grande puissance possible de $x$ en facteur commun !
    \item \textbf{Factorisation par groupement ($2$ par $2$) :} Pour un polynôme à $4$ termes $ax^3 + bx^2 + cx + d$, si $\frac{a}{b} = \frac{c}{d}$, groupez les deux premiers termes et les deux derniers :
    \[
        x^3 - 3x^2 - 4x + 12 = x^2(x - 3) - 4(x - 3) = (x - 3)(x^2 - 4) = (x - 3)(x - 2)(x + 2).
    \]
    \item \textbf{Polynôme bicarré ($ax^4 + bx^2 + c$) :} Posez $t = x^2 \ge 0$, factorisez $at^2 + bt + c$ en $t$, puis revenez à $x$.
    \item \textbf{Recherche d'une racine entière évidente ($\alpha \in \{\pm 1, \pm 2, \pm 3\}$ diviseur de $a_0$) :} Dès que vous trouvez $\alpha$ tel que $P(\alpha) = 0$, divisez par $(x - \alpha)$ via la grille de Horner.
\end{enumerate}
\end{astuce}

\begin{exemple}[Factorisations complètes -- UCLouvain, Juillet 2023 Série 2 Q1.5 \& Septembre 2025 Q1.6]
\begin{enumerate}
    \item \textbf{Juillet 2023 -- Série 2, Q1.5 :} \emph{Factoriser l'expression $P(x) = x^5 - 8x^3 + 16x$.}\\
    On commence par mettre $x$ en évidence : $P(x) = x(x^4 - 8x^2 + 16)$. Dans la parenthèse, on reconnaît le carré parfait $(x^2 - 4)^2$ (car $(x^2)^2 - 2\cdot 4\cdot x^2 + 4^2 = (x^2-4)^2$). Enfin, on applique la différence de deux carrés $x^2 - 4 = (x - 2)(x + 2)$ :
    \[
        P(x) = x\bigl(x^2 - 4\bigr)^2 = x\bigl((x - 2)(x + 2)\bigr)^2 = \boxed{x(x - 2)^2(x + 2)^2}.
    \]
    \item \textbf{Septembre 2025, Q1.6 :} \emph{Factoriser le polynôme $P(x) = x^3 - 3x^2 - 4x + 12$.}\\
    Par groupement des termes :
    \[
        P(x) = x^2(x - 3) - 4(x - 3) = (x - 3)(x^2 - 4) = \boxed{(x - 3)(x - 2)(x + 2)}.
    \]
    (On pouvait aussi regrouper $x(x^2 - 4) - 3(x^2 - 4) = (x^2 - 4)(x - 3)$).
\end{enumerate}
\end{exemple}

% ----------------------------------------------------------------------------
\section{Détermination de paramètres de divisibilité : les deux méthodes reines}
% ----------------------------------------------------------------------------

Un grand classique de la Question~1 de l'examen d'Algèbre (tombé en \textbf{Juillet 2024 Q1.4}, \textbf{Septembre 2024 Q1.1}, \textbf{Juillet 2025 Q1.7} et \textbf{Septembre 2026 Q1.6}) consiste à déterminer deux paramètres réels $a$ et $b$ pour qu'un polynôme de degré~4 (typiquement $P(x) = x^4 + ax + b$) soit divisible par un trinôme du second degré $D(x) = x^2 + px + q$.

Vous disposez de \textbf{deux méthodes infaillibles} qui donnent le résultat en moins de deux minutes :

\begin{theoreme}[Méthode A (Coefficients indéterminés) vs Méthode B (Division euclidienne)]
Pour imposer que $P(x) = x^4 + ax + b$ soit divisible par un trinôme unitaire $D(x) = x^2 + px + q$ :
\begin{itemize}
    \item \textbf{Méthode A (Identification d'un quotient $Q(x) = x^2 + ux + v$) :} Puisque $P$ est unitaire de degré~4 et $D$ est unitaire de degré~2, si le reste est nul, le quotient est nécessairement de la forme $Q(x) = x^2 + ux + v$. On écrit :
    \[
        x^4 + 0x^3 + 0x^2 + ax + b = (x^2 + px + q)(x^2 + ux + v).
    \]
    En développant le produit à droite et en identifiant les coefficients de $x^3$, $x^2$, $x^1$ et $x^0$ dans cet ordre, on trouve successivement $u$, puis $v$, puis $a$, puis $b$ sans résoudre aucun système compliqué !
    \item \textbf{Méthode B (Annulation du reste $R(x) = (\dots)x + (\dots)$ de la division euclidienne) :} On effectue la division euclidienne de $P(x)$ par $D(x)$ jusqu'au bout. On obtient un reste du premier degré $R(x) = \alpha(a) x + \beta(b)$. Le polynôme $P$ est divisible par $D$ si et seulement si $\alpha(a) = 0$ \textbf{et} $\beta(b) = 0$.
\end{itemize}
\end{theoreme}

\begin{exemple}[Application comparée des deux méthodes -- UCLouvain, Juillet 2024 Q1.4, Juillet 2025 Q1.7 \& Septembre 2026 Q1.6]
Résolvons simultanément les trois versions posées en 2024, 2025 et 2026 en écrivant la formule générale pour $P(x) = x^4 + ax + b$ divisible par $D(x) = x^2 + x + q$ :
posons $x^4 + 0x^3 + 0x^2 + ax + b = (x^2 + x + q)(x^2 + ux + v)$.
En développant le membre de droite :
\[
    (x^2 + x + q)(x^2 + ux + v) = x^4 + (u + 1)x^3 + (v + u + q)x^2 + (v + qu)x + qv.
\]
En identifiant avec $x^4 + 0x^3 + 0x^2 + ax + b$ degré par degré :
\begin{enumerate}
    \item Coefficient de $x^3$ : $u + 1 = 0 \implies \mathbf{u = -1}$.
    \item Coefficient de $x^2$ : $v + u + q = 0 \implies v - 1 + q = 0 \implies \mathbf{v = 1 - q}$.
    \item Coefficient de $x^1$ : $a = v + qu = (1 - q) - q = \mathbf{1 - 2q}$.
    \item Coefficient de $x^0$ : $b = qv = \mathbf{q(1 - q)}$.
\end{enumerate}
Appliquons ce schéma aux trois examens :
\begin{itemize}
    \item \textbf{Juillet 2024, Q1.4 ($D(x) = x^2 + x - 1$, soit $q = -1$) :}
    $u = -1$, $v = 1 - (-1) = 2$, d'où $a = 2 + (-1)(-1) = 3$ et $b = (-1)(2) = -2$.
    \[
        \boxed{a = 3 \quad \text{et} \quad b = -2} \qquad \bigl(\text{et } x^4 + 3x - 2 = (x^2 + x - 1)(x^2 - x + 2)\bigr).
    \]
    \item \textbf{Juillet 2025, Q1.7 ($D(x) = x^2 + x + 1$, soit $q = 1$) :}
    $u = -1$, $v = 1 - 1 = 0$, d'où $a = 0 + 1(-1) = -1$ et $b = 1 \cdot 0 = 0$.
    \[
        \boxed{a = -1 \quad \text{et} \quad b = 0} \qquad \bigl(\text{en effet, } x^4 - x = x(x^3 - 1) = x(x - 1)(x^2 + x + 1)\bigr).
    \]
    \item \textbf{Septembre 2026, Q1.6 ($D(x) = x^2 + x + 2$, soit $q = 2$) :}
    $u = -1$, $v = 1 - 2 = -1$, d'où $a = -1 + 2(-1) = -3$ et $b = 2(-1) = -2$.
    \[
        \boxed{a = -3 \quad \text{et} \quad b = -2} \qquad \bigl(\text{et } x^4 - 3x - 2 = (x^2 + x + 2)(x^2 - x - 1)\bigr).
    \]
\end{itemize}
\end{exemple}

\begin{exemple}[Identification du diviseur -- UCLouvain, Septembre 2024, Q1.1]
\textbf{Énoncé :} Déterminez les coefficients réels $a$ et $b$ pour que la division du polynôme $P(x) = 2x^4 - x^2 + x - 3$ par $D(x) = ax^2 + b$ donne comme quotient $Q(x) = x^2 - 1$ et comme reste $R(x) = x - 2$.

\textbf{Résolution :} Par définition de la division euclidienne $P(x) = D(x)Q(x) + R(x)$, on a :
\[
    D(x)\cdot Q(x) = P(x) - R(x) = (2x^4 - x^2 + x - 3) - (x - 2) = 2x^4 - x^2 - 1.
\]
Développons $D(x)Q(x) = (ax^2 + b)(x^2 - 1) = ax^4 + (b - a)x^2 - b$.
Par identification avec $2x^4 - x^2 - 1$ :
\begin{itemize}
    \item Terme en $x^4$ : $a = 2$.
    \item Terme constant : $-b = -1 \implies b = 1$ (et on vérifie que le terme en $x^2$ vaut bien $b - a = 1 - 2 = -1$).
\end{itemize}
On conclut que $\boxed{a = 2 \quad \text{et} \quad b = 1}$.
\end{exemple}

% ----------------------------------------------------------------------------
\section{Exercices d'entraînement du Chapitre 1}
% ----------------------------------------------------------------------------

\emph{Les corrigés détaillés de tous ces exercices se trouvent au Chapitre~10.}

\begin{enumerate}[label=\textbf{Exercice 1.\arabic* :}, leftmargin=*]
    \item \textbf{(UCLouvain -- Septembre 2021, Q1.6)} Fournir le quotient $Q(x)$ et le reste $R(x)$ de la division du polynôme $(x^4 - 1)$ par $(x^2 + 1)$.
    \item \textbf{(UCLouvain -- Juillet 2023, Série 1, Q1.6)} Quel est le quotient de la division du polynôme $x^4 + 4$ par $x^2 + 2x + 2$ ?
    \item \textbf{(Entraînement type Question 1)} Déterminer le quotient $Q(x)$ et le reste $R(x)$ de la division euclidienne de $P(x) = 2x^4 - 5x^3 + 4x^2 - 3x + 7$ par $x - 2$ à l'aide du schéma de Horner.
    \item \textbf{(Entraînement type Question 1)} Sachant que le polynôme $P(x) = 2x^3 - 3x^2 - 11x + 6$ admet $x_1 = 3$ pour racine, déterminer ses deux autres racines et sa forme complètement factorisée.
    \item \textbf{(Entraînement type Question 1)} Déterminer les réels $a$ et $b$ pour que le polynôme $P(x) = x^4 - 2x^3 + ax^2 + bx + 4$ soit divisible par $(x - 1)(x + 2)$.
    \item \textbf{(Entraînement type Question 1)} Soit $P(x)$ un polynôme dont le reste de la division par $(x - 1)$ vaut $5$ et dont le reste de la division par $(x + 2)$ vaut $-4$. Quel est le reste $R(x)$ de la division de $P(x)$ par $(x - 1)(x + 2)$ ?
    \item \textbf{(Entraînement type Question 1)} Factoriser complètement dans $\R$ le polynôme $P(x) = x^6 - 9x^4 - x^2 + 9$.
\end{enumerate}

% ============================================================================
% CHAPITRE 2 : ÉQUATIONS, VALEURS ABSOLUES, RADICAUX ET DOMAINES D'EXISTENCE
% ============================================================================
\chapter{Équations, Valeurs Absolues, Radicaux et Domaines d'Existence}

Dans ce deuxième chapitre, nous abordons trois thèmes incontournables de la Question~1 et de la Question~2 de l'examen d'admission de l'UCLouvain : la manipulation des \textbf{valeurs absolues} $|A(x)|$, la résolution rigoureuse des \textbf{équations irrationnelles} (comportant des racines carrées $\sqrt{A(x)}$) sans se faire piéger par les solutions parasites, et la détermination des \textbf{domaines de définition} (ou conditions d'existence) de fonctions composées mêlant fractions, radicaux et logarithmes.

% ----------------------------------------------------------------------------
\section{Valeur absolue : propriétés, équations et inéquations}
% ----------------------------------------------------------------------------

\subsection{Définition algébrique et interprétation géométrique}

\begin{definition}[Valeur absolue d'un nombre réel]
Pour tout réel $x \in \R$, la \textbf{valeur absolue} de $x$, notée $|x|$, représente la distance entre le point d'abscisse $x$ et l'origine $0$ sur la droite réelle :
\[
    \boxed{|x| = \begin{cases} x & \text{si } x \ge 0, \\ -x & \text{si } x < 0. \end{cases}}
\]
Plus généralement, $|x - a|$ est la \textbf{distance} entre les réels $x$ et $a$ sur la droite graduée.
\end{definition}

\begin{attention}[Le piège classique : $\sqrt{x^2} = |x|$ et non $x$ !]
Pour tout nombre réel $x \in \R$ (positif ou négatif), la racine carrée désigne toujours le réel \textbf{positif ou nul} dont le carré vaut $x^2$. Par conséquent :
\[
    \boxed{\forall x \in \R : \sqrt{x^2} = |x|} \qquad \text{(par exemple, } \sqrt{(-5)^2} = \sqrt{25} = 5 = |-5| \neq -5\text{)}.
\]
En revanche, $|x|^2 = x^2$ pour tout $x \in \R$.
\end{attention}

\begin{theoreme}[Propriétés fondamentales de la valeur absolue]
Pour tous réels $x, y \in \R$ :
\begin{enumerate}
    \item \textbf{Positivité et parité :} $|x| \ge 0$, $\;|x| = 0 \iff x = 0$, $\;$et $|-x| = |x|$.
    \item \textbf{Produit et quotient :} $|x \cdot y| = |x| \cdot |y|$ \quad et \quad $\left|\dfrac{x}{y}\right| = \dfrac{|x|}{|y|}$ (pour $y \neq 0$).
    \item \textbf{Inégalité triangulaire :} $|x + y| \le |x| + |y|$ \quad et \quad $\bigl||x| - |y|\bigr| \le |x - y|$.
\end{enumerate}
\end{theoreme}

\subsection{Résolution d'équations et d'inéquations avec valeurs absolues}

\begin{theoreme}[Équations et inéquations fondamentales avec $|A(x)|$]
Soit $A(x)$ une expression algébrique et $k \in \R$ un réel donné :
\begin{enumerate}
    \item \textbf{Équation $|A(x)| = k$ :}
    \begin{itemize}
        \item Si $k < 0$, l'équation n'a \textbf{aucune solution} ($S = \varnothing$), car une valeur absolue est toujours positive ou nulle.
        \item Si $k \ge 0$ : $\boxed{|A(x)| = k \iff A(x) = k \;\text{ \textbf{ou} }\; A(x) = -k}$.
    \end{itemize}
    \item \textbf{Équation $|A(x)| = |B(x)|$ :}
    \[
        \boxed{|A(x)| = |B(x)| \iff A(x) = B(x) \;\text{ \textbf{ou} }\; A(x) = -B(x)}.
    \]
    \item \textbf{Inéquation $|A(x)| < k$ (ou $\le k$) :}
    \begin{itemize}
        \item Si $k \le 0$ (pour $< k$) ou $k < 0$ (pour $\le k$), il n'y a \textbf{aucune solution} ($S = \varnothing$).
        \item Si $k > 0$ : $\boxed{|A(x)| < k \iff -k < A(x) < k \iff \begin{cases} A(x) > -k \\ \text{\textbf{et}} \\ A(x) < k \end{cases}}$ (intersection $\cap$).
    \end{itemize}
    \item \textbf{Inéquation $|A(x)| > k$ (ou $\ge k$) :}
    \begin{itemize}
        \item Si $k < 0$, l'inégalité est toujours vérifiée sur tout le domaine d'existence de $A(x)$.
        \item Si $k \ge 0$ : $\boxed{|A(x)| > k \iff A(x) < -k \;\text{ \textbf{ou} }\; A(x) > k}$ (union $\cup$).
    \end{itemize}
\end{enumerate}
\end{theoreme}

\begin{center}
\begin{tikzpicture}[scale=0.92, >=Stealth]
    % |X| < k
    \begin{scope}[shift={(0,0)}]
        \draw[->, thick, darkgray] (-3.5,0) -- (3.5,0) node[right] {\small $A(x)$};
        \draw[very thick, defgreen] (-1.8,0) -- (1.8,0);
        \fill[white] (-1.8,0) circle (2.2pt);
        \draw[thick, defgreen] (-1.8,0) circle (2.2pt) node[below=3pt] {\small $-k$};
        \fill[white] (1.8,0) circle (2.2pt);
        \draw[thick, defgreen] (1.8,0) circle (2.2pt) node[below=3pt] {\small $k$};
        \fill[darkgray] (0,0) circle (1.5pt) node[below=3pt] {\small $0$};
        \node[defgreen, font=\small\bfseries] at (0,0.45) {$|A(x)| < k \iff -k < A(x) < k$};
    \end{scope}

    % |X| > k
    \begin{scope}[shift={(8.2,0)}]
        \draw[->, thick, darkgray] (-3.5,0) -- (3.5,0) node[right] {\small $A(x)$};
        \draw[very thick, warnred] (-3.4,0) -- (-1.8,0);
        \draw[very thick, warnred, ->] (1.8,0) -- (3.5,0);
        \fill[white] (-1.8,0) circle (2.2pt);
        \draw[thick, warnred] (-1.8,0) circle (2.2pt) node[below=3pt] {\small $-k$};
        \fill[white] (1.8,0) circle (2.2pt);
        \draw[thick, warnred] (1.8,0) circle (2.2pt) node[below=3pt] {\small $k$};
        \fill[darkgray] (0,0) circle (1.5pt) node[below=3pt] {\small $0$};
        \node[warnred, font=\small\bfseries] at (0,0.45) {$|A(x)| > k \iff A(x) < -k \text{ ou } A(x) > k$};
    \end{scope}
\end{tikzpicture}
\end{center}

\begin{astuce}[Méthode générale de disjonction des cas (élimination des barres de valeur absolue)]
Lorsqu'une équation ou inéquation mêle $|A(x)|$ à d'autres termes dépendant de $x$, comme $|A(x)| \le B(x)$, commencez par \textbf{étudier le signe de l'expression $A(x)$ située à l'intérieur de la valeur absolue} pour découper $\R$ en deux régions :
\begin{itemize}
    \item \textbf{Cas 1 (Région où $A(x) \ge 0$) :} On remplace $|A(x)|$ par $A(x)$, on résout l'inéquation $A(x) \le B(x)$, puis on prend l'\textbf{intersection} avec la condition du Cas~1 ($A(x) \ge 0$).
    \item \textbf{Cas 2 (Région où $A(x) < 0$) :} On remplace $|A(x)|$ par $-A(x)$, on résout $-A(x) \le B(x)$, puis on prend l'\textbf{intersection} avec la condition du Cas~2 ($A(x) < 0$).
    \item \textbf{Synthèse finale :} On fait l'\textbf{union} ($S = S_1 \cup S_2$) des solutions obtenues dans les deux cas.
\end{itemize}
\end{astuce}

\begin{exemple}[Inéquation avec valeur absolue -- UCLouvain, Septembre 2021, Q2.1]
\textbf{Énoncé :} Déterminer l'ensemble des nombres réels $x$ qui vérifient l'inéquation suivante :
\[
    |x^2 - 2x - 6| \le 3x^2 - 2x - 6.
\]
\textbf{Résolution :} Étudions d'abord le signe de $A(x) = x^2 - 2x - 6$. Son discriminant vaut $\Delta = (-2)^2 - 4(1)(-6) = 4 + 24 = 28 = 4 \times 7$, d'où les deux racines :
\[
    r_1 = 1 - \sqrt{7} \approx -1,65 \qquad \text{et} \qquad r_2 = 1 + \sqrt{7} \approx 3,65.
\]
\begin{itemize}
    \item \textbf{1\ier{} cas : $x \in \,]-\infty, 1 - \sqrt{7}] \cup [1 + \sqrt{7}, +\infty[$ (région où $x^2 - 2x - 6 \ge 0$).}\\
    Dans cette région, $|x^2 - 2x - 6| = x^2 - 2x - 6$. L'inéquation devient :
    \[
        x^2 - 2x - 6 \le 3x^2 - 2x - 6 \iff 0 \le 2x^2 \iff x^2 \ge 0,
    \]
    ce qui est \textbf{toujours vrai} pour tout $x \in \R$. L'intersection avec le domaine du 1\ier{} cas donne donc :
    \[
        S_1 = \,]-\infty, 1 - \sqrt{7}] \cup [1 + \sqrt{7}, +\infty[.
    \]
    \item \textbf{2\ieme{} cas : $x \in \,]1 - \sqrt{7}, 1 + \sqrt{7}[$ (région où $x^2 - 2x - 6 < 0$).}\\
    Dans cette région, $|x^2 - 2x - 6| = -(x^2 - 2x - 6) = -x^2 + 2x + 6$. L'inéquation devient :
    \[
        -x^2 + 2x + 6 \le 3x^2 - 2x - 6 \iff 4x^2 - 4x - 12 \ge 0 \iff x^2 - x - 3 \ge 0.
    \]
    Le trinôme $x^2 - x - 3$ a pour discriminant $\Delta' = 1 + 12 = 13$, et ses racines sont $\alpha_1 = \frac{1 - \sqrt{13}}{2} \approx -1,30$ et $\alpha_2 = \frac{1 + \sqrt{13}}{2} \approx 2,30$. Il est positif ou nul à l'extérieur de ses racines : $x \le \frac{1-\sqrt{13}}{2}$ ou $x \ge \frac{1+\sqrt{13}}{2}$.
    Or, comme $4 < 7 < 13 < 28$, on a $1 - \sqrt{7} = \frac{2 - \sqrt{28}}{2} < \frac{1 - \sqrt{13}}{2} < \frac{1 + \sqrt{13}}{2} < \frac{2 + \sqrt{28}}{2} = 1 + \sqrt{7}$.
    L'intersection avec l'intervalle du 2\ieme{} cas donne donc :
    \[
        S_2 = \left]1 - \sqrt{7},\, \frac{1 - \sqrt{13}}{2}\right] \cup \left[\frac{1 + \sqrt{13}}{2},\, 1 + \sqrt{7}\right[.
    \]
\end{itemize}
\textbf{Synthèse ($S = S_1 \cup S_2$) :} Les bornes $1 - \sqrt{7}$ et $1 + \sqrt{7}$ se raccordent parfaitement ! On obtient :
\[
    \boxed{S = \left]-\infty,\, \frac{1 - \sqrt{13}}{2}\right] \cup \left[\frac{1 + \sqrt{13}}{2},\, +\infty\right[}.
\]
\end{exemple}

% ----------------------------------------------------------------------------
\section{Équations irrationnelles (avec racines carrées) et solutions parasites}
% ----------------------------------------------------------------------------

À l'examen d'admission (notamment en \textbf{Juillet 2023 Série 1 Q1.7}, \textbf{Juillet 2023 Série 2 Q1.7} et \textbf{Septembre 2023 Q1.7}), la Question~1 propose régulièrement une équation de la forme $\sqrt{A(x)} = B(x)$. C'est l'un des pièges les plus redoutables pour les candidats inattentifs !

\begin{theoreme}[Équivalence fondamentale pour résoudre $\sqrt{A(x)} = B(x)$]
Pour tout réel $x \in \R$, l'écriture $\sqrt{A(x)}$ désigne un nombre \textbf{positif ou nul}. Par conséquent, si $B(x) < 0$, l'égalité $\sqrt{A(x)} = B(x)$ est impossible ! L'équivalence exacte est :
\[
    \boxed{\sqrt{A(x)} = B(x) \iff \begin{cases} B(x) \ge 0 & \text{(condition de positivité du second membre)} \\ A(x) = \bigl(B(x)\bigr)^2 & \text{(élévation au carré)} \end{cases}}
\]
\emph{Remarque :} La condition d'existence du radical $A(x) \ge 0$ est automatiquement garantie dès lors que $A(x) = \bigl(B(x)\bigr)^2$, car un carré réel est toujours positif ou nul !
\end{theoreme}

\begin{attention}[Pourquoi l'élévation au carré seule crée-t-elle des solutions parasites ?]
Lorsque vous élevez les deux membres au carré sans condition :
\[
    A(x) = \bigl(B(x)\bigr)^2 \iff \sqrt{A(x)} = |B(x)| \iff \sqrt{A(x)} = B(x) \;\text{ \textbf{ou} }\; \textcolor{warnred}{\sqrt{A(x)} = -B(x)}.
\]
Les racines de $A(x) = (B(x))^2$ pour lesquelles $B(x) < 0$ sont en réalité les solutions de l'équation opposée $\sqrt{A(x)} = -B(x)$ : ce sont des \textbf{solutions parasites} qu'il faut impérativement éliminer !
\end{attention}

\begin{exemple}[Les trois équations irrationnelles de 2023 -- UCLouvain, Juillet \& Septembre 2023, Q1.7]
\begin{enumerate}
    \item \textbf{Juillet 2023 -- Série 1, Q1.7 :} \emph{Quelles sont les solutions dans $\R$ de l'équation $x = \sqrt{x + 1}$ ?}\\
    Appliquons le théorème avec $A(x) = x + 1$ et $B(x) = x$ :
    \[
        \sqrt{x + 1} = x \iff \begin{cases} x \ge 0 \\ x + 1 = x^2 \end{cases} \iff \begin{cases} x \ge 0 \\ x^2 - x - 1 = 0 \end{cases}
    \]
    Le trinôme $x^2 - x - 1 = 0$ ($\Delta = 1 + 4 = 5$) a deux racines : $x_1 = \frac{1 - \sqrt{5}}{2} < 0$ (car $\sqrt{5} > 1$) et $x_2 = \frac{1 + \sqrt{5}}{2} > 0$ (le célèbre \textbf{nombre d'or} $\varphi \approx 1,618$).
    La condition $x \ge 0$ élimine la solution parasite $x_1 = \frac{1-\sqrt{5}}{2}$, et il reste l'unique solution :
    \[
        \boxed{x = \frac{1 + \sqrt{5}}{2}}.
    \]
    \item \textbf{Juillet 2023 -- Série 2, Q1.7 :} \emph{Quelles sont les solutions dans $\R$ de l'équation $\sqrt{1 - x^2} = x - 1$ ?}\\
    Ici $A(x) = 1 - x^2$ et $B(x) = x - 1$. On pose :
    \[
        \sqrt{1 - x^2} = x - 1 \iff \begin{cases} x - 1 \ge 0 \iff x \ge 1 \\ 1 - x^2 = (x - 1)^2 \end{cases}
    \]
    Développons la seconde équation : $1 - x^2 = x^2 - 2x + 1 \iff 2x^2 - 2x = 0 \iff 2x(x - 1) = 0$, ce qui fournit deux candidats : $x = 0$ et $x = 1$.
    Or, pour $x = 0$, le membre de droite vaut $0 - 1 = -1 < 0$ (tandis que $\sqrt{1 - 0^2} = 1 \neq -1$) : $x = 0$ est une solution parasite ! Seul $x = 1 \ge 1$ convient :
    \[
        \boxed{x = 1}.
    \]
    \emph{Remarque rapide par les domaines :} La racine $\sqrt{1-x^2}$ exige $1-x^2 \ge 0 \iff x \in [-1, 1]$, tandis que la positivité du membre de droite exige $x - 1 \ge 0 \iff x \in [1, +\infty[$. L'intersection $[-1, 1] \cap [1, +\infty[ = \{1\}$ ne contient qu'un seul point possible : $x = 1$, qui vérifie bien $\sqrt{0} = 0$ !
    \item \textbf{Septembre 2023, Q1.7 :} \emph{Quelles sont les solutions dans $\R$ de l'équation $x - \sqrt{x + 1} = 1$ ?}\\
    Isolons d'abord le radical seul d'un côté avant d'élever au carré :
    \[
        \sqrt{x + 1} = x - 1 \iff \begin{cases} x \ge 1 \\ x + 1 = (x - 1)^2 \end{cases} \iff \begin{cases} x \ge 1 \\ x^2 - 3x = x(x - 3) = 0 \end{cases}
    \]
    Parmi les deux racines $x = 0$ et $x = 3$, la condition $x \ge 1$ rejette $x = 0$ (pour lequel $0 - \sqrt{1} = -1 \neq 1$) et conserve l'unique solution :
    \[
        \boxed{x = 3}.
    \]
\end{enumerate}
\end{exemple}

% ----------------------------------------------------------------------------
\section{Détermination du domaine de définition (Conditions d'existence)}
% ----------------------------------------------------------------------------

Presque chaque session d'examen d'Algèbre comporte à la Question~1 (souvent en Q1.2 ou Q1.5) une question demandant le \textbf{domaine de définition} $\dom f$ (ou les conditions d'existence) d'une fonction réelle $f(x)$.

\subsection{Les trois règles de base des conditions d'existence dans $\R$}

\begin{theoreme}[Catalogue des conditions d'existence (C.E.)]
Pour qu'une expression réelle $f(x)$ existe dans $\R$, toutes les opérations qui la composent doivent être simultanément licites (\textbf{intersection $\cap$} de toutes les conditions) :
\begin{enumerate}
    \item \textbf{Dénominateur d'une fraction $\dfrac{N(x)}{D(x)}$ :} Le dénominateur doit être \textbf{non nul} :
    \[
        \boxed{D(x) \neq 0}.
    \]
    \item \textbf{Racine carrée (ou d'indice pair $2n$) $\sqrt{A(x)}$ :} Le radicande sous la racine doit être \textbf{positif ou nul} :
    \[
        \boxed{A(x) \ge 0}.
    \]
    \emph{Cas fréquent :} Si la racine carrée est elle-même au dénominateur, comme $\dfrac{1}{\sqrt{A(x)}}$, la combinaison de $A(x) \ge 0$ et $\sqrt{A(x)} \neq 0$ donne une inégalité \textbf{stricte} : $\boxed{A(x) > 0}$.
    \item \textbf{Logarithme $\ln\bigl(U(x)\bigr)$ ou $\log_a\bigl(U(x)\bigr)$ :} L'argument du logarithme doit être \textbf{strictement positif} :
    \[
        \boxed{U(x) > 0} \qquad \text{(et si la base $a$ dépend d'un paramètre : } a > 0 \text{ et } a \neq 1\text{)}.
    \]
\end{enumerate}
\end{theoreme}

\subsection{Fonctions emboîtées : le motif récurrent $\sqrt{\ln(U(x))}$ et $\frac{1}{\sqrt{\ln(U(x))}}$}

Regardons attentivement les fonctions posées à l'examen en 2024, 2025 et 2026 : elles reposent presque toutes sur l'emboîtement d'une racine carrée et d'un logarithme népérien !

\begin{theoreme}[Conditions d'existence de $\sqrt{\ln(U(x))}$ et de $\dfrac{1}{\sqrt{\ln(U(x))}}$]
Puisque la fonction logarithme népérien $\ln : \,]0, +\infty[ \to \R$ est strictement croissante et s'annule en $1$ ($\ln(1) = 0$) :
\begin{itemize}
    \item Pour $f(x) = \sqrt{\ln\bigl(U(x)\bigr)}$, la racine impose $\ln\bigl(U(x)\bigr) \ge 0 = \ln(1)$, ce qui équivaut à :
    \[
        \boxed{\ln\bigl(U(x)\bigr) \ge 0 \iff U(x) \ge 1}.
    \]
    \item Pour $f(x) = \dfrac{1}{\sqrt{\ln\bigl(U(x)\bigr)}}$, la racine au dénominateur impose $\ln\bigl(U(x)\bigr) > 0 = \ln(1)$, ce qui équivaut à :
    \[
        \boxed{\ln\bigl(U(x)\bigr) > 0 \iff U(x) > 1}.
    \]
\end{itemize}
\emph{Observation très utile :} Dès que l'on impose $U(x) \ge 1$ (ou $U(x) > 1$), la condition d'existence du logarithme $U(x) > 0$ est \textbf{automatiquement satisfaite} puisque $1 > 0$ !
\end{theoreme}

\begin{exemple}[Les six questions de domaine d'existence (2021--2026)]
Passons en revue toutes les questions de domaine tombées à la Question~1 de 2021 à 2026 :
\begin{enumerate}
    \item \textbf{Septembre 2021, Q1.1 :} \emph{Domaine d'existence de $f(x) = \ln\left(\dfrac{x + 2}{x}\right)$.}\\
    Il faut $x \neq 0$ et l'argument du logarithme strictement positif : $\dfrac{x + 2}{x} > 0$. Le quotient de $(x+2)$ par $x$ a pour valeurs critiques $-2$ et $0$, et est strictement positif à l'extérieur de $[-2, 0]$ :
    \[
        \boxed{\dom f = \,]-\infty, -2[ \,\cup\, ]0, +\infty[}.
    \]

    \item \textbf{Juillet 2024, Q1.5 :} \emph{Domaine de définition de $f(x) = \dfrac{1}{\sqrt{\ln(x^2 - 1)}}$.}\\
    La racine au dénominateur impose $\ln(x^2 - 1) > 0 \iff x^2 - 1 > 1 \iff x^2 - 2 > 0 \iff (x - \sqrt{2})(x + \sqrt{2}) > 0$. Le trinôme est strictement positif à l'extérieur de ses racines $-\sqrt{2}$ et $\sqrt{2}$ :
    \[
        \boxed{\dom f = \,]-\infty, -\sqrt{2}[ \,\cup\, ]\sqrt{2}, +\infty[}.
    \]

    \item \textbf{Septembre 2024, Q1.2 :} \emph{Domaine de définition de $f(x) = \dfrac{1}{\sqrt{\ln(x^2 - 2x + 1)}}$.}\\
    On impose $\ln(x^2 - 2x + 1) > 0 \iff x^2 - 2x + 1 > 1 \iff x^2 - 2x > 0 \iff x(x - 2) > 0$ :
    \[
        \boxed{\dom f = \,]-\infty, 0[ \,\cup\, ]2, +\infty[}.
    \]

    \item \textbf{Juillet 2025, Q1.2 :} \emph{Conditions d'existence de l'inéquation $\dfrac{\ln(x + 1)}{\sqrt{x(x - 10)}} \ge 0$.}\\
    Ici, le logarithme et la racine sont séparés :
    \begin{itemize}
        \item Au numérateur, $\ln(x + 1)$ existe ssi $x + 1 > 0 \iff x \in \,]-1, +\infty[$.
        \item Au dénominateur, $\sqrt{x(x - 10)}$ est réel et non nul ssi $x(x - 10) > 0 \iff x \in \,]-\infty, 0[ \,\cup\, ]10, +\infty[$.
    \end{itemize}
    On prend l'\textbf{intersection} des deux ensembles : $\bigl]-1, +\infty\bigr[ \cap \bigl(]-\infty, 0[ \,\cup\, ]10, +\infty[\bigr)$ :
    \[
        \boxed{\text{C.E.} : x \in \,]-1, 0[ \,\cup\, ]10, +\infty[}.
    \]

    \item \textbf{Septembre 2025, Q1.2 :} \emph{Domaine de définition de $f(x) = \dfrac{1}{\sqrt{\ln(x^2 - 3)}}$.}\\
    On impose $\ln(x^2 - 3) > 0 \iff x^2 - 3 > 1 \iff x^2 - 4 > 0 \iff (x - 2)(x + 2) > 0$ :
    \[
        \boxed{\dom f = \,]-\infty, -2[ \,\cup\, ]2, +\infty[}.
    \]

    \item \textbf{Juillet 2026, Q1.2 :} \emph{Domaine de définition de $f(x) = \sqrt{\ln(4 - x^2)}$.}\\
    Attention : ici la racine est au \textbf{numérateur} ! La condition est donc large ($\ge 0$) :
    \[
        \ln(4 - x^2) \ge 0 \iff 4 - x^2 \ge 1 \iff x^2 - 3 \le 0 \iff (x + \sqrt{3})(x - \sqrt{3}) \le 0,
    \]
    ce qui donne l'intervalle fermé entre les racines :
    \[
        \boxed{\dom f = \bigl[-\sqrt{3},\, \sqrt{3}\bigr]}.
    \]

    \item \textbf{Septembre 2026, Q1.2 :} \emph{Domaine de définition de $f(x) = \sqrt{\ln\left(\dfrac{x - 1}{x + 2}\right)}$.}\\
    La racine au numérateur impose $\ln\left(\dfrac{x - 1}{x + 2}\right) \ge 0 \iff \dfrac{x - 1}{x + 2} \ge 1$ (avec $x \neq -2$).
    \textbf{Attention :} On ne multiplie surtout pas par $(x + 2)$ dont le signe est inconnu ! On soustrait $1$ et on met au même dénominateur :
    \[
        \frac{x - 1}{x + 2} - 1 \ge 0 \iff \frac{(x - 1) - (x + 2)}{x + 2} \ge 0 \iff \frac{-3}{x + 2} \ge 0.
    \]
    Comme le numérateur $-3$ est strictement négatif, la fraction ne peut jamais s'annuler et est positive si et seulement si son dénominateur est strictement négatif : $x + 2 < 0 \iff x < -2$. D'où :
    \[
        \boxed{\dom f = \,]-\infty, -2[}.
    \]
\end{enumerate}
\end{exemple}

% ----------------------------------------------------------------------------
\section{Exercices d'entraînement du Chapitre 2}
% ----------------------------------------------------------------------------

\emph{Les corrigés détaillés de tous ces exercices se trouvent au Chapitre~10.}

\begin{enumerate}[label=\textbf{Exercice 2.\arabic* :}, leftmargin=*]
    \item \textbf{(Équation avec valeurs absolues)} Résoudre dans $\R$ l'équation $|2x - 3| = |x + 4|$.
    \item \textbf{(Inéquation avec valeur absolue)} Résoudre dans $\R$ l'inéquation $|x^2 - 5x + 4| < 2$.
    \item \textbf{(Équations irrationnelles)} Résoudre dans $\R$ les équations suivantes en vérifiant soigneusement l'absence de solutions parasites :
    \begin{enumerate}[label=(\alph*)]
        \item $\sqrt{2x + 3} = x$
        \item $\sqrt{x + 4} + x = 2$
        \item $\sqrt{3x + 1} - \sqrt{x - 1} = 2$
    \end{enumerate}
    \item \textbf{(Domaines de définition type Question 1)} Déterminer le domaine de définition dans $\R$ de chacune des fonctions suivantes :
    \begin{enumerate}[label=(\alph*)]
        \item $f_1(x) = \dfrac{1}{\sqrt{\ln(x^2 - 4x + 5)}}$
        \item $f_2(x) = \sqrt{\ln\left(\dfrac{2x + 1}{x - 1}\right)}$
        \item $f_3(x) = \dfrac{\sqrt{9 - x^2}}{\ln(x + 2)}$
    \end{enumerate}
\end{enumerate}

% ============================================================================
% CHAPITRE 3 : EXPONENTIELLES ET LOGARITHMES : CALCULS, ÉQUATIONS ET BASES VARIABLES
% ============================================================================
\chapter{Exponentielles et Logarithmes : Calculs, Équations et Bases Variables}

Chaque questionnaire d'Algèbre de l'examen d'admission comporte à la Question~1 au moins une équation exponentielle ou une équation logarithmique (et parfois les deux), tandis que la Question~2 porte fréquemment sur une inéquation logarithmique avec paramètre (que nous approfondirons au Chapitre~7). L'examen se déroulant \textbf{sans calculatrice}, toutes ces questions sont conçues pour se résoudre exactement grâce aux règles algébriques des puissances, des logarithmes et de la formule de changement de base.

% ----------------------------------------------------------------------------
\section{Fonctions exponentielles et règles des puissances}
% ----------------------------------------------------------------------------

\subsection{Définition, propriétés algébriques et monotonie de $a^x$}

\begin{definition}[Exponentielle de base $a$]
Soit $a \in \R_0^+$ un réel strictement positif et différent de $1$ ($a > 0, a \neq 1$). La fonction \textbf{exponentielle de base $a$} associe à tout réel $x \in \R$ le réel strictement positif $a^x = e^{x \ln a} > 0$.
\end{definition}

\begin{theoreme}[Lois fondamentales des exposants]
Pour toutes bases $a > 0$, $b > 0$ et tous exposants réels $x, y \in \R$ :
\[
    \boxed{a^0 = 1, \;\; a^1 = a, \;\; a^{-x} = \frac{1}{a^x}, \;\; a^{x+y} = a^x \cdot a^y, \;\; a^{x-y} = \frac{a^x}{a^y}, \;\; \bigl(a^x\bigr)^y = a^{x \cdot y}, \;\; (a\cdot b)^x = a^x \cdot b^x}.
\]
Pour les exposants fractionnaires ($n \in \N_0$, $p \in \Z$) :
\[
    \boxed{a^{\frac{1}{n}} = \sqrt[n]{a} \qquad \text{et} \qquad a^{\frac{p}{n}} = \sqrt[n]{a^p} = \bigl(\sqrt[n]{a}\bigr)^p}.
\]
\end{theoreme}

\begin{theoreme}[Positivité, injectivité et sens de variation de $x \mapsto a^x$]
\begin{enumerate}
    \item \textbf{Positivité stricte :} Pour tout $x \in \R$, on a toujours $\boxed{a^x > 0}$. En particulier, une équation de la forme $a^x = k$ avec $k \le 0$ n'admet \textbf{aucune solution réelle}.
    \item \textbf{Injectivité (égalité de deux puissances de même base $a > 0, a \neq 1$) :}
    \[
        \boxed{a^{u(x)} = a^{v(x)} \iff u(x) = v(x)}.
    \]
    \item \textbf{Sens de variation (pour les inéquations exponentielles) :}
    \begin{itemize}
        \item Si la base est \textbf{supérieure à $1$ ($a > 1$)}, $x \mapsto a^x$ est \textbf{strictement croissante} (l'ordre est \textbf{conservé}) :
        \[
            \boxed{a^{u(x)} < a^{v(x)} \iff u(x) < v(x) \qquad (\text{pour } a > 1)}.
        \]
        \item Si la base est \textbf{comprise entre $0$ et $1$ ($0 < a < 1$)}, $x \mapsto a^x$ est \textbf{strictement décroissante} (l'ordre est \textbf{inversé}) :
        \[
            \boxed{a^{u(x)} < a^{v(x)} \iff u(x) > v(x) \qquad (\text{pour } 0 < a < 1)}.
        \]
    \end{itemize}
\end{enumerate}
\end{theoreme}

% ----------------------------------------------------------------------------
\section{Logarithmes en base quelconque et changement de base}
% ----------------------------------------------------------------------------

\subsection{Définition et identités de réciprocité}

\begin{definition}[Logarithme en base $a$ ($a > 0, a \neq 1$)]
Pour tout réel \textbf{strictement positif} $y > 0$, le \textbf{logarithme en base $a$ de $y$}, noté $\log_a(y)$, est l'unique exposant réel $x$ auquel il faut élever la base $a$ pour obtenir $y$ :
\[
    \boxed{x = \log_a(y) \iff a^x = y \qquad (\text{avec } y > 0 \text{ et } a > 0, a \neq 1)}.
\]
\begin{itemize}
    \item Lorsque la base est le nombre d'Euler $e \approx 2,71828$, on parle de \textbf{logarithme népérien}, noté $\ln(y) = \log_e(y)$.
    \item Lorsque la base est $10$, on parle de \textbf{logarithme décimal}, souvent noté $\log_{10}(y)$ ou simplement $\log(y)$.
\end{itemize}
\end{definition}

\begin{theoreme}[Valeurs particulières et identités de réciprocité]
Pour toute base $a > 0, a \neq 1$ :
\[
    \boxed{\log_a(1) = 0, \qquad \log_a(a) = 1, \qquad \log_a(a^x) = x \quad (\forall x \in \R), \qquad a^{\log_a(y)} = y \quad (\forall y > 0)}.
\]
\end{theoreme}

\subsection{Propriétés algébriques et formule de changement de base}

\begin{theoreme}[Règles de calcul des logarithmes]
Pour toute base $a > 0, a \neq 1$, tous réels \textbf{strictement positifs} $u > 0, v > 0$ et tout exposant $p \in \R$ :
\begin{enumerate}
    \item \textbf{Logarithme d'un produit :} $\boxed{\log_a(u \cdot v) = \log_a(u) + \log_a(v)}$.
    \item \textbf{Logarithme d'un inverse et d'un quotient :}
    \[
        \boxed{\log_a\left(\frac{1}{v}\right) = -\log_a(v)} \qquad \text{et} \qquad \boxed{\log_a\left(\frac{u}{v}\right) = \log_a(u) - \log_a(v)}.
    \]
    \item \textbf{Logarithme d'une puissance et d'une racine :}
    \[
        \boxed{\log_a\bigl(u^p\bigr) = p \cdot \log_a(u)} \qquad \text{et en particulier } \log_a\bigl(\sqrt{u}\bigr) = \frac{1}{2}\log_a(u).
    \]
    \item \textbf{Formule de changement de base (passage de la base $a$ à une base $b$ quelconque, ou à $\ln$) :}
    \[
        \boxed{\log_a(u) = \frac{\ln(u)}{\ln(a)} = \frac{\log_b(u)}{\log_b(a)}}, \qquad \boxed{\log_a(b) \cdot \log_b(a) = 1} \quad \text{et} \quad \boxed{\log_{a^k}(u) = \frac{1}{k}\log_a(u)}.
    \]
\end{enumerate}
\end{theoreme}

\begin{attention}[Trois pièges mortels sur les logarithmes !]
\begin{enumerate}
    \item \textbf{Il n'existe aucune formule pour le logarithme d'une somme :} $\ln(u + v) \neq \ln(u) + \ln(v)$ et $\ln(u)\cdot\ln(v) \neq \ln(u+v)$ !
    \item \textbf{Ne confondez pas $\ln(x^2)$ et $(\ln x)^2 = \ln^2(x)$ :}
    \begin{itemize}
        \item Dans $\ln(x^2)$, le carré porte sur $x$ : pour $x > 0$, $\ln(x^2) = 2\ln(x)$ (et pour $x \neq 0$ quelconque, $\ln(x^2) = 2\ln|x|$).
        \item Dans $(\ln x)^2$, c'est tout le logarithme qui est élevé au carré : $(\ln x)^2 = (\ln x)\cdot(\ln x) \neq 2\ln x$.
    \end{itemize}
    \item \textbf{Changement de sens dans une inéquation logarithmique lorsque la base $a$ est entre $0$ et $1$ :}
    \begin{center}
    \begin{tikzpicture}[scale=0.95, >=Stealth]
        % Base a > 1
        \begin{scope}[shift={(0,0)}]
            \draw[->, thick, darkgray] (-0.3,0) -- (3.4,0) node[right] {\small $x$};
            \draw[->, thick, darkgray] (0,-1.4) -- (0,1.6) node[above] {\small $\log_a(x)$};
            \draw[very thick, uclblue, domain=0.25:3.2, samples=60] plot (\x, {ln(\x)/ln(2)});
            \fill[warnred] (1,0) circle (2pt) node[below right] {\small $1$};
            \node[uclblue, font=\small\bfseries] at (2.2,1.35) {Base $a > 1$ (croissante)};
            \node[defgreen, font=\footnotesize] at (2.3,0.4) {$\log_a(x) > 0 \iff x > 1$};
            \node[warnred, font=\footnotesize] at (1.5,-0.9) {$\log_a(x) < 0 \iff 0 < x < 1$};
        \end{scope}
        % Base 0 < a < 1
        \begin{scope}[shift={(7.5,0)}]
            \draw[->, thick, darkgray] (-0.3,0) -- (3.4,0) node[right] {\small $x$};
            \draw[->, thick, darkgray] (0,-1.4) -- (0,1.6) node[above] {\small $\log_a(x)$};
            \draw[very thick, exorange, domain=0.25:3.2, samples=60] plot (\x, {-ln(\x)/ln(2)});
            \fill[warnred] (1,0) circle (2pt) node[above right] {\small $1$};
            \node[exorange, font=\small\bfseries] at (2.2,1.35) {Base $0 < a < 1$ (décroissante)};
            \node[defgreen, font=\footnotesize] at (1.6,0.95) {$\log_a(x) > 0 \iff 0 < x < 1$};
            \node[warnred, font=\footnotesize] at (2.3,-0.45) {$\log_a(x) < 0 \iff x > 1$};
        \end{scope}
    \end{tikzpicture}
    \end{center}
\end{enumerate}
\end{attention}

\begin{exemple}[Simplification par changement de base -- UCLouvain, Juillet 2021, Q1.1]
\textbf{Énoncé :} Simplifier l'opération $\log_a(27)\cdot\log_3(a)$ (où $a > 0, a \neq 1$).

\textbf{Résolution :} Passons tout en logarithme népérien (ou en base $3$) grâce à la formule de changement de base et au fait que $27 = 3^3$ :
\[
    \log_a(27)\cdot\log_3(a) = \frac{\ln(27)}{\ln(a)} \cdot \frac{\ln(a)}{\ln(3)} = \frac{\ln(3^3)}{\ln(3)} = \frac{3\ln(3)}{\ln(3)} = \boxed{3}.
\]
\end{exemple}

% ----------------------------------------------------------------------------
\section{Méthodes de résolution des équations exponentielles}
% ----------------------------------------------------------------------------

Les équations exponentielles posées à l'examen d'entrée de l'UCLouvain se répartissent en trois grandes familles que l'on reconnaît au premier coup d'œil :

\begin{astuce}[Les 3 familles d'équations exponentielles à l'examen]
\begin{enumerate}
    \item \textbf{Type 1 -- Produit de puissances égal à une constante ($a^{x+p}\cdot b^{x+q} = C$) :} On utilise $a^{x+p} = a^p \cdot a^x$ pour isoler $(a\cdot b)^x$, ou on décompose la constante $C$ en facteurs premiers ($C = a^\alpha \cdot b^\beta$).
    \item \textbf{Type 2 -- Équation trinôme en $a^x$ ($A \cdot a^{2x} + B \cdot a^x + C = 0$ ou $A \cdot a^x + B \cdot a^{-x} = C$) :}
    On pose l'\textbf{inconnue auxiliaire strictement positive $\boxed{t = a^x > 0}$} (avec $a^{2x} = (a^x)^2 = t^2$ et $a^{-x} = \frac{1}{t}$). On résout l'équation du second degré en $t$, on \textbf{rejette toute racine $t \le 0$}, puis pour chaque racine $t > 0$ on résout $a^x = t \iff x = \log_a(t)$.
    \item \textbf{Type 3 -- Puissances de logarithmes ($a^{\log_b(u)} = c$) :} On réécrit la base $a$ comme une puissance de $b$ (par exemple $4^{\log_2(6)} = (2^2)^{\log_2(6)} = \bigl(2^{\log_2(6)}\bigr)^2 = 6^2 = 36$).
\end{enumerate}
\end{astuce}

\begin{exemple}[Les cinq équations exponentielles des annales (2023--2026)]
\begin{enumerate}
    \item \textbf{Juillet 2023 -- Série 1, Q1.4 :} \emph{Résoudre dans $\R$ l'équation $4^{\log_2(6)} - 25^{\log_5(3x)} = 0$.}\\
    Condition d'existence du logarithme : $3x > 0 \iff x > 0$.
    Remarquons que $4 = 2^2$ et $25 = 5^2$. Grâce à l'identité $b^{\log_b(y)} = y$ :
    \begin{align*}
        4^{\log_2(6)} &= (2^2)^{\log_2(6)} = \bigl(2^{\log_2(6)}\bigr)^2 = 6^2 = 36, \\
        25^{\log_5(3x)} &= (5^2)^{\log_5(3x)} = \bigl(5^{\log_5(3x)}\bigr)^2 = (3x)^2 = 9x^2.
    \end{align*}
    L'équation se réduit donc à $36 - 9x^2 = 0 \iff x^2 = 4 \iff x = \pm 2$. Comme la C.E. impose $x > 0$, on rejette $-2$ et il reste :
    \[
        \boxed{x = 2}.
    \]

    \item \textbf{Septembre 2023, Q1.4 :} \emph{Résoudre dans $\R$ l'équation $3^{2x+1} - 8\cdot 3^{x+1} = 3^3$.}\\
    Détachons le $+1$ dans les exposants ($3^{2x+1} = 3\cdot (3^x)^2$ et $3^{x+1} = 3\cdot 3^x$) puis divisons toute l'équation par $3$ :
    \[
        3\cdot(3^x)^2 - 8\cdot 3\cdot 3^x = 27 \iff (3^x)^2 - 8\cdot 3^x - 9 = 0.
    \]
    Posons $t = 3^x > 0$ : l'équation devient $t^2 - 8t - 9 = 0 \iff (t - 9)(t + 1) = 0$, soit $t = 9$ ou $t = -1$.
    Puisque $t = 3^x > 0$, la valeur $t = -1$ est impossible. Il reste $3^x = 9 = 3^2 \iff \boxed{x = 2}$.

    \item \textbf{Juillet 2025, Q1.3 :} \emph{Résolvez dans $\R$ l'équation $16^x - 4^x = 2$.}\\
    Comme $16^x = (4^2)^x = (4^x)^2$, posons $t = 4^x > 0$. On obtient :
    \[
        t^2 - t - 2 = 0 \iff (t - 2)(t + 1) = 0 \implies t = 2 \quad (\text{car } t > 0).
    \]
    On revient à $x$ : $4^x = 2 \iff (2^2)^x = 2^{2x} = 2^1 \iff 2x = 1 \iff \boxed{x = \dfrac{1}{2}}$.

    \item \textbf{Septembre 2025, Q1.1 :} \emph{Résoudre dans $\R$ l'équation $9^x + 1 = 10\cdot 3^{x-1}$.}\\
    Écrivons $9^x = (3^x)^2$ et $3^{x-1} = \frac{3^x}{3}$, puis multiplions par $3$ et posons $t = 3^x > 0$ :
    \[
        3t^2 - 10t + 3 = 0 \iff (3t - 1)(t - 3) = 0 \iff t = \frac{1}{3} = 3^{-1} \;\text{ ou }\; t = 3 = 3^1.
    \]
    Les deux racines en $t$ sont strictement positives ! Elles donnent toutes deux une solution en $x$ :
    \[
        3^x = 3^{-1} \;\text{ ou }\; 3^x = 3^1 \iff \boxed{x = -1 \quad \text{ou} \quad x = 1 \quad (x = \pm 1)}.
    \]

    \item \textbf{Juillet 2026, Q1.1 :} \emph{Résoudre dans $\R$ l'équation $3^{x+1}\cdot 5^{x-1} + 1 = 2026$.}\\
    En soustrayant $1$ aux deux membres : $3^{x+1}\cdot 5^{x-1} = 2025$.
    \begin{itemize}
        \item \textbf{Méthode 1 (Décomposition en facteurs premiers de $2025$) :} $2025 = 45^2 = (9 \times 5)^2 = (3^2 \times 5)^2 = 3^4 \cdot 5^2$.
        En identifiant les exposants de $3$ et de $5$ dans $3^{x+1}\cdot 5^{x-1} = 3^4 \cdot 5^2$, on voit immédiatement que $x + 1 = 4$ et $x - 1 = 2$, ce qui donne $\boxed{x = 3}$.
        \item \textbf{Méthode 2 (Regroupement sous un seul exposant $x$) :}
        \[
            3^{x+1}\cdot 5^{x-1} = \frac{3}{5}\cdot 15^x = 2025 \iff 15^x = \frac{5 \times 2025}{3} = 3375 = 15^3 \iff \boxed{x = 3}.
        \]
    \end{itemize}

    \item \textbf{Septembre 2026, Q1.1 :} \emph{Trouvez tous les réels $x$ tels que $2^x + 2^{3-x} = 6$.}\\
    Écrivons $2^{3-x} = \dfrac{2^3}{2^x} = \dfrac{8}{2^x}$ et posons $t = 2^x > 0$ :
    \[
        t + \frac{8}{t} = 6 \iff t^2 - 6t + 8 = 0 \iff (t - 2)(t - 4) = 0 \iff t = 2 \;\text{ ou }\; t = 4.
    \]
    En revenant à $2^x = t$ ($2^x = 2^1$ ou $2^x = 2^2$), on obtient les deux solutions réelles :
    \[
        \boxed{x = 1 \quad \text{ou} \quad x = 2}.
    \]
\end{enumerate}
\end{exemple}

% ----------------------------------------------------------------------------
\section{Méthodes de résolution des équations logarithmiques}
% ----------------------------------------------------------------------------

\begin{astuce}[Méthodologie en 3 étapes pour toute équation logarithmique]
\begin{enumerate}
    \item \textbf{Étape 1 -- Conditions d'existence (C.E.) :} Avant toute transformation algébrique, imposez que l'argument de \textbf{chaque} logarithme présent dans l'énoncé initial soit \textbf{strictement positif} (le domaine de $\log_a(u) + \log_a(v)$ est $\{u > 0 \text{ et } v > 0\}$, plus restreint que celui de $\log_a(u\cdot v)$ !).
    \item \textbf{Étape 2 -- Réduction :}
    \begin{itemize}
        \item \emph{Cas A (Combinaison linéaire de logarithmes) :} Utilisez $\log_a(u) + \log_a(v) = \log_a(uv)$ et $\log_a(u) - \log_a(v) = \log_a(u/v)$ pour ramener l'équation à la forme $\log_a\bigl(A(x)\bigr) = \log_a\bigl(B(x)\bigr) \iff A(x) = B(x)$ (ou $\log_a\bigl(A(x)\bigr) = k \iff A(x) = a^k$).
        \item \emph{Cas B (Polynôme en $\log_a(x)$ avec un terme $(\log_a x)^2$) :} Posez l'inconnue auxiliaire $t = \log_a(x) \in \R$ (attention, ici $t$ peut être de signe quelconque dans $\R$ !), résolvez en $t$, puis revenez à $x = a^t$.
    \end{itemize}
    \item \textbf{Étape 3 -- Filtrage par les C.E. :} Testez chaque solution candidate dans les conditions d'existence de l'Étape~1 et éliminez les valeurs interdites.
\end{enumerate}
\end{astuce}

\begin{exemple}[Les cinq équations logarithmiques des annales (2021--2024)]
\begin{enumerate}
    \item \textbf{Juillet 2022, Q1.1 :} \emph{Résoudre dans $\R$ l'équation $\log_4(8x^2) + \log_4(4x) = 1$.}\\
    \textbf{C.E. :} $8x^2 > 0$ et $4x > 0 \iff \mathbf{x > 0}$.
    Regroupons la somme des deux logarithmes en un seul logarithme de produit :
    \[
        \log_4\bigl(8x^2 \cdot 4x\bigr) = 1 \iff \log_4\bigl(32x^3\bigr) = 1 \iff 32x^3 = 4 \iff x^3 = \frac{1}{8} = \left(\frac{1}{2}\right)^3.
    \]
    On en tire $x = \frac{1}{2}$, qui vérifie bien la C.E. $x > 0$. Donc $\boxed{x = \dfrac{1}{2}}$.

    \item \textbf{Juillet 2023 -- Série 2, Q1.4 :} \emph{Résoudre dans $\R$ l'équation $\log_5(x) + \log_5(x - 4) = 1$.}\\
    \textbf{C.E. :} $x > 0$ et $x - 4 > 0 \iff \mathbf{x > 4}$.
    En regroupant :
    \[
        \log_5\bigl(x(x - 4)\bigr) = 1 \iff x(x - 4) = 5 \iff x^2 - 4x - 5 = 0 \iff (x - 5)(x + 1) = 0.
    \]
    On obtient deux candidats : $x = 5$ et $x = -1$. La condition d'existence $x > 4$ \textbf{élimine $x = -1$} (pour lequel $\log_5(-1)$ n'existe pas dans $\R$ !) et conserve l'unique solution :
    \[
        \boxed{x = 5}.
    \]

    \item \textbf{Juillet 2024, Q1.2 :} \emph{Résoudre dans $\R$ l'équation $\log x^2 - \log 2x + \log 3 = \log 4$ (où $\log$ est en base $10$).}\\
    \textbf{C.E. :} $x^2 > 0$ et $2x > 0 \iff \mathbf{x > 0}$.
    En combinant les logarithmes du membre de gauche ($\log u - \log v + \log w = \log\frac{uw}{v}$) :
    \[
        \log\left(\frac{3x^2}{2x}\right) = \log(4) \iff \frac{3x}{2} = 4 \iff \boxed{x = \frac{8}{3}} \quad (\text{qui vérifie bien } x > 0).
    \]

    \item \textbf{Septembre 2024, Q1.6 :} \emph{Résolvez dans $\R$ l'équation $\log(x^2) + \log x - (\log x)^2 = \log(100)$ (où $\log$ est le logarithme en base $10$).}\\
    \textbf{C.E. :} $x > 0$. Pour $x > 0$, on a $\log(x^2) = 2\log x$ et $\log(100) = \log(10^2) = 2$. L'équation devient :
    \[
        2\log x + \log x - (\log x)^2 = 2 \iff (\log x)^2 - 3\log x + 2 = 0.
    \]
    Posons $t = \log x \in \R$ : on obtient $t^2 - 3t + 2 = 0 \iff (t - 1)(t - 2) = 0 \iff t = 1 \text{ ou } t = 2$.
    En revenant à $x = 10^t$ :
    \[
        \log_{10}(x) = 1 \iff x = 10^1 = 10 \qquad \text{ou} \qquad \log_{10}(x) = 2 \iff x = 10^2 = 100.
    \]
    Les deux valeurs sont strictement positives : $\boxed{x = 10 \quad \text{et} \quad x = 100}$.

    \item \textbf{Septembre 2021, Q2.2 :} \emph{Résoudre dans $\R$ l'équation $\log_5(1 + 5^x) = \log_5(10) - x - 1$.}\\
    \textbf{C.E. :} Comme $5^x > 0$ pour tout $x \in \R$, on a toujours $1 + 5^x > 1 > 0$, donc $\dom = \R$.
    Exprimons $x + 1$ comme un logarithme en base $5$ : $x + 1 = \log_5\bigl(5^{x+1}\bigr)$. Le membre de droite s'écrit :
    \[
        \log_5(10) - \log_5\bigl(5^{x+1}\bigr) = \log_5\left(\frac{10}{5^{x+1}}\right) = \log_5\left(\frac{2}{5^x}\right).
    \]
    Par injectivité de $\log_5$, l'équation équivaut à :
    \[
        1 + 5^x = \frac{2}{5^x} \iff (5^x)^2 + 5^x - 2 = 0.
    \]
    En posant $y = 5^x > 0$, on a $y^2 + y - 2 = (y - 1)(y + 2) = 0$. Comme $y > 0$, on rejette $y = -2$ et on garde $y = 1$, soit $5^x = 1 \iff \boxed{x = 0}$.
\end{enumerate}
\end{exemple}

% ----------------------------------------------------------------------------
\section{Exercices d'entraînement du Chapitre 3}
% ----------------------------------------------------------------------------

\emph{Les corrigés détaillés de tous ces exercices se trouvent au Chapitre~10.}

\begin{enumerate}[label=\textbf{Exercice 3.\arabic* :}, leftmargin=*]
    \item \textbf{(Calculs logarithmiques sans calculatrice)} Simplifier les nombres suivants sous forme d'un entier ou d'une fraction irréductible :
    \begin{enumerate}[label=(\alph*)]
        \item $A = \log_2(24) - \log_2(3) + \log_4(9) - \log_2(\sqrt{3})$
        \item $B = \log_3(5)\cdot\log_5(7)\cdot\log_7(81)$
        \item $C = 9^{\log_3(5)} - 8^{\log_2(3)}$
    \end{enumerate}
    \item \textbf{(Équations exponentielles type Question 1)} Résoudre dans $\R$ les équations suivantes :
    \begin{enumerate}[label=(\alph*)]
        \item $4^{x+1} - 9\cdot 2^x + 2 = 0$
        \item $3^x + 3^{2-x} = 10$
        \item $2^{x+2}\cdot 3^{x-1} = 144$
    \end{enumerate}
    \item \textbf{(Équations logarithmiques type Question 1)} Résoudre dans $\R$ les équations suivantes :
    \begin{enumerate}[label=(\alph*)]
        \item $\log_3(x + 2) + \log_3(x - 4) = 3$
        \item $\ln^2(x) - \ln(x^4) + 3 = 0$
        \item $\log_2(x) + \log_x(2) = \dfrac{5}{2}$
    \end{enumerate}
\end{enumerate}

% ============================================================================
% CHAPITRE 4 : NOMBRES COMPLEXES DANS C
% ============================================================================
\chapter{Nombres Complexes dans $\C$}

Les nombres complexes sont systématiquement présents à l'examen d'admission d'Algèbre de l'UCLouvain (chaque session de 2021 à 2026 comporte une sous-question de la Question~1 sur $\C$, et parfois une question complète en Question~2 comme en Juillet 2021). Le programme exige une maîtrise parfaite du passage entre la \textbf{forme cartésienne (algébrique)} $z = a + bi$ et la \textbf{forme trigonométrique (polaire)} $z = \rho(\cos\theta + i\sin\theta) = \rho e^{i\theta}$, ainsi que la résolution d'équations avec le conjugué $z^*$ et le calcul des racines $n$-ièmes d'un nombre complexe.

% ----------------------------------------------------------------------------
\section{Forme cartésienne (algébrique), conjugué et opérations dans $\C$}
% ----------------------------------------------------------------------------

\subsection{L'unité imaginaire $i$ et l'écriture $z = a + bi$}

\begin{definition}[Nombre complexe, partie réelle et partie imaginaire]
L'ensemble $\C$ des \textbf{nombres complexes} est constitué de tous les nombres de la forme :
\[
    \boxed{z = a + bi \qquad \text{avec } a \in \R,\; b \in \R \text{ et } i^2 = -1}.
\]
\begin{itemize}
    \item Cette écriture unique s'appelle la \textbf{forme cartésienne} (ou \emph{forme algébrique}, ou \emph{forme rectangulaire}) de $z$.
    \item Le réel $a = \Reop(z)$ est la \textbf{partie réelle} de $z$, et le réel $b = \Imop(z)$ est la \textbf{partie imaginaire} de $z$ (\textbf{attention :} la partie imaginaire $\Imop(z) = b$ est un \emph{nombre réel}, sans le $i$ !).
    \item Si $b = 0$, $z = a \in \R$ est un \textbf{réel}. Si $a = 0$, $z = bi$ est dit \textbf{imaginaire pur} (ensemble $i\R$).
\end{itemize}
\end{definition}

\begin{theoreme}[Égalité de deux nombres complexes et puissances cycliques de $i$]
\begin{enumerate}
    \item \textbf{Principe d'identification (parties réelle et imaginaire) :} Deux nombres complexes $a + bi$ et $c + di$ (avec $a,b,c,d \in \R$) sont égaux \textbf{si et seulement si} leurs parties réelles sont égales \textbf{et} leurs parties imaginaires sont égales :
    \[
        \boxed{a + bi = c + di \iff \begin{cases} a = c \\ b = d \end{cases}} \quad \text{et en particulier } \; a + bi = 0 \iff a = 0 \text{ et } b = 0.
    \]
    \item \textbf{Cycle de période 4 des puissances de $i$ :} Pour tout entier $k \in \Z$ :
    \[
        \boxed{i^{4k} = 1, \qquad i^{4k+1} = i, \qquad i^{4k+2} = -1, \qquad i^{4k+3} = -i \qquad \left(\text{et } \frac{1}{i} = -i\right)}.
    \]
\end{enumerate}
\end{theoreme}

\subsection{Nombre complexe conjugué $z^*$ et division en forme cartésienne}

\begin{definition}[Conjugué $z^*$ d'un nombre complexe]
Dans les énoncés d'examen de l'UCLouvain, le \textbf{conjugué} du nombre complexe $z = a + bi$ ($a, b \in \R$) est noté $z^*$ (on rencontre aussi la notation $\bar{z}$ dans la littérature). Il est défini en changeant le signe de la partie imaginaire :
\[
    \boxed{z^* = a - bi}.
\]
Géométriquement, dans le plan complexe de Gauss, le point image de $z^*$ est le \textbf{symétrique} du point image de $z$ par rapport à l'\textbf{axe réel horizontal}.
\end{definition}

\begin{theoreme}[Propriétés algébriques du conjugué $z^*$]
Pour tous nombres complexes $z = a + bi$ et $w \in \C$, et tout entier $n \in \Z$ :
\begin{enumerate}
    \item \textbf{Somme et différence de $z$ et $z^*$ :}
    \[
        \boxed{z + z^* = 2a = 2\Reop(z) \in \R} \qquad \text{et} \qquad \boxed{z - z^* = 2bi = 2i\Imop(z)}.
    \]
    En particulier : $z \in \R \iff z^* = z$, et $z$ est imaginaire pur $\iff z^* = -z$.
    \item \textbf{Produit d'un complexe par son conjugué (toujours réel positif !) :}
    \[
        \boxed{z \cdot z^* = (a + bi)(a - bi) = a^2 - (bi)^2 = a^2 + b^2 = |z|^2 \ge 0}.
    \]
    \item \textbf{Compatibilité avec toutes les opérations :}
    \[
        (z^*)^* = z, \quad (z \pm w)^* = z^* \pm w^*, \quad (z \cdot w)^* = z^* \cdot w^*, \quad \left(\frac{z}{w}\right)^* = \frac{z^*}{w^*}, \quad (z^n)^* = (z^*)^n.
    \]
\end{enumerate}
\end{theoreme}

\begin{astuce}[Comment diviser deux nombres complexes en forme cartésienne ?]
Pour mettre un quotient $\dfrac{z_1}{z_2} = \dfrac{a + bi}{c + di}$ sous forme cartésienne $X + Yi$, \textbf{multipliez le numérateur et le dénominateur par le conjugué $z_2^* = c - di$ du dénominateur} afin de rendre le dénominateur réel ($c^2 + d^2$) :
\[
    \boxed{\frac{a + bi}{c + di} = \frac{(a + bi)(c - di)}{(c + di)(c - di)} = \frac{(ac + bd) + (bc - ad)i}{c^2 + d^2} = \frac{ac + bd}{c^2 + d^2} + \frac{bc - ad}{c^2 + d^2}\,i}.
\]
\end{astuce}

\begin{exemple}[Calculs avec le conjugué et équations cartésiennes -- UCLouvain, Juillet 2021 Q1.7, Septembre 2023 Q1.1 \& Juillet 2024 Q1.7]
\begin{enumerate}
    \item \textbf{Juillet 2021, Q1.7 :} \emph{Soit $z = 1 - i$. Donner le module et l'argument du produit de $z$ avec son conjugué.}\\
    Le produit vaut $P = z \cdot z^* = (1 - i)(1 + i) = 1^2 + (-1)^2 = 2$. Comme $2$ est un nombre \textbf{réel strictement positif}, son module vaut $\boxed{|P| = 2}$ et son argument vaut $\boxed{\arg(P) = 0 \pmod{2\pi}}$.
    \item \textbf{Septembre 2023, Q1.1 :} \emph{Soit $z = 3 + 4i$, et $z^*$ le complexe conjugué de $z$. Que vaut $\dfrac{z z^*}{z + z^*}$ ?}\\
    Sans aucun calcul compliqué, utilisons les deux formules du théorème pour $z = a + bi$ ($a = 3, b = 4$) :
    \[
        z z^* = a^2 + b^2 = 3^2 + 4^2 = 9 + 16 = 25 \qquad \text{et} \qquad z + z^* = 2a = 2 \times 3 = 6,
    \]
    d'où immédiatement : $\dfrac{z z^*}{z + z^*} = \boxed{\dfrac{25}{6}}$.
    \item \textbf{Juillet 2024, Q1.7 :} \emph{Résoudre dans $\C$ l'équation $(z - 2)z^* = -iz$. Exprimez les solutions sous forme cartésienne.}\\
    Posons $z = a + bi$ avec $a, b \in \R$ (donc $z^* = a - bi$ et $z z^* = a^2 + b^2$). Développons l'équation $(z - 2)z^* = -iz \iff z z^* - 2z^* = -iz$ :
    \[
        (a^2 + b^2) - 2(a - bi) = -i(a + bi) \iff (a^2 + b^2 - 2a) + 2bi = b - ai.
    \]
    Par identification des parties réelle et imaginaire :
    \[
        \begin{cases} a^2 + b^2 - 2a = b & (1) \\ 2b = -a \iff a = -2b & (2) \end{cases}
    \]
    Remplaçons $a = -2b$ dans l'équation $(1)$ :
    \[
        (-2b)^2 + b^2 - 2(-2b) = b \iff 5b^2 + 3b = 0 \iff b(5b + 3) = 0.
    \]
    On obtient deux valeurs réelles pour $b$ :
    \begin{itemize}
        \item Si $b = 0$, alors $a = -2(0) = 0$, ce qui donne la première solution $z_1 = 0$.
        \item Si $b = -\dfrac{3}{5}$, alors $a = -2\left(-\dfrac{3}{5}\right) = \dfrac{6}{5}$, ce qui donne la seconde solution $z_2 = \dfrac{6}{5} - \dfrac{3}{5}i$.
    \end{itemize}
    L'ensemble des solutions est donc $\boxed{S = \left\{0,\; \frac{6}{5} - \frac{3}{5}i\right\}}$.
\end{enumerate}
\end{exemple}

% ----------------------------------------------------------------------------
\section{Forme trigonométrique (polaire), module, argument et formule de De Moivre}
% ----------------------------------------------------------------------------

\subsection{Module $\rho = |z|$ et argument $\theta = \arg(z)$ dans le plan de Gauss}

Dans le plan complexe muni d'un repère orthonormé $(O; \vec{e}_1, \vec{e}_2)$, tout nombre complexe $z = a + bi$ est représenté par le point $M(a, b)$ (ou le vecteur $\overrightarrow{OM}$). Ses coordonnées polaires $(\rho, \theta)$ définissent son \textbf{module} et son \textbf{argument}.

\begin{center}
\begin{tikzpicture}[scale=1.15, >=Stealth]
    \draw[step=0.5, gray!20, very thin] (-1.2,-0.8) grid (3.3,2.6);
    \draw[->, thick, darkgray] (-1.3,0) -- (3.5,0) node[right] {$\Reop(z)$};
    \draw[->, thick, darkgray] (0,-0.9) -- (0,2.7) node[above] {$\Imop(z)$};

    \coordinate (O) at (0,0);
    \coordinate (M) at (2.6,1.8);
    \coordinate (Mx) at (2.6,0);
    \coordinate (My) at (0,1.8);

    \fill[uclaccent!15] (0,0) -- (0.85,0) arc (0:34.7:0.85) -- cycle;
    \draw[thick, uclaccent, ->] (0.85,0) arc (0:34.7:0.85);
    \node[uclaccent] at (17.5:1.15) {$\theta = \arg(z)$};

    \draw[dashed, gray] (M) -- (Mx);
    \draw[dashed, gray] (M) -- (My);
    \draw[very thick, exorange] (O) -- (Mx) node[midway, below] {$a = \rho\cos\theta$};
    \draw[very thick, defgreen] (O) -- (My) node[midway, left] {$b = \rho\sin\theta$};
    \draw[very thick, uclblue, ->] (O) -- (M) node[midway, above left] {$\rho = |z| = \sqrt{a^2+b^2}$};

    \fill[warnred] (M) circle (2.2pt) node[above right] {$M(z = a+bi = \rho e^{i\theta})$};
\end{tikzpicture}
\end{center}

\begin{definition}[Module, argument et formes trigonométrique / exponentielle]
Soit $z = a + bi \in \C_0$ un nombre complexe non nul :
\begin{enumerate}
    \item Le \textbf{module} de $z$, noté $\rho = |z|$, est la distance $\overline{OM} > 0$ :
    \[
        \boxed{\rho = |z| = \sqrt{a^2 + b^2} = \sqrt{z z^*}}.
    \]
    \item L'\textbf{argument} de $z$, noté $\theta = \arg(z)$ (défini à $2k\pi$ près, $k \in \Z$), est l'angle orienté entre le demi-axe réel positif et le vecteur $\overrightarrow{OM}$, caractérisé par le système :
    \[
        \boxed{\begin{cases} \cos\theta = \dfrac{a}{\rho} = \dfrac{a}{\sqrt{a^2+b^2}} \\[1.5ex] \sin\theta = \dfrac{b}{\rho} = \dfrac{b}{\sqrt{a^2+b^2}} \end{cases}}
    \]
    \item La \textbf{forme trigonométrique} et la \textbf{forme exponentielle (ou polaire)} de $z$ s'écrivent :
    \[
        \boxed{z = \rho(\cos\theta + i\sin\theta) = \rho e^{i\theta}} \qquad \text{où } e^{i\theta} = \cos\theta + i\sin\theta \text{ est la formule d'Euler}.
    \]
\end{enumerate}
\end{definition}

\begin{attention}[Ne calculez JAMAIS $\theta = \arctan(b/a)$ les yeux fermés sans regarder le quadrant !]
La fonction $\arctan(b/a)$ renvoie toujours un angle compris dans $\left]-\frac{\pi}{2}, \frac{\pi}{2}\right[$ (quadrants I et IV). Si la partie réelle $a$ est \textbf{négative} ($a < 0$, quadrants II et III), il faut \textbf{ajouter $\pi$ (ou retrancher $\pi$)} à $\arctan(b/a)$ !
Le moyen le plus sûr et le plus rapide sans calculatrice est d'\textbf{identifier directement le couple $(\cos\theta, \sin\theta) = \left(\frac{a}{\rho}, \frac{b}{\rho}\right)$ sur le cercle trigonométrique} :
\begin{center}
\renewcommand{\arraystretch}{1.45}
\begin{tabular}{@{}lcccccccc@{}}
\toprule
\textbf{Complexe de base} & $1$ & $i$ & $-1$ & $-i$ & $1 + i$ & $1 - i$ & $1 + i\sqrt{3}$ & $\sqrt{3} + i$ \\
\midrule
\textbf{Module $\rho$} & $1$ & $1$ & $1$ & $1$ & $\sqrt{2}$ & $\sqrt{2}$ & $2$ & $2$ \\
\textbf{Argument $\theta$} & $0$ & $\dfrac{\pi}{2}$ & $\pi$ & $-\dfrac{\pi}{2}$ & $\dfrac{\pi}{4}$ & $-\dfrac{\pi}{4}$ & $\dfrac{\pi}{3}$ & $\dfrac{\pi}{6}$ \\
\bottomrule
\end{tabular}
\end{center}
\end{attention}

\subsection{Opérations en forme polaire et formule de De Moivre}

Si la forme cartésienne $a + bi$ est idéale pour les additions et soustractions, la forme polaire $\rho e^{i\theta}$ est d'une puissance redoutable pour les \textbf{produits}, les \textbf{quotients}, les \textbf{puissances} et les \textbf{racines} !

\begin{theoreme}[Produit, quotient et formule de De Moivre en forme polaire]
Soient $z_1 = \rho_1 e^{i\theta_1}$ et $z_2 = \rho_2 e^{i\theta_2}$ deux nombres complexes non nuls, et $n \in \Z$ :
\begin{enumerate}
    \item \textbf{Produit (on multiplie les modules, on additionne les arguments) :}
    \[
        \boxed{z_1 \cdot z_2 = \rho_1 \rho_2 \, e^{i(\theta_1 + \theta_2)}} \implies |z_1 z_2| = |z_1|\cdot|z_2| \;\text{ et }\; \arg(z_1 z_2) = \arg(z_1) + \arg(z_2).
    \]
    \item \textbf{Quotient (on divise les modules, on soustrait les arguments) :}
    \[
        \boxed{\frac{z_1}{z_2} = \frac{\rho_1}{\rho_2} \, e^{i(\theta_1 - \theta_2)}} \implies \left|\frac{z_1}{z_2}\right| = \frac{|z_1|}{|z_2|} \;\text{ et }\; \arg\left(\frac{z_1}{z_2}\right) = \arg(z_1) - \arg(z_2).
    \]
    \item \textbf{Puissance $n$-ième et formule de De Moivre :}
    \[
        \boxed{z^n = \rho^n e^{in\theta} = \rho^n\bigl(\cos(n\theta) + i\sin(n\theta)\bigr)} \implies |z^n| = |z|^n \;\text{ et }\; \arg(z^n) = n\arg(z).
    \]
\end{enumerate}
\end{theoreme}

\begin{astuce}[L'astuce magique $(1 \pm i)^2 = \pm 2i$ pour les puissances de $(1 \pm i)$]
Pour calculer rapidement $(1 + i)^n$, $(1 - i)^n$ ou $(i - 1)^n$ en forme cartésienne sans même passer par la trigonométrie, retenez par cœur que le carré de $(1 \pm i)$ est \textbf{imaginaire pur} :
\[
    \boxed{(1 + i)^2 = 1 + 2i + i^2 = 2i} \qquad \text{et} \qquad \boxed{(1 - i)^2 = (i - 1)^2 = 1 - 2i + i^2 = -2i}.
\]
Ainsi, $(1 + i)^{10} = \bigl((1 + i)^2\bigr)^5 = (2i)^5 = 32 i^5 = 32i$ en trois secondes chrono !
\end{astuce}

\begin{exemple}[Les sept questions de puissances et formes polaires des annales (2022--2026)]
\begin{enumerate}
    \item \textbf{Juillet 2022, Q1.4 :} \emph{Donner le module et l'argument de $z = (1 + \sqrt{3}i)^4$.}\\
    Posons $w = 1 + \sqrt{3}i$. Son module vaut $|w| = \sqrt{1^2 + (\sqrt{3})^2} = \sqrt{4} = 2$, et comme $\cos\theta = \frac{1}{2}$ et $\sin\theta = \frac{\sqrt{3}}{2}$, son argument vaut $\arg(w) = \frac{\pi}{3}$.
    Par la formule de De Moivre pour $z = w^4$ :
    \[
        |z| = |w|^4 = 2^4 = 16 \qquad \text{et} \qquad \arg(z) = 4\arg(w) = \frac{4\pi}{3} \equiv -\frac{2\pi}{3} \pmod{2\pi}.
    \]
    Réponse : $\boxed{\text{module} = 16, \quad \text{argument} = -\dfrac{2\pi}{3} \text{ (ou } \dfrac{4\pi}{3}\text{)}}$.

    \item \textbf{Septembre 2022, Q1.3 :} \emph{Donner le module et l'argument de $z = (1 - i)^6$.}\\
    \textbf{Méthode 1 (Polaire) :} $|1 - i| = \sqrt{2}$ et $\arg(1 - i) = -\frac{\pi}{4}$. Donc $|z| = (\sqrt{2})^6 = 2^3 = 8$ et $\arg(z) = 6\left(-\frac{\pi}{4}\right) = -\frac{3\pi}{2} \equiv \frac{\pi}{2} \pmod{2\pi}$.\\
    \textbf{Méthode 2 (Carré $(1-i)^2 = -2i$) :} $(1 - i)^6 = \bigl((1 - i)^2\bigr)^3 = (-2i)^3 = -8i^3 = 8i = 8e^{i\pi/2}$.\\
    Réponse : $\boxed{\text{module} = 8 \quad \text{et} \quad \text{argument} = \dfrac{\pi}{2}}$.

    \item \textbf{Juillet 2023 -- Série 2, Q1.1 :} \emph{Valeur sous forme cartésienne $(a + bi)$ de $E = (i - 1)^6 + (i - 1)^4 + (i - 1)^2$.}\\
    Utilisons l'astuce $(i - 1)^2 = i^2 - 2i + 1 = -2i$ :
    \begin{itemize}
        \item $(i - 1)^2 = -2i$
        \item $(i - 1)^4 = (-2i)^2 = 4i^2 = -4$
        \item $(i - 1)^6 = (-2i)^3 = -8i^3 = 8i$
    \end{itemize}
    En additionnant : $E = 8i - 4 - 2i = \boxed{-4 + 6i}$.

    \item \textbf{Septembre 2024, Q1.4 :} \emph{Donnez le module $\rho$ et l'argument $\theta$ de $z = \dfrac{1 + i\sqrt{3}}{i + \sqrt{3}}$.}\\
    \textbf{Méthode 1 (Quotient des formes polaires) :} Au numérateur, $z_1 = 1 + i\sqrt{3}$ a pour module $\rho_1 = 2$ et argument $\theta_1 = \frac{\pi}{3}$. Au dénominateur, $z_2 = \sqrt{3} + i$ a pour module $\rho_2 = 2$ et argument $\theta_2 = \frac{\pi}{6}$. Par division :
    \[
        \rho = \frac{\rho_1}{\rho_2} = \frac{2}{2} = 1 \qquad \text{et} \qquad \theta = \theta_1 - \theta_2 = \frac{\pi}{3} - \frac{\pi}{6} = \frac{\pi}{6}.
    \]
    \textbf{Méthode 2 (Astuce de factorisation par $i$) :} Remarquons que $1 + i\sqrt{3} = i(\sqrt{3} - i)$... ou encore mieux : calculons en cartésien $\frac{(1+i\sqrt{3})(\sqrt{3}-i)}{3+1} = \frac{2\sqrt{3} + 2i}{4} = \frac{\sqrt{3}}{2} + \frac{1}{2}i = e^{i\pi/6}$.\\
    Réponse : $\boxed{\rho = 1 \quad \text{et} \quad \theta = \dfrac{\pi}{6}}$.

    \item \textbf{Juillet 2025, Q1.1 :} \emph{Exprimez $z = \dfrac{-6i}{1 + i}$ sous forme trigonométrique (module et argument).}\\
    Multiplions haut et bas par le conjugué $1 - i$ du dénominateur :
    \[
        z = \frac{-6i(1 - i)}{(1 + i)(1 - i)} = \frac{-6i + 6i^2}{1^2 + 1^2} = \frac{-6 - 6i}{2} = -3 - 3i = -3(1 + i).
    \]
    Le point $(-3, -3)$ est situé dans le \textbf{troisième quadrant} ($a < 0$ et $b < 0$). Son module et son argument valent :
    \[
        |z| = \sqrt{(-3)^2 + (-3)^2} = \sqrt{18} = 3\sqrt{2} \quad \text{et} \quad \arg(z) = -\frac{3\pi}{4} \;\left(\cos\theta = \sin\theta = -\frac{\sqrt{2}}{2}\right).
    \]
    Réponse : $\boxed{\text{module} = 3\sqrt{2}, \quad \text{argument} = -\dfrac{3\pi}{4}}$.

    \item \textbf{Juillet 2026, Q1.3 :} \emph{Calculer dans $\C$ : $z = (i - \sqrt{3})^4$ (forme rectangulaire ou polaire au choix).}\\
    Posons $w = -\sqrt{3} + i$ (deuxième quadrant : $a = -\sqrt{3} < 0, b = 1 > 0$). On a $|w| = \sqrt{3 + 1} = 2$, $\cos\theta = -\frac{\sqrt{3}}{2}$ et $\sin\theta = \frac{1}{2}$, donc $\theta = \frac{5\pi}{6}$.
    En élevant à la puissance $4$ par De Moivre :
    \[
        z = w^4 = 2^4 e^{i \cdot 4 \cdot \frac{5\pi}{6}} = 16 e^{i \frac{4\pi}{3}} = 16\left(-\frac{1}{2} - i\frac{\sqrt{3}}{2}\right) = \boxed{-8 - 8\sqrt{3}\,i \;\;\left(\text{ou } 16 e^{i\frac{4\pi}{3}}\right)}.
    \]

    \item \textbf{Septembre 2026, Q1.3 :} \emph{Donnez la forme cartésienne de $(1 + i)^{10}$.}\\
    Avec $(1 + i)^2 = 2i$, on a directement $(1 + i)^{10} = \bigl((1 + i)^2\bigr)^5 = (2i)^5 = 32 i^5 = \boxed{32i}$.
\end{enumerate}
\end{exemple}

% ----------------------------------------------------------------------------
\section{Racines $n$-ièmes d'un nombre complexe (Équations $z^n = w$)}
% ----------------------------------------------------------------------------

\begin{theoreme}[Les $n$ racines $n$-ièmes d'un nombre complexe non nul]
Soit $w = R e^{i\alpha}$ un nombre complexe non nul ($R = |w| > 0$, $\alpha = \arg(w)$) et $n \in \N_0$ un entier naturel non nul. L'équation $\boxed{z^n = w}$ possède exactement \textbf{$n$ solutions distinctes dans $\C$}, données par :
\[
    \boxed{z_k = \sqrt[n]{R} \, e^{i\left(\frac{\alpha}{n} + \frac{2k\pi}{n}\right)} = \sqrt[n]{R}\left[\cos\left(\frac{\alpha + 2k\pi}{n}\right) + i\sin\left(\frac{\alpha + 2k\pi}{n}\right)\right], \quad k \in \{0, \dots, n-1\}}.
\]
\textbf{Interprétation géométrique dans le plan complexe :} Pour $n \ge 3$, les points images des $n$ racines $z_0, z_1, \dots, z_{n-1}$ sont les sommets d'un \textbf{polygone régulier à $n$ côtés} (triangle équilatéral pour $n=3$, carré pour $n=4$, hexagone régulier pour $n=6$) inscrit dans le cercle centré à l'origine de rayon $\rho = \sqrt[n]{R}$, obtenu en tournant de $\dfrac{2\pi}{n}$ en $\dfrac{2\pi}{n}$ à partir de la première racine $z_0$.
\end{theoreme}

\begin{exemple}[Racines cubiques et quatrièmes -- UCLouvain, Septembre 2021 Q1.7, Juillet 2023 Série 1 Q1.1 \& Septembre 2025 Q1.3]
\begin{enumerate}
    \item \textbf{Septembre 2021, Q1.7 :} \emph{Représenter dans le plan complexe les solutions de $z^4 + (3\sqrt{2})^4 = 0$.}\\
    Écrivons l'équation sous la forme $z^4 = -(3\sqrt{2})^4 = (3\sqrt{2})^4 \cdot e^{i\pi}$ (car $-1 = e^{i\pi}$).
    Le module commun des $4$ racines vaut $\rho = \sqrt[4]{(3\sqrt{2})^4} = 3\sqrt{2}$, et les arguments valent $\theta_k = \frac{\pi}{4} + \frac{2k\pi}{4} = \frac{\pi}{4} + k\frac{\pi}{2}$ pour $k \in \{0, 1, 2, 3\}$ (soit $\frac{\pi}{4}, \frac{3\pi}{4}, \frac{5\pi}{4}, \frac{7\pi}{4}$).
    Comme $\cos\left(\pm\frac{\pi}{4}\right) = \frac{\sqrt{2}}{2}$, la conversion en forme cartésienne est immédiate :
    \[
        z_k = 3\sqrt{2}\left(\pm\frac{\sqrt{2}}{2} \pm i\frac{\sqrt{2}}{2}\right) = \boxed{\pm 3 \pm 3i \;\in\; \{3+3i,\; -3+3i,\; -3-3i,\; 3-3i\}},
    \]
    qui forment le carré de sommets $(\pm 3, \pm 3)$ sur la grille du plan complexe.

    \item \textbf{Juillet 2023 -- Série 1, Q1.1 :} \emph{Résoudre dans $\C$ l'équation $z^3 + 8i = 0$ sous forme cartésienne $(a + bi)$.}\\
    On isole $z^3$ et on met le second membre sous forme polaire :
    \[
        z^3 = -8i = 8 \cdot (-i) = 2^3 e^{-i\frac{\pi}{2}}.
    \]
    On applique la formule des racines cubiques ($n = 3$) avec $\rho = \sqrt[3]{2^3} = 2$ et $\theta_k = -\frac{\pi}{6} + \frac{2k\pi}{3}$ ($k \in \{0, 1, 2\}$) :
    \begin{itemize}
        \item Pour $k = 0$ ($\theta_0 = -\frac{\pi}{6}$) : $z_0 = 2\left(\cos\left(-\frac{\pi}{6}\right) + i\sin\left(-\frac{\pi}{6}\right)\right) = 2\left(\frac{\sqrt{3}}{2} - \frac{1}{2}i\right) = \sqrt{3} - i$.
        \item Pour $k = 1$ ($\theta_1 = -\frac{\pi}{6} + \frac{2\pi}{3} = \frac{\pi}{2}$) : $z_1 = 2\left(\cos\frac{\pi}{2} + i\sin\frac{\pi}{2}\right) = 2i$.
        \item Pour $k = 2$ ($\theta_2 = -\frac{\pi}{6} + \frac{4\pi}{3} = \frac{7\pi}{6} \equiv -\frac{5\pi}{6}$) : $z_2 = 2\left(-\frac{\sqrt{3}}{2} - \frac{1}{2}i\right) = -\sqrt{3} - i$.
    \end{itemize}
    \emph{Alternative algébrique :} Comme $(2i)^3 = -8i$, on factorise $z^3 - (2i)^3 = (z - 2i)(z^2 + 2iz - 4) = 0$.\\
    Réponse : $\boxed{S = \bigl\{\sqrt{3} - i,\; -\sqrt{3} - i,\; 2i\bigr\}}$.

    \item \textbf{Septembre 2025, Q1.3 :} \emph{Résoudre dans $\C$ : $z^4 + 2 - 2\sqrt{3}i = 0$ (forme rectangulaire ou polaire).}\\
    On écrit $z^4 = -2 + 2\sqrt{3}i$. Le complexe $w = -2 + 2\sqrt{3}i$ (quadrant II : $a = -2 < 0, b = 2\sqrt{3} > 0$) a pour module $R = \sqrt{4 + 12} = 4 = (\sqrt{2})^4$ et pour argument $\alpha = \frac{2\pi}{3}$ (car $\cos\alpha = -\frac{1}{2}$ et $\sin\alpha = \frac{\sqrt{3}}{2}$).
    Les $4$ racines quatrièmes ont donc pour module $\rho = \sqrt[4]{4} = \sqrt{2}$ et pour arguments $\theta_k = \frac{2\pi/3}{4} + \frac{2k\pi}{4} = \frac{\pi}{6} + k\frac{\pi}{2}$ :
    \[
        \boxed{\text{Forme polaire : } z_k = \sqrt{2}\, e^{i\left(\frac{\pi}{6} + k\frac{\pi}{2}\right)} \quad (k \in \{0, 1, 2, 3\})}.
    \]
    En forme rectangulaire, comme $z_0 = \sqrt{2}\left(\frac{\sqrt{3}}{2} + i\frac{1}{2}\right) = \frac{\sqrt{6}}{2} + i\frac{\sqrt{2}}{2}$ et que multiplier par $e^{i\pi/2} = i$ fait tourner de $90\degree$ ($z_{k+1} = i z_k$) :
    \[
        \boxed{S = \left\{\frac{\sqrt{6}}{2} + i\frac{\sqrt{2}}{2},\;\; -\frac{\sqrt{2}}{2} + i\frac{\sqrt{6}}{2},\;\; -\frac{\sqrt{6}}{2} - i\frac{\sqrt{2}}{2},\;\; \frac{\sqrt{2}}{2} - i\frac{\sqrt{6}}{2}\right\}}.
    \]
\end{enumerate}
\end{exemple}

\begin{exemple}[Équation polynomiale de degré $2n$ dans $\C$ -- UCLouvain, Juillet 2021, Q2.1]
\textbf{Énoncé :} Soit $z \in \C$ et $n \in \N_0$, résoudre $(z^2 + 1)^n = (z - i)^{2n}$.

\textbf{Résolution :} Dans $\C$, la somme de deux carrés $z^2 + 1 = z^2 - i^2$ se factorise en $(z - i)(z + i)$ ! L'équation devient :
\[
    (z - i)^n (z + i)^n - (z - i)^{2n} = 0 \iff (z - i)^n \Bigl[(z + i)^n - (z - i)^n\Bigr] = 0.
\]
Par la règle du produit nul, on distingue deux cas :
\begin{itemize}
    \item \textbf{Cas 1 :} $(z - i)^n = 0 \iff \mathbf{z = i}$ (solution valable pour tout $n \in \N_0$).
    \item \textbf{Cas 2 :} $(z + i)^n = (z - i)^n$ avec $z \neq i$. En divisant par $(z - i)^n \neq 0$, on se ramène aux \textbf{racines $n$-ièmes de l'unité} :
    \[
        \left(\frac{z + i}{z - i}\right)^n = 1 = e^{i \cdot 2k\pi} \iff \frac{z + i}{z - i} = e^{i\frac{2k\pi}{n}} \qquad \text{pour } k \in \{1, 2, \dots, n - 1\}
    \]
    (la valeur $k = 0$ donnerait $\frac{z+i}{z-i} = 1 \iff i = -i$, impossible). Isolons $z$ :
    \[
        z + i = (z - i)e^{i\frac{2k\pi}{n}} \iff z\left(1 - e^{i\frac{2k\pi}{n}}\right) = -i\left(1 + e^{i\frac{2k\pi}{n}}\right) \iff z_k = -i\,\frac{1 + e^{i\frac{2k\pi}{n}}}{1 - e^{i\frac{2k\pi}{n}}}.
    \]
    \emph{Remarque (factorisation par l'angle moitié d'Euler) :} En mettant $e^{i\frac{k\pi}{n}}$ en évidence en haut et en bas, $z_k = -i\frac{e^{-ik\pi/n} + e^{ik\pi/n}}{e^{-ik\pi/n} - e^{ik\pi/n}} = -i\frac{2\cos(k\pi/n)}{-2i\sin(k\pi/n)} = \cotg\left(\frac{k\pi}{n}\right) \in \R$ !
\end{itemize}
L'ensemble des solutions est donc :
\[
    \boxed{S = \{i\} \cup \left\{-i\,\frac{1 + e^{i\frac{2k\pi}{n}}}{1 - e^{i\frac{2k\pi}{n}}} = \cotg\left(\frac{k\pi}{n}\right) \;\middle|\; k \in \{1, \dots, n-1\}\right\}}.
\]
\end{exemple}

% ----------------------------------------------------------------------------
\section{Exercices d'entraînement du Chapitre 4}
% ----------------------------------------------------------------------------

\emph{Les corrigés détaillés de tous ces exercices se trouvent au Chapitre~10.}

\begin{enumerate}[label=\textbf{Exercice 4.\arabic* :}, leftmargin=*]
    \item \textbf{(Formes cartésienne et polaire)} Déterminer le module, l'argument principal et la forme cartésienne des nombres complexes suivants :
    \begin{enumerate}[label=(\alph*)]
        \item $z_1 = \dfrac{(1 + i)^8}{(1 - i\sqrt{3})^3}$
        \item $z_2 = \dfrac{2 - 2i}{1 + i\sqrt{3}}$
    \end{enumerate}
    \item \textbf{(Équations avec conjugué $z^*$)} Résoudre dans $\C$ (et donner les solutions sous forme cartésienne $a + bi$) :
    \begin{enumerate}[label=(\alph*)]
        \item $2z + iz^* = 5 - 4i$
        \item $z^2 + 2z^* + 1 = 0$
    \end{enumerate}
    \item \textbf{(Racines $n$-ièmes dans $\C$)} Résoudre dans $\C$ sous forme polaire et sous forme cartésienne :
    \begin{enumerate}[label=(\alph*)]
        \item $z^3 - 27 = 0$
        \item $z^4 + 16 = 0$
        \item $z^2 - (3 + 4i) = 0$ (astuce : poser $z = a + bi$ et utiliser $a^2 - b^2 = 3$, $2ab = 4$ et $a^2 + b^2 = \sqrt{3^2+4^2} = 5$).
    \end{enumerate}
\end{enumerate}

% ============================================================================
% CHAPITRE 5 : ANALYSE COMBINATOIRE, BINÔME DE NEWTON ET PROBABILITÉS
% ============================================================================
\chapter{Analyse Combinatoire, Binôme de Newton et Probabilités}

Ce chapitre couvre pas moins de \textbf{deux à trois sous-questions sur les sept} de la Question~1 à chaque session de l'examen d'Algèbre de l'UCLouvain ! Vous y trouverez systématiquement :
\begin{itemize}
    \item une question sur le \textbf{développement du binôme de Newton} $(a + b)^n$ (calcul d'un coefficient précis, du terme constant d'ordre nul en $x^0$, ou réduction de $(x+1)^n \pm (x-1)^n$) ;
    \item une ou deux questions de \textbf{dénombrement et de calcul de probabilités} (lancers de dés, tirages de cartes ou de boules dans une urne, files d'attente, sélection d'ouvrages ou de questions d'examen), ainsi que les paramètres statistiques élémentaires (moyenne et écart-type).
\end{itemize}

% ----------------------------------------------------------------------------
\section{Principes fondamentaux du dénombrement (Analyse combinatoire)}
% ----------------------------------------------------------------------------

\subsection{Principe multiplicatif (« ET ») et principe additif (« OU »)}

\begin{theoreme}[Les deux règles d'or du dénombrement]
\begin{enumerate}
    \item \textbf{Principe multiplicatif (étapes successives reliées par « ET ») :} Si une opération globale se décompose en $k$ étapes successives telles que la 1\iere{} étape offre $n_1$ choix possibles, la 2\ieme{} étape offre $n_2$ choix possibles (quel que soit le choix précédent), \dots, et la $k$-ième étape offre $n_k$ choix possibles, alors le nombre total de configurations est le \textbf{produit} :
    \[
        \boxed{N = n_1 \times n_2 \times \dots \times n_k}.
    \]
    \item \textbf{Principe additif (cas disjoints reliés par « OU exclusif ») :} Si l'ensemble des configurations se partitionne en plusieurs cas mutuellement exclusifs (disjoints) comportant respectivement $N_1, N_2, \dots, N_p$ possibilités, alors le nombre total de possibilités est la \textbf{somme} :
    \[
        \boxed{N = N_1 + N_2 + \dots + N_p}.
    \]
    \item \textbf{Principe du complémentaire (« au moins un ») :} Pour dénombrer les configurations vérifiant \emph{« au moins une fois la condition $C$ »}, il est presque toujours plus rapide de calculer le nombre total de configurations sans contrainte et d'en \textbf{soustraire} celles qui ne vérifient \emph{aucune} fois la condition $C$ :
    \[
        \boxed{N(\text{au moins un}) = N_{\text{total}} - N(\text{aucun})}.
    \]
\end{enumerate}
\end{theoreme}

\subsection{Permutations, arrangements et combinaisons : quel outil choisir ?}

Lorsqu'on choisit $k$ éléments parmi $n$ éléments distincts, posez-vous toujours \textbf{deux questions} :
\begin{enumerate}
    \item \textbf{L'ordre des éléments choisis a-t-il de l'importance ?} (Ex. un podium Or-Argent-Bronze ou un code de cadenas $\implies$ \textbf{OUI} ; une main de cartes ou un comité de 3 personnes $\implies$ \textbf{NON}).
    \item \textbf{Peut-on répéter un même élément plusieurs fois (tirage avec remise) ou non (sans remise) ?}
\end{enumerate}

\begin{theoreme}[Formulaire complet de l'analyse combinatoire ($n, k \in \N$)]
\begin{enumerate}
    \item \textbf{Factorielle $n!$ et Permutations de $n$ objets distincts :} Le nombre de façons d'ordonner $n$ objets distincts en file indienne est :
    \[
        \boxed{P_n = n! = n \times (n-1) \times (n-2) \times \dots \times 2 \times 1} \qquad \text{avec par convention } \boxed{0! = 1}.
    \]
    \emph{Valeurs à connaître par cœur :} $1! = 1$, $2! = 2$, $3! = 6$, $4! = 24$, $5! = 120$, $6! = 720$.\\
    \emph{Permutations avec répétitions :} Si parmi les $n$ objets, il y a $n_1$ objets identiques de type~1, \dots, $n_r$ objets identiques de type~$r$ ($n_1 + \dots + n_r = n$), le nombre d'anagrammes distincts vaut $\dfrac{n!}{n_1!\,n_2!\dots n_r!}$.
    \item \textbf{Arrangements avec répétition (ordre important, avec remise) :} Nombre de listes ordonnées de longueur $k$ formées à partir d'un alphabet à $n$ éléments :
    \[
        \boxed{n^k = \underbrace{n \times n \times \dots \times n}_{k \text{ facteurs}}}.
    \]
    \item \textbf{Arrangements sans répétition $A_n^k$ (ordre important, sans remise, $0 \le k \le n$) :} Nombre de podiums ou listes ordonnées de $k$ éléments \textbf{distincts} choisis parmi $n$ :
    \[
        \boxed{A_n^k = n \times (n-1) \times \dots \times (n - k + 1) = \frac{n!}{(n - k)!}}.
    \]
    \item \textbf{Combinaisons sans répétition $C_n^k = \binom{n}{k}$ (ordre SANS importance, sans remise, $0 \le k \le n$) :} Nombre de sous-ensembles (poignées simultanées) de $k$ éléments choisis parmi $n$ (obtenu en divisant $A_n^k$ par les $k!$ ordres possibles des $k$ éléments choisis) :
    \[
        \boxed{C_n^k = \binom{n}{k} = \frac{A_n^k}{k!} = \frac{n!}{k!\,(n - k)!} = \frac{n \times (n-1) \times \dots \times (n - k + 1)}{k \times (k-1) \times \dots \times 2 \times 1}}.
    \]
\end{enumerate}
\end{theoreme}

\begin{astuce}[Calcul mental rapide de $\binom{n}{k}$ sans calculatrice]
Ne calculez \textbf{jamais} $n!$ en entier ! Écrivez les $k$ facteurs décroissants à partir de $n$ au numérateur, les $k$ facteurs de $k!$ au dénominateur, et \textbf{simplifiez avant de multiplier}. De plus, si $k > \frac{n}{2}$, utilisez la symétrie $\binom{n}{k} = \binom{n}{n-k}$ :
\begin{align*}
    \binom{10}{3} &= \frac{10 \times 9 \times 8}{3 \times 2 \times 1} = 10 \times 3 \times 4 = 120, \qquad \binom{8}{6} = \binom{8}{2} = \frac{8 \times 7}{2} = 28, \\
    \binom{12}{4} &= \frac{12 \times 11 \times 10 \times 9}{4 \times 3 \times 2 \times 1} = 11 \times 5 \times 9 = 495.
\end{align*}
\end{astuce}

\subsection{Propriétés des coefficients binomiaux et triangle de Pascal}

\begin{theoreme}[Identités de symétrie et de Pascal]
Pour tous entiers $0 \le k \le n$ :
\begin{enumerate}
    \item \textbf{Valeurs aux bords :}
    \[
        \boxed{C_n^0 = \binom{n}{0} = 1, \quad C_n^n = \binom{n}{n} = 1, \quad C_n^1 = \binom{n}{1} = n, \quad C_n^2 = \binom{n}{2} = \frac{n(n-1)}{2}}.
    \]
    \item \textbf{Symétrie :}Choisir les $k$ éléments que l'on prend revient à choisir les $n - k$ éléments que l'on laisse :
    \[
        \boxed{C_n^k = C_n^{n-k} \qquad \text{soit} \qquad \binom{n}{k} = \binom{n}{n-k}}.
    \]
    \item \textbf{Relation de récurrence de Pascal :} Chaque coefficient du triangle de Pascal est la somme des deux coefficients situés juste au-dessus de lui :
    \[
        \boxed{C_n^k = C_{n-1}^{k-1} + C_{n-1}^k \qquad \Longleftrightarrow \qquad C_n^k - C_{n-1}^{k-1} = C_{n-1}^k}.
    \]
\end{enumerate}
\end{theoreme}

\begin{center}
\renewcommand{\arraystretch}{1.25}
\begin{tabular}{c|ccccccccc}
$n \setminus k$ & $0$ & $1$ & $2$ & $3$ & $4$ & $5$ & $6$ & $7$ & $8$ \\
\hline
$n = 0$ & $1$ & & & & & & & & \\
$n = 1$ & $1$ & $1$ & & & & & & & \\
$n = 2$ & $1$ & $2$ & $1$ & & & & & & \\
$n = 3$ & $1$ & $3$ & $3$ & $1$ & & & & & \\
$n = 4$ & $1$ & $4$ & $\mathbf{6}$ & $4$ & $1$ & & & & \\
$n = 5$ & $1$ & $5$ & $10$ & $10$ & $5$ & $1$ & & & \\
$n = 6$ & $1$ & $6$ & $15$ & $\mathbf{20}$ & $15$ & $6$ & $1$ & & \\
$n = 7$ & $1$ & $7$ & $21$ & $35$ & $35$ & $21$ & $7$ & $1$ & \\
$n = 8$ & $1$ & $8$ & $28$ & $56$ & $\mathbf{70}$ & $56$ & $28$ & $8$ & $1$ \\
\end{tabular}
\end{center}

\begin{exemple}[Les six problèmes de dénombrement des annales (2021--2024)]
\begin{enumerate}
    \item \textbf{Juillet 2021, Q1.3 :} \emph{Simplifier l'opération $C_n^k - C_{n-1}^{k-1}$.}\\
    Par la relation du triangle de Pascal $C_n^k = C_{n-1}^{k-1} + C_{n-1}^k$, on obtient immédiatement $\boxed{C_{n-1}^k = \binom{n-1}{k}}$.

    \item \textbf{Septembre 2021, Q1.4 :} \emph{Sur un cadenas à $5$ chiffres (de $0$ à $9$), combien peut-on former de combinaisons où le chiffre $1$ ne figure qu'une seule fois ?}\\
    Décomposons en deux étapes par le principe multiplicatif :
    \begin{itemize}
        \item \emph{Étape 1 :} Choisir laquelle des $5$ positions du cadenas est occupée par l'unique chiffre $1$ $\implies \binom{5}{1} = 5$ choix.
        \item \emph{Étape 2 :} Pour chacune des $4$ autres positions, choisir librement un chiffre parmi les $9$ chiffres restants $\{0, 2, 3, 4, 5, 6, 7, 8, 9\}$ (avec répétition autorisée) $\implies 9^4$ choix.
    \end{itemize}
    Comme $9^2 = 81$ et $9^4 = 81^2 = (80 + 1)^2 = 6400 + 160 + 1 = 6561$, le nombre total vaut :
    \[
        N = 5 \times 9^4 = 5 \times 6561 = \boxed{32\,805}.
    \]

    \item \textbf{Juillet 2022, Q1.6 :} \emph{Si un étudiant doit répondre à $4$ des $8$ questions que comporte un examen tout en choisissant (exactement) $2$ questions parmi les $5$ premières, combien a-t-il de possibilités ?}\\
    Pour former son ensemble de $4$ questions, l'étudiant doit :
    \begin{itemize}
        \item choisir $2$ questions parmi les $5$ premières : $\binom{5}{2} = \frac{5 \times 4}{2} = 10$ possibilités ;
        \item \textbf{ET} choisir les $4 - 2 = 2$ questions restantes parmi les $8 - 5 = 3$ dernières questions : $\binom{3}{2} = \binom{3}{1} = 3$ possibilités.
    \end{itemize}
    Par le principe multiplicatif : $N = \binom{5}{2} \times \binom{3}{2} = 10 \times 3 = \boxed{30}$.

    \item \textbf{Juillet 2023 -- Série 1, Q1.2 :} \emph{On compose un livret de quatre feuilles recto-verso, soit huit pages de texte. Combien y a-t-il de compositions possibles, selon l'ordre et l'orientation (recto ou verso) des pages ?}\\
    \begin{itemize}
        \item Il y a $4! = 24$ façons d'ordonner les $4$ feuilles les unes à la suite des autres.
        \item Pour chacune des $4$ feuilles, il y a $2$ orientations possibles (placer la feuille côté recto ou la retourner côté verso), soit $2^4 = 16$ choix d'orientation.
    \end{itemize}
    Par le principe multiplicatif : $N = 2^4 \cdot 4! = 16 \times 24 = \boxed{384}$.

    \item \textbf{Septembre 2023, Q1.3 :} \emph{Cinq athlètes participent à une compétition pour laquelle des médailles d'or, d'argent et de bronze sont attribuées aux trois meilleurs scores. Combien y a-t-il de façons différentes de décerner ces médailles ?}\\
    L'ordre compte (Or $\neq$ Argent $\neq$ Bronze) et un athlète ne reçoit qu'une médaille (sans remise). C'est un arrangement de $3$ éléments parmi $5$ :
    \[
        A_5^3 = 5 \times 4 \times 3 = \boxed{60}.
    \]

    \item \textbf{Septembre 2024, Q1.3 :} \emph{Une bibliothécaire veut présenter $3$ ouvrages parmi un ensemble de $8$. L'ordre n'a pas d'importance, mais elle doit obligatoirement inclure \textbf{au moins un} des deux ouvrages les plus récents. De combien de manières peut-elle faire sa sélection ?}\\
    \textbf{Méthode 1 (Complémentaire -- la plus rapide !) :}
    \begin{itemize}
        \item Nombre total de sélections de $3$ livres parmi $8$ sans contrainte : $\binom{8}{3} = \frac{8 \times 7 \times 6}{3 \times 2 \times 1} = 56$.
        \item Nombre de sélections interdites (ne contenant \textbf{aucun} des $2$ livres récents, donc choisies uniquement parmi les $8 - 2 = 6$ livres anciens) : $\binom{6}{3} = \frac{6 \times 5 \times 4}{3 \times 2 \times 1} = 20$.
    \end{itemize}
    Par soustraction : $N = \binom{8}{3} - \binom{6}{3} = 56 - 20 = \boxed{36}$.\\
    \textbf{Méthode 2 (Disjonction des cas) :} Soit exactement $1$ livre récent et $2$ anciens ($\binom{2}{1}\binom{6}{2} = 2 \times 15 = 30$), soit les $2$ livres récents et $1$ ancien ($\binom{2}{2}\binom{6}{1} = 1 \times 6 = 6$), d'où $30 + 6 = 36$.
\end{enumerate}
\end{exemple}

% ----------------------------------------------------------------------------
\section{La formule du binôme de Newton et ses applications à l'examen}
% ----------------------------------------------------------------------------

\begin{theoreme}[Formule du binôme de Newton]
Pour tous nombres réels ou complexes $a$ et $b$, et tout entier naturel $n \in \N$ :
\[
    \boxed{(a + b)^n = \sum_{k=0}^n \binom{n}{k} a^{n-k} b^k = \sum_{k=0}^n \binom{n}{k} a^k b^{n-k}}.
\]
Le \textbf{terme général} de rang $k$ ($0 \le k \le n$) dans le développement de $(a + b)^n$ est donc :
\[
    \boxed{T_k = \binom{n}{k} a^{n-k} b^k}.
\]
\end{theoreme}

\begin{astuce}[Méthode universelle pour trouver le coefficient de $x^p$ (ou le terme constant en $x^0$)]
À l'examen, lorsqu'on vous demande le coefficient de $x^p$ (ou le \emph{terme constant / terme d'ordre nul / de degré zéro}, c'est-à-dire $p = 0$) dans une expression de la forme $\left(\alpha x^r + \beta x^s\right)^n$ (où $\frac{1}{x} = x^{-1}$) :
\begin{enumerate}
    \item Écrivez le \textbf{terme général} $T_k = \binom{n}{k} \bigl(\alpha x^r\bigr)^{n-k} \bigl(\beta x^s\bigr)^k$.
    \item Séparez les constantes numériques (avec leurs signes !) et la puissance de $x$ :
    \[
        T_k = \underbrace{\binom{n}{k} \alpha^{n-k} \beta^k}_{\text{coefficient numérique}} \cdot \, x^{r(n-k) + sk}.
    \]
    \item Résolvez la petite équation du premier degré sur l'exposant : $\boxed{r(n - k) + sk = p}$ pour trouver l'entier $k \in \{0, \dots, n\}$.
    \item Remplacez cette valeur de $k$ dans le coefficient numérique $\binom{n}{k} \alpha^{n-k} \beta^k$.
\end{enumerate}
\end{astuce}

\begin{exemple}[Les dix questions sur le binôme de Newton des annales (2021--2026)]
Regroupons les trois types d'exercices posés à la Question~1 entre 2021 et 2026 :

\paragraph{Type A : Coefficient d'un monôme $x^p$ dans $(ax + b)^n$.}
\begin{enumerate}
    \item \textbf{Juillet 2021, Q1.2 :} \emph{Coefficient du terme en $a^5 b^4$ du polynôme $(a + 2b)^9$.}\\
    Le terme général est $\binom{9}{k} a^{9-k} (2b)^k = \binom{9}{k} 2^k a^{9-k} b^k$. Pour avoir $a^5 b^4$, il faut $k = 4$ :
    \[
        \text{Coeff} = \binom{9}{4} \cdot 2^4 = \frac{9 \times 8 \times 7 \times 6}{4 \times 3 \times 2 \times 1} \times 16 = 126 \times 16 = \boxed{2016}.
    \]
    \item \textbf{Juillet 2022, Q1.2 :} \emph{Coefficient du terme de degré $9$ pour $x$ du binôme $(x - 2)^{12}$.}\\
    Terme général : $\binom{12}{k} x^{12-k} (-2)^k$. Pour avoir $x^9$, on prend $12 - k = 9 \iff k = 3$ (n'oubliez pas le signe de $(-2)^3 = -8$ !) :
    \[
        \text{Coeff} = \binom{12}{3} (-2)^3 = \frac{12 \times 11 \times 10}{3 \times 2 \times 1} \times (-8) = 220 \times (-8) = \boxed{-1760}.
    \]
    \item \textbf{Septembre 2022, Q1.1 :} \emph{Coefficient du terme de degré $14$ pour $x$ de $(3x^2 + 2x^3)^6$.}\\
    Terme général : $\binom{6}{k} (3x^2)^{6-k} (2x^3)^k = \binom{6}{k} 3^{6-k} 2^k \, x^{2(6-k) + 3k} = \binom{6}{k} 3^{6-k} 2^k \, x^{12 + k}$.
    On veut l'exposant $12 + k = 14 \iff k = 2$ :
    \[
        \text{Coeff} = \binom{6}{2} \cdot 3^4 \cdot 2^2 = 15 \times 81 \times 4 = 60 \times 81 = \boxed{4860}.
    \]
    \item \textbf{Septembre 2023, Q1.2 :} \emph{Coefficient du terme de degré $7$ dans $(x - 2)^{10}$.}\\
    Pour $x^7$, il faut $10 - k = 7 \iff k = 3$, d'où $\binom{10}{3} (-2)^3 = 120 \times (-8) = \boxed{-960}$.
    \item \textbf{Juillet 2024, Q1.3 :} \emph{Coefficient du terme en $x^4$ de $(x + 2)^7 + (x - 1)^7$.}\\
    Dans $(x + a)^7$, le coefficient de $x^4$ ($7 - k = 4 \iff k = 3$) vaut $\binom{7}{3} a^3$. En additionnant pour $a = 2$ et $a = -1$ :
    \[
        \text{Coeff} = \binom{7}{3} 2^3 + \binom{7}{3} (-1)^3 = \binom{7}{3}\bigl(8 - 1\bigr) = 35 \times 7 = \boxed{245}.
    \]
    \item \textbf{Septembre 2026, Q1.5 :} \emph{Coefficient du terme en $x^2$ dans $(x + 2)^6$.}\\
    Pour $x^2$, le terme $\binom{6}{2} x^2 \cdot 2^{6-2}$ donne $\binom{6}{2} \cdot 2^4 = 15 \times 16 = \boxed{240}$.
\end{enumerate}

\paragraph{Type B : Terme constant (d'ordre nul / de degré zéro en $x^0$) dans $\left(ax^p + \frac{b}{x}\right)^n$.}
\begin{enumerate}[resume]
    \item \textbf{Juillet 2023 -- Série 2, Q1.3 :} \emph{Terme constant du développement de $\left(x + \dfrac{1}{x}\right)^8$.}\\
    Terme général : $\binom{8}{k} x^{8-k} (x^{-1})^k = \binom{8}{k} x^{8-2k}$. Terme constant ($x^0$) : $8 - 2k = 0 \iff k = 4$, d'où $\binom{8}{4} = \boxed{70}$.
    \item \textbf{Septembre 2024, Q1.5 :} \emph{Terme de degré zéro du développement de $\left(x + \dfrac{2}{x}\right)^6$.}\\
    Terme général : $\binom{6}{k} x^{6-k} (2x^{-1})^k = \binom{6}{k} 2^k x^{6-2k}$. Pour $6 - 2k = 0 \iff k = 3$ :
    \[
        T_3 = \binom{6}{3} \cdot 2^3 = 20 \times 8 = \boxed{160}.
    \]
    \item \textbf{Septembre 2025, Q1.5 :} \emph{Terme d'ordre nul résultant du développement de $\left(x^2 + \dfrac{1}{x}\right)^6$.}\\
    Terme général : $\binom{6}{k} (x^2)^{6-k} (x^{-1})^k = \binom{6}{k} x^{12-3k}$. Pour $12 - 3k = 0 \iff k = 4$ (ou $k = 2$ en inversant les rôles) :
    \[
        \binom{6}{4} = \binom{6}{2} = \boxed{15}.
    \]
    \item \textbf{Juillet 2026, Q1.5 :} \emph{Coefficient du terme d'ordre nul dans $\left(2x + \dfrac{1}{x}\right)^8$.}\\
    Terme général : $\binom{8}{k} (2x)^{8-k} (x^{-1})^k = \binom{8}{k} 2^{8-k} x^{8-2k}$. Pour $8 - 2k = 0 \iff k = 4$ :
    \[
        \binom{8}{4} \cdot 2^4 = 70 \times 16 = \boxed{1120}.
    \]
\end{enumerate}

\paragraph{Type C : Simplification de $(x + 1)^n + (x - 1)^n$ et $(x + 1)^n - (x - 1)^n$.}
\begin{itemize}
    \item Lorsqu'on \textbf{additionne} $(x + 1)^n + (x - 1)^n$, tous les termes où $(-1)^k$ est impair s'annulent et les termes où $k$ est pair doublent !
    \item Lorsqu'on \textbf{soustrait} $(x + 1)^n - (x - 1)^n$, tous les termes où $k$ est pair s'annulent et les termes où $k$ est impair doublent !
\end{itemize}
\begin{enumerate}[resume]
    \item \textbf{Juillet 2023 -- Série 1, Q1.3 :} \emph{Forme polynomiale de $(x + 1)^5 + (x - 1)^5$.}\\
    La ligne $n = 5$ du triangle de Pascal est $1, 5, 10, 10, 5, 1$. En additionnant $(x+1)^5 + (x-1)^5$, seuls les termes de degrés impairs ($x^5, x^3, x^1$) survivent et sont multipliés par $2$ :
    \[
        (x + 1)^5 + (x - 1)^5 = 2\left(\binom{5}{0}x^5 + \binom{5}{2}x^3 + \binom{5}{4}x\right) = 2(x^5 + 10x^3 + 5x) = \boxed{2x^5 + 20x^3 + 10x}.
    \]
    \item \textbf{Juillet 2025, Q1.5 :} \emph{Réduisez $P(x) = (x + 1)^6 - (x - 1)^6 - 12x(x^2 + 1)^2$ à sa plus simple expression.}\\
    La ligne $n = 6$ de Pascal est $1, 6, 15, 20, 15, 6, 1$. Dans la différence $(x + 1)^6 - (x - 1)^6$, seuls les termes de puissances impaires ($k = 1, 3, 5$, soit $x^5, x^3, x^1$) survivent, multipliés par $2$ :
    \[
        (x + 1)^6 - (x - 1)^6 = 2\bigl(6x^5 + 20x^3 + 6x\bigr) = 12x^5 + 40x^3 + 12x.
    \]
    Or, d'autre part, $12x(x^2 + 1)^2 = 12x(x^4 + 2x^2 + 1) = 12x^5 + 24x^3 + 12x$. En soustrayant les deux expressions :
    \[
        P(x) = (12x^5 + 40x^3 + 12x) - (12x^5 + 24x^3 + 12x) = \boxed{16x^3}.
    \]
\end{enumerate}
\end{exemple}

% ----------------------------------------------------------------------------
\section{Calcul des probabilités et paramètres statistiques}
% ----------------------------------------------------------------------------

\subsection{Loi de Laplace (équiprobabilité) et règles de probabilité}

\begin{theoreme}[Formule de Laplace et probabilités fondamentales]
Lorsqu'une expérience aléatoire possède un univers fini $\Omega$ d'issues \textbf{équiprobables} (dés équilibrés, cartes bien battues, boules indiscernables au toucher), la probabilité d'un événement $A \subset \Omega$ est le rapport :
\[
    \boxed{P(A) = \frac{\text{Nombre de cas favorables à } A}{\text{Nombre total de cas possibles}} = \frac{\text{card}(A)}{\text{card}(\Omega)}}.
\]
\textbf{Attention à la cohérence du modèle :} Si vous comptez les cas possibles au dénominateur \textbf{avec ordre} (par exemple $6^k$ pour $k$ dés numérotés de $1$ à $k$), vous devez impérativement compter les cas favorables au numérateur \textbf{avec ordre également} ! Si vous comptez au dénominateur \textbf{sans ordre} (par exemple $\binom{N}{k}$ pour un tirage simultané sans remise), comptez aussi le numérateur \textbf{sans ordre}.
\begin{itemize}
    \item \textbf{Événement contraire $\bar{A}$ :} $\boxed{P(\bar{A}) = 1 - P(A)}$ (idéal pour « au moins un »).
    \item \textbf{Union de deux événements :} $\boxed{P(A \cup B) = P(A) + P(B) - P(A \cap B)}$.
    \item \textbf{Événements indépendants :} $A$ et $B$ sont indépendants ssi $\boxed{P(A \cap B) = P(A)\cdot P(B)}$.
\end{itemize}
\end{theoreme}

\begin{exemple}[Les neuf problèmes de probabilités des annales (2021--2026)]
Étudions méthodiquement toutes les questions de probabilités posées à l'examen d'admission de 2021 à 2026 :

\paragraph{1. Lancers de pièces et de dés (toujours distinguer les dés au dénominateur : $6^n$ issues équiprobables !).}
\begin{enumerate}
    \item \textbf{Juillet 2021, Q1.4 :} \emph{Probabilité que trois pièces de monnaie équilibrées lancées simultanément tombent du même côté.}\\
    Il y a $2^3 = 8$ issues équiprobables au total ($\text{PPP}, \text{PPF}, \dots, \text{FFF}$). Seules $2$ issues sont favorables ($\text{PPP}$ et $\text{FFF}$), d'où :
    \[
        P = \frac{2}{2^3} = \frac{2}{8} = \boxed{\frac{1}{4}}.
    \]

    \item \textbf{Septembre 2023, Q1.5 :} \emph{On lance deux pièces simultanément. Quelle est la probabilité d'obtenir un « double face » ($\text{FF}$) en un maximum de quatre lancers ?}\\
    Lors d'un lancer de deux pièces, il y a $4$ issues équiprobables ($\text{PP}, \text{PF}, \text{FP}, \text{FF}$), donc $P(\text{FF}) = \frac{1}{4}$ et la probabilité de \textbf{ne pas} obtenir un double face sur un lancer vaut $P(\overline{\text{FF}}) = 1 - \frac{1}{4} = \frac{3}{4}$.
    Par passage à l'\textbf{événement contraire} : ne pas obtenir un seul double face au cours des $4$ lancers indépendants a pour probabilité $\left(\frac{3}{4}\right)^4 = \frac{81}{256}$. La probabilité d'en obtenir \textbf{au moins un} en $4$ lancers vaut donc :
    \[
        P = 1 - \left(\frac{3}{4}\right)^4 = 1 - \frac{81}{256} = \boxed{\frac{175}{256}}.
    \]

    \item \textbf{Septembre 2022, Q1.2 :} \emph{On lance deux fois d'affilée deux dés non distinguables. Quelle est la probabilité que les deux tirages soient identiques (l'ordre des deux dés n'ayant pas d'importance) ?}\\
    \emph{Attention au piège d'équiprobabilité :} même si l'on ne regarde pas l'ordre des deux dés d'un lancer, un « double » comme $\{3, 3\}$ n'a qu'$1$ chance sur $36$ de sortir, alors qu'une paire de chiffres différents comme $\{2, 5\}$ correspond à $2$ issues ordonnées $(2,5)$ et $(5,2)$, soit $\frac{2}{36}$ !
    Au premier lancer de deux dés, deux cas disjoints se présentent :
    \begin{itemize}
        \item \emph{Cas 1 : Le 1\ier{} lancer est un double} (probabilité $\frac{6}{36} = \frac{1}{6}$). Pour que le 2\ieme{} lancer soit identique, il doit reproduire exactement ce même double (probabilité $\frac{1}{36}$). Contribution : $\frac{6}{36} \times \frac{1}{36} = \frac{6}{1296}$.
        \item \emph{Cas 2 : Le 1\ier{} lancer donne deux chiffres distincts} (probabilité $\frac{30}{36} = \frac{5}{6}$). Pour que le 2\ieme{} lancer donne les deux mêmes chiffres (dans n'importe quel ordre), il y a $2$ issues favorables sur $36$ (probabilité $\frac{2}{36}$). Contribution : $\frac{30}{36} \times \frac{2}{36} = \frac{60}{1296}$.
    \end{itemize}
    En additionnant : $P = \frac{6 + 60}{1296} = \frac{66}{1296} = \boxed{\dfrac{11}{216}}$.

    \item \textbf{Juillet 2024, Q1.1 :} \emph{On lance $5$ dés à six faces équilibrés. Quelle est la probabilité d'obtenir un \textbf{full} (c'est-à-dire $3$ dés identiques d'une valeur et $2$ autres dés identiques d'une autre valeur, ex. 3-3-3-5-5) ?}\\
    Au dénominateur, en distinguant les $5$ dés, il y a $6^5 = 7776$ issues équiprobables.
    Au numérateur, construisons un full étape par étape (principe multiplicatif) :
    \begin{itemize}
        \item Choisir la valeur faciale (parmi $6$) du brelan ($3$ dés identiques) : $6$ choix.
        \item Choisir la valeur faciale (parmi les $5$ restantes) de la paire ($2$ dés identiques) : $5$ choix.
        \item Choisir lesquels des $5$ dés forment la paire (les $3$ autres formant automatiquement le brelan) : $\binom{5}{2} = \binom{5}{3} = 10$ choix.
    \end{itemize}
    D'où :
    \[
        P(\text{full}) = \frac{6 \times 5 \times \binom{5}{2}}{6^5} = \frac{30 \times 10}{6^5} = \frac{5 \times 10}{6^4} = \frac{50}{1296} = \boxed{\frac{25}{648}} \quad \left(\text{ou } \frac{300}{7776}\right).
    \]

    \item \textbf{Juillet 2026, Q1.4 :} \emph{On lance $3$ dés équilibrés à six faces. Quelle est la probabilité d'obtenir \textbf{exactement une paire} de dés identiques ?}\\
    Nombre total d'issues : $6^3 = 216$. Pour former exactement une paire (2 dés d'une valeur et le 3\ieme{} dé d'une valeur différente) :
    \begin{itemize}
        \item Choisir lesquels des $3$ dés portent la paire : $\binom{3}{2} = 3$ choix.
        \item Choisir la valeur faciale de la paire : $6$ choix.
        \item Choisir la valeur faciale du 3\ieme{} dé (différente de la paire) : $5$ choix.
    \end{itemize}
    D'où : $P = \dfrac{3 \times 6 \times 5}{6^3} = \dfrac{15}{36} = \boxed{\dfrac{5}{12}}$.
\end{enumerate}

\paragraph{2. Tirages simultanés sans remise (urnes, cartes, chaussettes) et files d'attente.}
\begin{enumerate}[resume]
    \item \textbf{Juillet 2022, Q1.3 :} \emph{Une file d'attente est formée au hasard avec $n$ personnes ($n \ge 2$). Quelle est la probabilité que Jean et Alice soient situés directement l'un à la suite de l'autre ?}\\
    Il y a $n!$ façons d'ordonner les $n$ personnes. Si l'on « soude » Jean et Alice en un seul bloc $[JA]$, il reste $(n - 1)$ objets à permuter, soit $(n - 1)!$ ordres, à multiplier par $2$ car le bloc peut être $[JA]$ ou $[AJ]$. D'où :
    \[
        P = \frac{2 \cdot (n - 1)!}{n!} = \boxed{\frac{2}{n}}.
    \]

    \item \textbf{Juillet 2023 -- Série 2, Q1.2 :} \emph{Une manne contient $4$ paires de chaussettes mélangées (soit $8$ chaussettes). On en tire $4$ au hasard. Quelle est la probabilité de tirer quatre chaussettes de paires toutes différentes ?}\\
    \begin{itemize}
        \item Nombre total de poignées de $4$ chaussettes parmi $8$ : $\binom{8}{4} = \frac{8 \times 7 \times 6 \times 5}{4 \times 3 \times 2 \times 1} = 70$.
        \item Pour avoir $4$ chaussettes issues des $4$ paires différentes, on doit prendre exactement $1$ chaussette dans la paire~1 ($2$ choix : gauche ou droite), $1$ dans la paire~2 ($2$ choix), $1$ dans la paire~3 ($2$ choix) et $1$ dans la paire~4 ($2$ choix), soit $2^4 = 16$ poignées favorables.
    \end{itemize}
    D'où : $P = \dfrac{2^4}{\binom{8}{4}} = \dfrac{16}{70} = \boxed{\dfrac{8}{35}}$.

    \item \textbf{Juillet 2025, Q1.4 :} \emph{Une urne contient $6$ boules rouges et $4$ boules vertes (soit $10$ boules). On tire $3$ boules simultanément sans remise. Quelle est la probabilité d'obtenir exactement $2$ rouges et $1$ verte ?}\\
    On choisit $2$ rouges parmi $6$ \textbf{et} $1$ verte parmi $4$, sur un total de $\binom{10}{3}$ tirages possibles :
    \[
        P = \frac{\binom{6}{2} \cdot \binom{4}{1}}{\binom{10}{3}} = \frac{15 \times 4}{120} = \frac{60}{120} = \boxed{\frac{1}{2}}.
    \]

    \item \textbf{Septembre 2025, Q1.4 :} \emph{On tire sans remise $4$ cartes d'un jeu ne contenant que les figures (rois, dames, valets) des quatre couleurs (soit $3 \times 4 = 12$ cartes). Quelle est la probabilité de tirer la même figure dans les quatre couleurs (un carré de rois, de dames ou de valets) ?}\\
    Il y a $\binom{12}{4} = \frac{12 \times 11 \times 10 \times 9}{24} = 495$ mains de $4$ cartes possibles, et seulement $3$ mains favorables (les 4 rois, les 4 dames ou les 4 valets) :
    \[
        P = \frac{3}{\binom{12}{4}} = \frac{3}{495} = \boxed{\frac{1}{165}}.
    \]

    \item \textbf{Septembre 2026, Q1.4 :} \emph{On tire simultanément $2$ cartes d'un jeu complet de $52$ cartes. Quelle est la probabilité d'obtenir deux as ?}\\
    Il y a $4$ as dans le jeu de $52$ cartes, donc $\binom{4}{2} = 6$ paires d'as favorables sur $\binom{52}{2} = \frac{52 \times 51}{2} = 26 \times 51 = 1326$ paires de cartes possibles :
    \[
        P = \frac{\binom{4}{2}}{\binom{52}{2}} = \frac{6}{26 \times 51} = \frac{1}{13 \times 17} = \boxed{\frac{1}{221}}.
    \]
\end{enumerate}
\end{exemple}

\subsection{Paramètres statistiques : moyenne et écart-type d'une distribution}

\begin{theoreme}[Moyenne $\bar{x}$, variance $\sigma^2$ et écart-type $\sigma$]
Soit une distribution statistique où les valeurs $x_1, x_2, \dots, x_k$ apparaissent avec les effectifs respectifs $n_1, n_2, \dots, n_k$ (effectif total $N = \sum_{i=1}^k n_i$) :
\begin{enumerate}
    \item \textbf{Moyenne arithmétique pondérée :}
    \[
        \boxed{\bar{x} = \frac{1}{N}\sum_{i=1}^k n_i x_i = \frac{n_1 x_1 + n_2 x_2 + \dots + n_k x_k}{n_1 + n_2 + \dots + n_k}}.
    \]
    \item \textbf{Variance $\sigma^2$ et écart-type $\sigma$ (formule de définition et formule de König-Huygens) :}
    \[
        \boxed{\sigma^2 = \frac{1}{N}\sum_{i=1}^k n_i (x_i - \bar{x})^2 = \left(\frac{1}{N}\sum_{i=1}^k n_i x_i^2\right) - \bar{x}^2} \qquad \text{et} \qquad \boxed{\sigma = \sqrt{\sigma^2}}.
    \]
\end{enumerate}
\end{theoreme}

\begin{exemple}[Calcul de moyenne et d'écart-type -- UCLouvain, Septembre 2021, Q1.2]
\textbf{Énoncé :} Une étude sur Louvain-la-Neuve a relevé la distribution suivante de la superficie des chambres d'étudiants :
\begin{center}
\begin{tabular}{|l|c|c|c|c|c|}
\hline
Superficie $x_i$ [$\text{m}^2$] & $6$ & $9$ & $12$ & $15$ & $18$ \\ \hline
Nombre de chambres $n_i$ [-] & $500$ & $1400$ & $1200$ & $800$ & $300$ \\ \hline
\end{tabular}
\end{center}
Déterminer la moyenne et l'écart-type de la superficie des chambres.

\textbf{Résolution :} Simplifions d'abord tous les effectifs par $100$ (les fréquences $f_i = \frac{n_i}{N}$ ne changent pas !) : les poids réduits sont $5, 14, 12, 8, 3$, de somme $N' = 5 + 14 + 12 + 8 + 3 = 42$.
\begin{itemize}
    \item \textbf{Moyenne :}
    \[
        \bar{x} = \frac{5(6) + 14(9) + 12(12) + 8(15) + 3(18)}{42} = \frac{30 + 126 + 144 + 120 + 54}{42} = \frac{474}{42} = \frac{79}{7} \approx \boxed{11,29\text{ m}^2}.
    \]
    \item \textbf{Variance et écart-type :} Par König-Huygens, $\frac{1}{42}\sum n'_i x_i^2 = \frac{5(36) + 14(81) + 12(144) + 8(225) + 3(324)}{42} = \frac{180 + 1134 + 1728 + 1800 + 972}{42} = \frac{5814}{42} = \frac{969}{7} \approx 138,43$. D'où $\sigma^2 = \frac{969}{7} - \left(\frac{79}{7}\right)^2 = \frac{6783 - 6241}{49} = \frac{542}{49} \approx 11,06$ et $\sigma = \sqrt{\frac{542}{49}} = \frac{\sqrt{542}}{7} \approx \boxed{3,33\text{ m}^2}$.
\end{itemize}
\end{exemple}

% ----------------------------------------------------------------------------
\section{Exercices d'entraînement du Chapitre 5}
% ----------------------------------------------------------------------------

\emph{Les corrigés détaillés de tous ces exercices se trouvent au Chapitre~10.}

\begin{enumerate}[label=\textbf{Exercice 5.\arabic* :}, leftmargin=*]
    \item \textbf{(UCLouvain -- Septembre 2021, Q1.5)} Soit un jeu de $52$ cartes que l'on distribue en quatre paquets de $13$ cartes. Quelle est la probabilité que l'un (au moins) des quatre paquets comporte $13$ des $26$ cartes rouges que compte le jeu ?
    \item \textbf{(Binôme de Newton)}
    \begin{enumerate}[label=(\alph*)]
        \item Déterminer le coefficient du terme en $x^5$ dans le développement de $(2x - 1)^8$.
        \item Déterminer le terme constant (d'ordre nul en $x^0$) dans le développement de $\left(3x^2 - \dfrac{1}{3x}\right)^9$.
    \end{enumerate}
    \item \textbf{(Probabilités aux dés et aux cartes)}
    \begin{enumerate}[label=(\alph*)]
        \item On lance $4$ dés équilibrés à $6$ faces. Quelle est la probabilité d'obtenir un \textbf{carré} ($4$ dés identiques) ? Et celle d'obtenir \textbf{quatre chiffres tous différents} ?
        \item On tire simultanément $3$ boules d'une urne contenant $5$ boules blanches, $3$ boules noires et $2$ boules rouges. Quelle est la probabilité d'obtenir $3$ boules de la même couleur ? Et celle d'obtenir $1$ boule de chaque couleur ?
    \end{enumerate}
\end{enumerate}

% ============================================================================
% CHAPITRE 6 : SUITES ARITHMÉTIQUES ET GÉOMÉTRIQUES
% ============================================================================
\chapter{Suites Arithmétiques et Géométriques}

À chaque session de l'examen d'admission d'Algèbre de l'UCLouvain, la Question~1 se termine (en Q1.5, Q1.6 ou Q1.7) par un problème concret ou direct sur les \textbf{suites arithmétiques} ou les \textbf{suites géométriques}. De plus, la Question~3 (le grand problème textuel) mobilise très souvent une suite géométrique (dégradation annuelle de panneaux photovoltaïques en Juillet~2022, perte de capacité d'une batterie de robot-tondeuse en Juillet~2021) ou une suite arithmétique (barres de suspension d'un escalier en Septembre~2023).

% ----------------------------------------------------------------------------
\section{Suites arithmétiques (progression linéaire à pas constant $+r$)}
% ----------------------------------------------------------------------------

\subsection{Définition, terme général et moyenne arithmétique}

\begin{definition}[Suite arithmétique de raison $r$]
Une suite de nombres réels $(u_n)_{n \ge 1}$ est dite \textbf{arithmétique} si l'on passe d'un terme au suivant en \textbf{ajoutant toujours la même constante} $r \in \R$, appelée la \textbf{raison} de la suite :
\[
    \boxed{\forall n \ge 1 : u_{n+1} = u_n + r \quad \Longleftrightarrow \quad u_{n+1} - u_n = r = \text{constante}}.
\]
\begin{itemize}
    \item Si $r > 0$, la suite est strictement \textbf{croissante}.
    \item Si $r < 0$, la suite est strictement \textbf{décroissante}.
\end{itemize}
\end{definition}

\begin{theoreme}[Terme général et moyenne arithmétique]
Soit $(u_n)$ une suite arithmétique de raison $r$ :
\begin{enumerate}
    \item \textbf{Terme général à partir de $u_1$ (ou à partir d'un terme quelconque $u_p$) :}
    \[
        \boxed{u_n = u_1 + (n - 1)r} \qquad \text{et plus généralement} \qquad \boxed{u_n = u_p + (n - p)r}.
    \]
    \emph{Remarque :} Si la suite est indexée à partir de $u_0$, alors $u_n = u_0 + n r$.
    \item \textbf{Propriété caractéristique (Moyenne arithmétique) :} Trois nombres consécutifs $a, b, c$ sont en progression arithmétique si et seulement si le terme du milieu est la \textbf{moyenne arithmétique} des deux termes qui l'encadrent :
    \[
        \boxed{b = \frac{a + c}{2} \quad \Longleftrightarrow \quad a + c = 2b}.
    \]
\end{enumerate}
\end{theoreme}

\subsection{Somme des $n$ premiers termes d'une suite arithmétique (Formule de Gauss)}

\begin{theoreme}[Somme de $n$ termes consécutifs d'une suite arithmétique]
La somme $S_n = u_1 + u_2 + \dots + u_n = \sum_{k=1}^n u_k$ des $n$ premiers termes d'une suite arithmétique est égale au \textbf{nombre de termes} multiplié par la \textbf{moyenne du premier et du dernier terme} :
\[
    \boxed{S_n = n \cdot \frac{u_1 + u_n}{2} = \frac{n}{2}\bigl(2u_1 + (n - 1)r\bigr)}.
\]
\textbf{Cas particuliers indispensables pour le calcul mental :}
\begin{itemize}
    \item \textbf{Somme des $n$ premiers entiers naturels :}
    \[
        \boxed{1 + 2 + 3 + \dots + n = \frac{n(n + 1)}{2}}.
    \]
    \item \textbf{Somme des $n$ premiers entiers pairs :}
    \[
        \boxed{2 + 4 + 6 + \dots + 2n = 2(1 + 2 + \dots + n) = n(n + 1)}.
    \]
    \item \textbf{Somme des $n$ premiers entiers impairs :}
    \[
        \boxed{1 + 3 + 5 + \dots + (2n - 1) = n^2}.
    \]
\end{itemize}
\end{theoreme}

\begin{center}
\begin{tikzpicture}[scale=0.65, >=Stealth]
    % Illustration visuelle de la formule de Gauss : escalier retourné
    \foreach \i in {1,...,5} {
        \pgfmathsetmacro\h{0.6*\i + 0.4}
        \fill[uclblue!75, draw=white, thick] (\i-1, 0) rectangle (\i, \h);
        \pgfmathsetmacro\htop{0.6*6 + 0.8}
        \fill[uclaccent!35, draw=white, thick] (\i-1, \h) rectangle (\i, \htop);
    }
    \draw[thick, darkgray] (0,0) rectangle (5, 4.4);
    \draw[<->, thick, warnred] (-0.4, 0) -- (-0.4, 1.0) node[midway, left] {\small $u_1$};
    \draw[<->, thick, warnred] (-0.4, 1.0) -- (-0.4, 4.4) node[midway, left] {\small $u_n$};
    \draw[<->, thick, uclblue] (5.4, 0) -- (5.4, 4.4) node[midway, right] {\small $u_1 + u_n$};
    \draw[<->, thick, darkgray] (0, -0.5) -- (5, -0.5) node[midway, below] {\small $n$ colonnes : $2 S_n = n(u_1 + u_n)$};
\end{tikzpicture}
\end{center}

\begin{exemple}[Les cinq problèmes de suites arithmétiques des annales (2022--2026)]
\begin{enumerate}
    \item \textbf{Juillet 2022, Q1.7 :} \emph{La somme de $12$ termes d'une suite arithmétique de raison égale à $3$ vaut $2022$. Quel est le premier terme de cette suite ?}\\
    On applique la formule de la somme avec $n = 12$, $r = 3$ et $S_{12} = 2022$ :
    \[
        S_{12} = \frac{12}{2}\bigl(2u_1 + (12 - 1)\cdot 3\bigr) = 6\bigl(2u_1 + 33\bigr) = 2022.
    \]
    Divisons par $6$ : $2u_1 + 33 = \dfrac{2022}{6} = 337 \iff 2u_1 = 304 \iff \boxed{u_1 = 152}$.

    \item \textbf{Juillet 2023 -- Série 2, Q1.6 (L'échelle du menuisier) :} \emph{Un menuisier construit une échelle en bois avec $9$ échelons dont la taille décroît linéairement depuis le bas de l'échelle jusqu'à son sommet. Quelle est la largeur du dernier échelon si le premier (en bas) mesure $50\text{ cm}$ et que l'ensemble des échelons est coupé sans déchet à partir d'une latte longue de $3,6\text{ m}$ ?}\\
    Les longueurs $u_1, u_2, \dots, u_9$ des $9$ échelons forment une suite arithmétique avec $n = 9$, $u_1 = 50\text{ cm}$ et une somme totale $S_9 = 3,6\text{ m} = 360\text{ cm}$.
    Utilisons directement la formule $S_9 = 9 \cdot \dfrac{u_1 + u_9}{2}$ sans même avoir besoin de calculer la raison $r$ :
    \[
        9 \cdot \frac{50 + u_9}{2} = 360 \iff 50 + u_9 = \frac{360 \times 2}{9} = 40 \times 2 = 80 \iff \boxed{u_9 = 30\text{ cm}}.
    \]

    \item \textbf{Septembre 2024, Q1.7 :} \emph{Quel est le premier terme d'une suite arithmétique de raison $3$ dont la somme des $10$ premiers termes vaut $185$ ?}\\
    Avec $n = 10$, $r = 3$ et $S_{10} = 185$ :
    \[
        S_{10} = \frac{10}{2}\bigl(2u_1 + 9\cdot 3\bigr) = 5(2u_1 + 27) = 185 \iff 2u_1 + 27 = 37 \iff \boxed{u_1 = 5}.
    \]

    \item \textbf{Juillet 2025, Q1.6 (La soirée étudiante chantante) :} \emph{Lors d'une soirée étudiante, les participants chantent à intervalles réguliers. Après chaque période de trois quarts d'heure ($45\text{ min}$), deux étudiants quittent la fête de sorte que les deux derniers clôturent la fête après $9\text{ heures}$. À chaque période, chaque étudiant encore présent entame à son tour une chanson. Combien de chansons ont été entonnées durant la soirée et combien d'étudiants ont pris part à la fête ?}\\
    \begin{itemize}
        \item \emph{Nombre de périodes :} Une durée de $9\text{ h}$ divisée par des périodes de $\frac{3}{4}\text{ h}$ ($45\text{ min}$) donne $n = \frac{9}{3/4} = 9 \times \frac{4}{3} = 12$ périodes.
        \item \emph{Nombre initial d'étudiants :} À la $12$\ieme{} et dernière période, il reste $u_{12} = 2$ étudiants. Comme $2$ étudiants sont partis après chacune des $11$ périodes précédentes ($r = -2$), le nombre initial d'étudiants à la 1\iere{} période était :
        \[
            u_1 = u_{12} + 11 \times 2 = 2 + 22 = \boxed{24\text{ étudiants}}.
        \]
        \item \emph{Nombre total de chansons :} C'est la somme des $12$ termes de cette suite arithmétique :
        \[
            S_{12} = 12 \cdot \frac{u_1 + u_{12}}{2} = 6 \times (24 + 2) = 6 \times 26 = \boxed{156\text{ chansons}}.
        \]
    \end{itemize}

    \item \textbf{Juillet 2026, Q1.7 (Les tronçons de corde) :} \emph{On a $240\text{ m}$ de corde. On coupe des tronçons de tailles croissantes successives de $2, 4, 6, \dots\text{ m}$. Combien de tronçons entiers pourra-t-on couper ?}\\
    La longueur totale de $n$ tronçons est la somme des $n$ premiers entiers pairs :
    \[
        S_n = 2 + 4 + 6 + \dots + 2n = n(n + 1).
    \]
    On cherche le plus grand entier $n$ tel que $n(n + 1) \le 240$. Or $15 \times 16 = 240$ tombe pile ! On peut donc couper exactement $\boxed{15\text{ tronçons}}$.
\end{enumerate}
\end{exemple}

% ----------------------------------------------------------------------------
\section{Suites géométriques (progression exponentielle à facteur constant $\times q$)}
% ----------------------------------------------------------------------------

\subsection{Définition, terme général et moyenne géométrique}

\begin{definition}[Suite géométrique de raison $q$]
Une suite de nombres réels $(u_n)_{n \ge 1}$ est dite \textbf{géométrique} si l'on passe d'un terme au suivant en \textbf{multipliant toujours par la même constante non nulle} $q \in \R_0$, appelée la \textbf{raison} de la suite :
\[
    \boxed{\forall n \ge 1 : u_{n+1} = q \cdot u_n \quad \Longleftrightarrow \quad \frac{u_{n+1}}{u_n} = q = \text{constante}}.
\]
\textbf{Attention aux évolutions en pourcentage :}
\begin{itemize}
    \item Une grandeur qui \textbf{augmente de $p\,\%$} chaque année est une suite géométrique de raison $\boxed{q = 1 + \frac{p}{100}}$.
    \item Une grandeur qui \textbf{diminue (se dégrade) de $p\,\%$} chaque année est une suite géométrique de raison $\boxed{q = 1 - \frac{p}{100}}$ (par exemple, une baisse de $5\,\%$/an correspond à $q = 0,95$).
\end{itemize}
\end{definition}

\begin{theoreme}[Terme général et moyenne géométrique]
Soit $(u_n)$ une suite géométrique de raison $q \neq 0$ :
\begin{enumerate}
    \item \textbf{Terme général à partir de $u_1$ (ou à partir d'un terme quelconque $u_p$) :}
    \[
        \boxed{u_n = u_1 \cdot q^{n-1}} \qquad \text{et plus généralement} \qquad \boxed{u_n = u_p \cdot q^{n-p}}.
    \]
    \item \textbf{Propriété caractéristique (Moyenne géométrique) :} Trois nombres strictement positifs $a, b, c$ sont en progression géométrique si et seulement si :
    \[
        \boxed{b^2 = a \cdot c \quad \Longleftrightarrow \quad b = \sqrt{a \cdot c}}.
    \]
\end{enumerate}
\end{theoreme}

\subsection{Somme des $n$ premiers termes d'une suite géométrique}

\begin{theoreme}[Somme de $n$ termes consécutifs d'une suite géométrique ($q \neq 1$)]
Pour toute suite géométrique de premier terme $u_1$ et de raison $q \neq 1$, la somme des $n$ premiers termes $S_n = u_1 + u_2 + \dots + u_n = u_1(1 + q + q^2 + \dots + q^{n-1})$ vaut :
\[
    \boxed{S_n = u_1 \cdot \frac{1 - q^n}{1 - q} = u_1 \cdot \frac{q^n - 1}{q - 1}}.
\]
\textbf{Astuce de signe pour le calcul mental :}
\begin{itemize}
    \item Si $\mathbf{q > 1}$ (ex. $q = 2$ ou $q = 3/2$), utilisez la seconde forme $\boxed{S_n = u_1 \cdot \frac{q^n - 1}{q - 1}}$ pour que le numérateur et le dénominateur soient tous deux \textbf{positifs} ! En particulier, pour $q = 2$ : $\boxed{1 + 2 + 4 + \dots + 2^{n-1} = 2^n - 1}$.
    \item Si $\mathbf{0 < q < 1}$ (ex. $q = 3/5$ ou $q = 0,95$), utilisez la première forme $\boxed{S_n = u_1 \cdot \frac{1 - q^n}{1 - q}}$.
\end{itemize}
\end{theoreme}

\begin{astuce}[Valeurs des puissances de $2$ et de $3$ à connaître par cœur le jour de l'examen]
Sans calculatrice, reconnaître instantanément les puissances de $2$ et de $3$ fait gagner un temps précieux :
\begin{center}
\renewcommand{\arraystretch}{1.3}
\begin{tabular}{@{}lcccccccccc@{}}
\toprule
$n$ & $1$ & $2$ & $3$ & $4$ & $5$ & $6$ & $7$ & $8$ & $9$ & $10$ \\
\midrule
$2^n$ & $2$ & $4$ & $8$ & $16$ & $32$ & $64$ & $128$ & $256$ & $512$ & $1024$ \\
$3^n$ & $3$ & $9$ & $27$ & $81$ & $243$ & $729$ & $2187$ & & & \\
\bottomrule
\end{tabular}
\end{center}
Ainsi, dans les problèmes où une quantité double chaque heure ou chaque année ($q = 2$), la somme sur $n$ périodes vaut $S_n = u_1(2^n - 1)$ :
\begin{itemize}
    \item Sur $7$ heures ($n = 7$) : $2^7 - 1 = 128 - 1 = \mathbf{127}$.
    \item Sur $10$ ans ($n = 10$) : $2^{10} - 1 = 1024 - 1 = \mathbf{1023}$.
\end{itemize}
\end{astuce}

\begin{exemple}[Les six problèmes de suites géométriques des annales (2021--2026)]
\begin{enumerate}
    \item \textbf{Septembre 2021, Q1.3 :} \emph{Donner la valeur du premier terme $u_1$ et la raison $q$ d'une suite géométrique dont les $11$\ieme{} et $21$\ieme{} termes valent respectivement $3\,072$ et $3\,145\,728$.}\\
    Utilisons la relation entre $u_{21}$ et $u_{11}$ :
    \[
        u_{21} = u_{11} \cdot q^{21-11} \iff q^{10} = \frac{u_{21}}{u_{11}} = \frac{3\,145\,728}{3\,072} = 1\,024 = 2^{10},
    \]
    ce qui donne $q = 2$ (ou $q = \pm 2$). Ensuite, à partir de $u_{11} = u_1 \cdot q^{10} = 1024\,u_1 = 3072$ :
    \[
        u_1 = \frac{3072}{1024} = 3 \implies \boxed{u_1 = 3 \quad \text{et} \quad q = 2}.
    \]

    \item \textbf{Septembre 2022, Q1.5 :} \emph{La somme des six premiers termes d'une série géométrique de raison $3/2$ vaut $1\,330$. Quel est le premier terme de cette série ?}\\
    Appliquons la formule de la somme avec $n = 6$, $q = \dfrac{3}{2} > 1$ et $S_6 = 1330$. Comme $q - 1 = \dfrac{1}{2}$ et $q^6 - 1 = \dfrac{3^6}{2^6} - 1 = \dfrac{729 - 64}{64} = \dfrac{665}{64}$ :
    \[
        S_6 = u_1 \cdot \frac{\left(\frac{3}{2}\right)^6 - 1}{\frac{3}{2} - 1} = u_1 \cdot \frac{665 / 64}{1 / 2} = u_1 \cdot \frac{665}{32} = 1330.
    \]
    Or $1330 = 2 \times 665$ ! En simplifiant par $665$, il vient $\dfrac{u_1}{32} = 2 \iff \boxed{u_1 = 64}$.

    \item \textbf{Juillet 2023 -- Série 1, Q1.5 (L'attraction du parc touristique) :} \emph{Le nombre de personnes ayant visité l'attraction pendant l'heure écoulée doublait par rapport à l'heure précédente. À la fin de la journée, qui a duré $7$ heures, un total de $12\,700$ personnes ont visité l'attraction. Combien de personnes étaient présentes lors de la première heure ?}\\
    C'est une suite géométrique de raison $q = 2$ sur $n = 7$ heures, avec $S_7 = 12\,700$ :
    \[
        S_7 = u_1 \cdot \frac{2^7 - 1}{2 - 1} = u_1(128 - 1) = 127\,u_1 = 12\,700 \implies \boxed{u_1 = \frac{12\,700}{127} = 100\text{ personnes}}.
    \]

    \item \textbf{Juillet 2024, Q1.6 (La collection de pierres de Laurent) :} \emph{La première année, Laurent achète un certain nombre de pierres, puis chaque année suivante, il double la quantité achetée l'année précédente. Au bout de dix ans, il possède un total de $5\,115$ pierres. Combien de pierres a-t-il achetées la première année ?}\\
    Suite géométrique de raison $q = 2$ sur $n = 10$ ans, avec $S_{10} = 5115$ :
    \[
        S_{10} = u_1(2^{10} - 1) = 1023\,u_1 = 5115 \implies \boxed{u_1 = \frac{5115}{1023} = 5\text{ pierres}}.
    \]

    \item \textbf{Septembre 2025, Q1.7 :} \emph{Soit une série géométrique de raison $3/2$. La somme des cinq premiers termes donne $211$. Quel est le premier terme de cette série ?}\\
    Avec $n = 5$ et $q = \dfrac{3}{2}$, on a $q^5 - 1 = \dfrac{3^5}{2^5} - 1 = \dfrac{243 - 32}{32} = \dfrac{211}{32}$ et $q - 1 = \dfrac{1}{2}$. D'où :
    \[
        S_5 = u_1 \cdot \frac{211 / 32}{1 / 2} = u_1 \cdot \frac{211}{16} = 211 \implies \boxed{u_1 = 16}.
    \]

    \item \textbf{Septembre 2026, Q1.7 (Les prix du concours) :} \emph{Lors d'un concours, le premier classé reçoit un prix de $5\,000$\,€ et les suivants reçoivent chacun un prix égal aux $3/5$ du précédent. Calculez le montant total des prix pour les cinq premiers (sachant que $6^5 = 7\,776$).}\\
    Ici $u_1 = 5000$, $q = \dfrac{3}{5} = 0,6 < 1$ et $n = 5$. Pourquoi l'énoncé donne-t-il $6^5 = 7776$ ? Parce que $q = 0,6 = \dfrac{6}{10}$, donc $q^5 = 0,6^5 = \dfrac{6^5}{10^5} = 0,07776$ !
    Appliquons la formule de la somme avec $1 - q = 1 - \dfrac{3}{5} = \dfrac{2}{5}$ :
    \[
        S_5 = 5000 \cdot \frac{1 - (3/5)^5}{1 - 3/5} = \frac{5000}{2/5}\bigl(1 - 0,07776\bigr) = 12\,500 \times 0,92224 = \frac{92\,224}{8} = \boxed{11\,528\,\text{€}}.
    \]
    \emph{Calcul fractionnaire encore plus rapide sans décimales :} Comme $(3/5)^5 = \frac{243}{3125}$, on a $S_5 = 5000 \cdot \frac{5}{2} \cdot \frac{3125 - 243}{3125} = \frac{25\,000}{2 \times 3125} \times 2882 = 4 \times 2882 = 11\,528$\,€ !
\end{enumerate}
\end{exemple}

% ----------------------------------------------------------------------------
\section{Exercices d'entraînement du Chapitre 6}
% ----------------------------------------------------------------------------

\emph{Les corrigés détaillés de tous ces exercices se trouvent au Chapitre~10.}

\begin{enumerate}[label=\textbf{Exercice 6.\arabic* :}, leftmargin=*]
    \item \textbf{(Suite arithmétique)} Le $4$\ieme{} terme d'une suite arithmétique vaut $14$ et le $11$\ieme{} terme vaut $35$. Déterminer la raison $r$, le premier terme $u_1$ ainsi que la somme $S_{20}$ des $20$ premiers termes.
    \item \textbf{(Suite géométrique)} La somme des trois premiers termes d'une suite géométrique strictement croissante vaut $26$ et leur produit vaut $216$. Déterminer ces trois termes.
    \item \textbf{(Problème mixte : puits de forage)} Pour forer un puits, une entreprise facture $100$\,€ pour le premier mètre, puis chaque mètre supplémentaire coûte $15$\,€ de plus que le mètre précédent. Quelle profondeur maximale peut-on forer avec un budget de $5\,200$\,€ ?
    \item \textbf{(UCLouvain -- Juillet 2021, Q3 : Le robot-tondeuse)} Une tondeuse automatique tond une pelouse à la vitesse $v_t = 1\text{ dm}^3/\text{h}$, tandis que l'herbe pousse en continu à la vitesse $v_p = 0,1\text{ dm}^3/\text{h}$. La batterie a une capacité initiale $C_{b,0} = 100\text{ Wh}$ au cycle $n = 1$, mais perd $d_b = 0,5\,\%$ de sa capacité maximale à chaque cycle : $C_b(n) = C_{b,0}(1 - d_b)^{n-1}$. La tondeuse consomme une puissance $p_t = 5\text{ W}$ pendant la tonte et retourne se charger dès qu'il ne lui reste plus que $C_{b,m} = 5\text{ Wh}$. Le temps de charge est fixe et vaut $t_c = 1\text{ h}$.
    \begin{enumerate}[label=(\alph*)]
        \item Exprimer la durée de tonte $t_t(n)$ au cycle $n$, puis déterminer l'équation donnant le cycle limite $n_\ell$ à partir duquel le temps d'attente après la charge tombe à zéro (la tondeuse ne parvient plus à compenser la pousse de l'herbe).
        \item Exprimer sous forme close (via la somme d'une suite géométrique) le temps total d'utilisation sur les $n_\ell$ cycles.
    \end{enumerate}
\end{enumerate}

% ============================================================================
% CHAPITRE 7 : DISCUSSION D'ÉQUATIONS ET D'INÉQUATIONS PARAMÉTRIQUES
% ============================================================================
\chapter{Discussion d'Équations et d'Inéquations Paramétriques}

La \textbf{Question~2} de l'épreuve d'Algèbre de l'examen d'admission à l'EPL (UCLouvain) est dédiée à la résolution et à la discussion complète dans $\R$ d'une inéquation (ou d'une équation) dépendant d'un \textbf{paramètre réel} noté $m$ ou $a$. Contrairement à la Question~1, ici \textbf{toute la démarche et la clarté de l'explication sont évaluées} : conditions d'existence, factorisation, étude du signe du discriminant $\Delta(m)$, comparaison de l'ordre des valeurs critiques en fonction du paramètre, tableaux de signes et conclusion récapitulative.

Ce chapitre vous donne une méthodologie infaillible en 5 étapes, illustrée par \textbf{l'intégralité des Questions~2 tombées entre 2021 et 2026}.

% ----------------------------------------------------------------------------
\section{Méthodologie générale en 5 étapes pour les inéquations rationnelles paramétriques}
% ----------------------------------------------------------------------------

Face à une inéquation rationnelle dépendant d'un paramètre réel $m \in \R$, suivez scrupuleusement ce plan de bataille :

\begin{theoreme}[Les 5 étapes d'une discussion d'inéquation rationnelle paramétrique]
\begin{enumerate}
    \item \textbf{Étape 1 -- Mise sous forme standard (membre de droite égal à $0$) :}
    Ramenez tous les termes à gauche et réduisez au même dénominateur pour obtenir un unique quotient comparé à $0$ :
    \[
        \boxed{\frac{N(x, m)}{D(x, m)} \ge 0 \qquad \left(\text{ou } > 0,\; \le 0,\; < 0\right)}.
    \]
    \textbf{Règle absolue :} Ne multipliez \textbf{JAMAIS} les deux membres d'une inéquation par un dénominateur $D(x,m)$ dont le signe dépend de $x$ ou de $m$ !
    \item \textbf{Étape 2 -- Conditions d'existence (Valeurs interdites) :}
    Résolvez $D(x, m) = 0$ pour identifier toutes les \textbf{valeurs interdites} de $x$ (qui seront marquées par une double barre $\|$ ou une croix $\times$ dans le tableau de signes et \textbf{toujours exclues} de l'ensemble des solutions $S$, même si l'inégalité est large $\ge 0$). Si un dénominateur s'annule pour une valeur particulière du paramètre (ex. $m = 0$ dans $\frac{1}{mx}$), traitez ce cas d'emblée ($S = \varnothing$).
    \item \textbf{Étape 3 -- Recherche des racines (valeurs critiques) du numérateur et du dénominateur :}
    Factorisez $N(x,m)$ et $D(x,m)$. Si un facteur est du second degré $ax^2 + bx + c$, calculez son discriminant $\Delta(m) = b^2 - 4ac$ et étudiez le \textbf{signe de $\Delta(m)$ en fonction de $m$} :
    \begin{itemize}
        \item Pour les valeurs de $m$ où $\Delta(m) < 0$, ce trinôme n'a pas de racine réelle et garde un \textbf{signe constant} (celui de $a$) sur tout $\R$.
        \item Pour les valeurs de $m$ où $\Delta(m) \ge 0$, notez ses racines $x_1(m) \le x_2(m)$ (\emph{attention : si le coefficient dominant $a$ dépend de $m$ et change de signe, comme dans $mx^2 - mx - 1$, l'ordre entre $\frac{-b-\sqrt{\Delta}}{2m}$ et $\frac{-b+\sqrt{\Delta}}{2m}$ s'inverse lorsque $m < 0$ !}).
    \end{itemize}
    \item \textbf{Étape 4 -- Classement (ordre) des valeurs critiques selon les plages de $m$ et Tableaux de signes :}
    Pour placer les valeurs critiques dans l'ordre croissant sur la première ligne du tableau de signes, \textbf{comparez les racines entre elles} (en résolvant $x_{\text{num}}(m) = x_{\text{dén}}(m)$ pour trouver les valeurs de basculement de $m$). Dressez ensuite un \textbf{tableau de signes propre pour chaque cas}.
    \item \textbf{Étape 5 -- Synthèse finale encadrée :}
    Récapitulez clairement l'ensemble des solutions $S$ pour chaque intervalle ou valeur particulière du paramètre $m$.
\end{enumerate}
\end{theoreme}

% ----------------------------------------------------------------------------
\section{Inéquations rationnelles paramétriques : études complètes des annales}
% ----------------------------------------------------------------------------

\subsection{Cas où un facteur du premier degré change de place par rapport aux racines}

\begin{exemple}[Inéquation rationnelle avec paramètre positif -- UCLouvain, Septembre 2026, Q2]
\textbf{Énoncé :} Résoudre dans $\R$, en discutant selon le paramètre réel strictement positif $m > 0$, l'inéquation :
\[
    \frac{x - m}{x^3 - x} > 0.
\]
\textbf{Résolution détaillée :}
\begin{enumerate}
    \item \textbf{Factorisation et domaine d'existence :} Le dénominateur se factorise en $x^3 - x = x(x^2 - 1) = x(x - 1)(x + 1)$. L'inéquation s'écrit :
    \[
        \frac{x - m}{x(x - 1)(x + 1)} > 0.
    \]
    Les valeurs interdites (qui annulent le dénominateur) sont $x \in \{-1, 0, 1\}$, donc $\dom = \R \setminus \{-1, 0, 1\}$.
    \item \textbf{Valeurs critiques et ordre en fonction de $m > 0$ :}
    Le dénominateur a trois racines simples fixes : $-1 < 0 < 1$. Le numérateur s'annule en $x = m$.
    Puisque l'énoncé précise que $\mathbf{m > 0}$, la racine $m$ est toujours strictement à droite de $0$. Il suffit donc de comparer $m$ avec la racine $1$, ce qui donne \textbf{trois cas} :
    \begin{itemize}
        \item \textbf{1\ier{} cas : $0 < m < 1$.} L'ordre des valeurs critiques est $-1 < 0 < m < 1$ :
        \begin{center}
        \renewcommand{\arraystretch}{1.2}
        \begin{tabular}{c|ccccccccccc}
            $x$ & $-\infty$ & & $-1$ & & $0$ & & $m$ & & $1$ & & $+\infty$ \\ \hline
            $x - m$ & & $-$ & $-$ & $-$ & $-$ & $-$ & $0$ & $+$ & $+$ & $+$ & \\
            $x(x-1)(x+1)$ & & $-$ & $0$ & $+$ & $0$ & $-$ & $-$ & $-$ & $0$ & $+$ & \\ \hline
            Quotient & & $+$ & $\|$ & $-$ & $\|$ & $+$ & $0$ & $-$ & $\|$ & $+$ &
        \end{tabular}
        \end{center}
        Comme l'inégalité est stricte ($> 0$), on prend les régions marquées $+$ en excluant toutes les bornes :
        \[
            S = \,]-\infty, -1[ \,\cup\, ]0, m[ \,\cup\, ]1, +\infty[.
        \]

        \item \textbf{2\ieme{} cas : $m = 1$.} Le facteur $(x - 1)$ apparaît en haut et en bas ! Pour $x \neq 1$, la fraction se simplifie en $\dfrac{1}{x(x + 1)} > 0$, qui est strictement positive pour $x \in \,]-\infty, -1[ \,\cup\, ]0, +\infty[$. \textbf{Attention :} $x = 1$ reste une valeur interdite du dénominateur initial, il faut donc l'exclure de $]0, +\infty[$ :
        \[
            S = \,]-\infty, -1[ \,\cup\, ]0, 1[ \,\cup\, ]1, +\infty[.
        \]

        \item \textbf{3\ieme{} cas : $m > 1$.} L'ordre des valeurs critiques est $-1 < 0 < 1 < m$ :
        \begin{center}
        \renewcommand{\arraystretch}{1.2}
        \begin{tabular}{c|ccccccccccc}
            $x$ & $-\infty$ & & $-1$ & & $0$ & & $1$ & & $m$ & & $+\infty$ \\ \hline
            $x - m$ & & $-$ & $-$ & $-$ & $-$ & $-$ & $-$ & $-$ & $0$ & $+$ & \\
            $x(x-1)(x+1)$ & & $-$ & $0$ & $+$ & $0$ & $-$ & $0$ & $+$ & $+$ & $+$ & \\ \hline
            Quotient & & $+$ & $\|$ & $-$ & $\|$ & $+$ & $\|$ & $-$ & $0$ & $+$ &
        \end{tabular}
        \end{center}
        On lit directement : $S = \,]-\infty, -1[ \,\cup\, ]0, 1[ \,\cup\, ]m, +\infty[$.
    \end{itemize}
\end{enumerate}
\textbf{Conclusion :}
\[
    \boxed{\begin{cases}
        \text{Si } 0 < m < 1 : & S = \,]-\infty, -1[ \,\cup\, ]0, m[ \,\cup\, ]1, +\infty[ \\
        \text{Si } m = 1 : & S = \,]-\infty, -1[ \,\cup\, ]0, 1[ \,\cup\, ]1, +\infty[ \\
        \text{Si } m > 1 : & S = \,]-\infty, -1[ \,\cup\, ]0, 1[ \,\cup\, ]m, +\infty[
    \end{cases}}
\]
\end{exemple}

\subsection{Cas où un trinôme du second degré dépend du paramètre}

\begin{exemple}[Mise sous forme standard et discussion -- UCLouvain, Juillet 2024, Q2]
\textbf{Énoncé :} Résoudre dans $\R$ et discuter en fonction du paramètre réel $m$ l'inéquation suivante :
\[
    \frac{(m - 1)(x - 1)}{x + m - 1} \le x.
\]
\textbf{Résolution détaillée :}
\begin{enumerate}
    \item \textbf{Étape 1 (Mise sous forme standard) :} On fait passer $x$ à gauche et on réduit au même dénominateur (sans multiplier par $x + m - 1$ !) :
    \[
        \frac{(m - 1)(x - 1) - x(x + m - 1)}{x + m - 1} \le 0 \iff \frac{-x^2 - m + 1}{x + m - 1} \le 0.
    \]
    En multipliant les deux membres par $-1$ (ce qui \textbf{inverse le sens de l'inégalité} : $\le 0$ devient $\ge 0$) :
    \[
        \boxed{\frac{x^2 + m - 1}{x + m - 1} \ge 0}.
    \]
    \item \textbf{Étape 2 (Condition d'existence) :} Le dénominateur s'annule en $x_D = 1 - m$, donc $\text{C.E. } : \mathbf{x \neq 1 - m}$.
    \item \textbf{Étape 3 (Racines du numérateur $N(x) = x^2 - (1 - m)$) :}
    \begin{itemize}
        \item Si $1 - m < 0 \iff \mathbf{m > 1}$ : $x^2 + (m - 1) > 0$ pour tout $x \in \R$ (aucune racine réelle, numérateur toujours strictement positif).
        \item Si $1 - m = 0 \iff \mathbf{m = 1}$ : $x^2 = 0$ admet une racine double $x = 0$ (qui coïncide avec la valeur interdite $1 - m = 0$).
        \item Si $1 - m > 0 \iff \mathbf{m < 1}$ : le numérateur possède deux racines réelles opposées $x_1 = -\sqrt{1 - m}$ et $x_2 = +\sqrt{1 - m}$.
    \end{itemize}
    \item \textbf{Étape 4 (Comparaison de la racine du dénominateur $1 - m$ avec $x_1 = -\sqrt{1-m}$ et $x_2 = \sqrt{1-m}$ pour $m < 1$) :}
    Posons $u = \sqrt{1 - m} > 0$, de sorte que $x_1 = -u$, $x_2 = u$, et la valeur interdite du dénominateur vaut $x_D = 1 - m = u^2 > 0$.
    Comme $u > 0$, on a toujours $-u < 0 < u^2$, donc $x_1 < x_D$. En revanche, la position de $x_D = u^2$ par rapport à $x_2 = u$ dépend de la comparaison de $u$ avec $1$ ($u^2 = u \iff u = 1 \iff 1 - m = 1 \iff m = 0$) :
    \begin{itemize}
        \item Si $\mathbf{0 < m < 1}$ ($0 < u < 1$, donc $u^2 < u$) : l'ordre est $\mathbf{-\sqrt{1 - m} < 1 - m < \sqrt{1 - m}}$.
        \item Si $\mathbf{m = 0}$ ($u = 1$) : $x_D = 1$ coïncide avec $x_2 = 1$ (ordre : $-1 < 1$).
        \item Si $\mathbf{m < 0}$ ($u > 1$, donc $u < u^2$) : l'ordre est $\mathbf{-\sqrt{1 - m} < \sqrt{1 - m} < 1 - m}$.
    \end{itemize}
    Dressons le tableau de signes pour le cas $0 < m < 1$ (et de même pour $m < 0$) :
    \begin{center}
    \renewcommand{\arraystretch}{1.2}
    \begin{tabular}{c|ccccccccc}
        $x$ & $-\infty$ & & $-\sqrt{1-m}$ & & $1-m$ & & $\sqrt{1-m}$ & & $+\infty$ \\ \hline
        $x^2 + m - 1$ & & $+$ & $0$ & $-$ & $-$ & $-$ & $0$ & $+$ & \\
        $x + m - 1$ & & $-$ & $-$ & $-$ & $0$ & $+$ & $+$ & $+$ & \\ \hline
        Quotient & & $-$ & $0$ & $+$ & $\|$ & $-$ & $0$ & $+$ &
    \end{tabular}
    \end{center}
\end{enumerate}
\textbf{Étape 5 (Synthèse finale des 5 cas) :}
\[
    \boxed{\begin{cases}
        \text{Si } m < 0 : & S = \bigl[-\sqrt{1 - m},\, \sqrt{1 - m}\bigr] \cup \,]1 - m, +\infty[ \\
        \text{Si } m = 0 : & S = [-1, 1[ \,\cup\, ]1, +\infty[ \\
        \text{Si } 0 < m < 1 : & S = \bigl[-\sqrt{1 - m},\, 1 - m\bigr[ \cup \bigl[\sqrt{1 - m},\, +\infty\bigr[ \\
        \text{Si } m = 1 : & S = \,]0, +\infty[ \\
        \text{Si } m > 1 : & S = \,]1 - m, +\infty[
    \end{cases}}
\]
\end{exemple}

\begin{exemple}[Cinq autres discussions rationnelles et valeurs absolues des annales (2021--2025)]
\begin{enumerate}
    \item \textbf{Septembre 2024, Q2 :} \emph{Résoudre et discuter en fonction de $m \in \R$ : $\dfrac{x^2 - 2mx + 1}{x - m} \ge 0$.}\\
    \begin{itemize}
        \item \textbf{C.E. :} $x \neq m$.
        \item \textbf{Numérateur $N(x) = x^2 - 2mx + 1$ :} $\Delta = 4m^2 - 4 = 4(m^2 - 1)$.
        \item Si $\mathbf{-1 < m < 1}$ ($\Delta < 0$) : $N(x) > 0$ pour tout $x \in \R$. Le quotient a le signe de $x - m$, donc $S = \,]m, +\infty[$.
        \item Si $\mathbf{m = \pm 1}$ ($\Delta = 0$) : $N(x) = (x - m)^2 \ge 0$, donc pour $x \neq m$ le quotient vaut $x - m > 0 \iff S = \,]m, +\infty[$.
        \item Si $\mathbf{m < -1 \text{ ou } m > 1}$ ($\Delta > 0$) : les racines du numérateur sont $x_1 = m - \sqrt{m^2 - 1}$ et $x_2 = m + \sqrt{m^2 - 1}$. Comme $\sqrt{m^2 - 1} > 0$, la valeur interdite $m$ est \textbf{toujours située pile au milieu des deux racines} : $\mathbf{x_1 < m < x_2}$ ! Le tableau de signes est identique pour $m < -1$ et $m > 1$ :
    \end{itemize}
    \[
        \boxed{\begin{cases}
            \text{Si } -1 \le m \le 1 : & S = \,]m, +\infty[ \\
            \text{Si } m < -1 \text{ ou } m > 1 : & S = \bigl[m - \sqrt{m^2 - 1},\, m\bigr[ \,\cup\, \bigl[m + \sqrt{m^2 - 1},\, +\infty\bigr[
        \end{cases}}
    \]

    \item \textbf{Juillet 2022, Q2 :} \emph{Résoudre et discuter selon $m \in \R$ : $\dfrac{x^2 + 2x - 2m}{x^2 - m^2} > 0$.}\\
    \textbf{C.E. :} $x \notin \{-m, m\}$. Le numérateur a pour discriminant $\Delta = 4 + 8m = 4(2m + 1)$ :
    \begin{itemize}
        \item Si $\mathbf{m < -1/2}$ ($\Delta < 0$) : numérateur toujours $> 0$. Comme $m < 0$, on a $m < -m$, et $x^2 - m^2 > 0$ à l'extérieur de $[m, -m]$ : $S = \,]-\infty, m[ \,\cup\, ]-m, +\infty[$.
        \item Si $\mathbf{m \ge -1/2}$ ($\Delta \ge 0$) : le numérateur a pour racines $x_1 = -1 - \sqrt{2m+1}$ et $x_2 = -1 + \sqrt{2m+1}$. Quand $x_1$ ou $x_2$ croisent-elles $\pm m$ ?
        L'équation $(x \pm m) = 0$ dans $x^2 + 2x - 2m = 0$ donne $m^2 + 2(\pm m) - 2m = 0$, soit $m^2 = 0 \implies m = 0$ (pour $x = m$) et $m^2 - 4m = 0 \implies m = 0 \text{ ou } m = 4$ (pour $x = -m$). Cela découpe la discussion en trois sous-cas :
        \begin{itemize}
            \item Pour $\mathbf{-1/2 \le m < 0}$ (ordre $x_1 < x_2 < m < -m$) : $S = \,]-\infty, x_1[ \,\cup\, ]x_2, m[ \,\cup\, ]-m, +\infty[$.
            \item Pour $\mathbf{0 \le m < 4}$ (ordre $x_1 < -m \le x_2 < m$, avec $x_2 = m = 0$ si $m=0$) : $S = \,]-\infty, x_1[ \,\cup\, ]-m, x_2[ \,\cup\, ]m, +\infty[$ (et pour $m = 0$ : $]-\infty, -2[ \,\cup\, ]0, +\infty[$).
            \item Pour $\mathbf{m \ge 4}$ (ordre $-m \le x_1 < x_2 < m$) : $S = \,]-\infty, -m[ \,\cup\, ]x_1, x_2[ \,\cup\, ]m, +\infty[$ (et pour $m = 4$, $x_1 = -m = -4$ est exclu : $]-\infty, -4[ \,\cup\, ]-4, 2[ \,\cup\, ]4, +\infty[$).
        \end{itemize}
    \end{itemize}

    \item \textbf{Juillet 2025, Q2 (Avec valeur absolue $|x|$) :} \emph{Résoudre et discuter selon $m \in \R$ : $\dfrac{|x| - m}{x^2 - mx + 1} \ge 0$.}\\
    \begin{itemize}
        \item \textbf{Numérateur $N(x) = |x| - m$ :} Si $m < 0$, alors $|x| - m > 0$ pour tout $x \in \R$ (toujours strictement positif). Si $m \ge 0$, $|x| - m \ge 0 \iff x \le -m \text{ ou } x \ge m$ (racines $-m$ et $m$).
        \item \textbf{Dénominateur $D(x) = x^2 - mx + 1$ :} $\Delta = m^2 - 4$.
        Si $-2 < m < 2$, $\Delta < 0$ donc $D(x) > 0$ sur tout $\R$. Si $m \le -2$ ou $m \ge 2$, $D(x)$ a deux racines $x_{1,2} = \dfrac{m \pm \sqrt{m^2 - 4}}{2}$ (pour $m \ge 2$, on vérifie que $-m < 0 < x_1 \le x_2 < m$).
    \end{itemize}
    Cela donne les 4 plages de discussion :
    \[
        \boxed{\begin{cases}
            \text{Si } m \le -2 : & S = \left]-\infty,\, \frac{m - \sqrt{m^2-4}}{2}\right[ \cup \left]\frac{m + \sqrt{m^2-4}}{2},\, +\infty\right[ \\
            \text{Si } -2 < m < 0 : & S = \R \\
            \text{Si } 0 \le m < 2 : & S = \,]-\infty, -m] \cup [m, +\infty[ \\
            \text{Si } m \ge 2 : & S = \,]-\infty, -m] \cup \left]\frac{m - \sqrt{m^2-4}}{2},\, \frac{m + \sqrt{m^2-4}}{2}\right[ \cup [m, +\infty[
        \end{cases}}
    \]

    \item \textbf{Septembre 2025, Q2 (Coefficient dominant $m$ au carré !) :} \emph{Résoudre et discuter selon $m \in \R$ : $x - 1 \ge \dfrac{1}{mx}$.}\\
    \begin{itemize}
        \item \textbf{Cas $m = 0$ :} Division par zéro dans l'énoncé $\implies \mathbf{S = \varnothing}$.
        \item \textbf{Pour $m \neq 0$ et $x \neq 0$ :} Mise sous forme standard :
        \[
            x - 1 - \frac{1}{mx} \ge 0 \iff \frac{mx^2 - mx - 1}{mx} \ge 0 \qquad \text{avec } P(x) = mx^2 - mx - 1.
        \]
        Le discriminant de $P(x)$ vaut $\Delta = m^2 + 4m = m(m + 4)$.
        \item \textbf{Si $-4 < m < 0$ ($\Delta < 0$) :} Comme $m < 0$, $P(x) < 0$ pour tout $x \in \R$, et le dénominateur $mx$ est $< 0$ lorsque $x > 0$. Le quotient $\frac{-}{-}$ est donc positif pour $\mathbf{S = \,]0, +\infty[}$.
        \item \textbf{Si $m > 0$ ($\Delta > 0$) :} Deux racines de signes contraires (car produit $-\frac{1}{m} < 0$) : $x_1 = \frac{m - \sqrt{m^2+4m}}{2m} < 0 < x_2 = \frac{m + \sqrt{m^2+4m}}{2m}$. Le tableau de signes donne $\mathbf{S = [x_1, 0[ \,\cup\, [x_2, +\infty[}$.
        \item \textbf{Si $m \le -4$ ($\Delta \ge 0$ et $m < 0$) :} \emph{Attention à l'ordre des racines quand on divise par $2m < 0$ !} La plus petite racine est $x_1 = \frac{m + \sqrt{m^2+4m}}{2m}$ et la plus grande est $x_2 = \frac{m - \sqrt{m^2+4m}}{2m}$, avec $0 < x_1 \le x_2 < 1$. Le tableau de signes (avec coefficient dominant $m < 0$) donne $\mathbf{S = \,]0, x_1] \cup [x_2, +\infty[}$ (et pour $m = -4$, $x_1 = x_2 = \frac{1}{2}$, donc $S = \,]0, +\infty[$).
    \end{itemize}

    \item \textbf{Juillet 2023 -- Série 1, Q2 (Valeur absolue paramétrique) :} \emph{Résoudre et discuter selon $a \in \R$ : $|x^2 - a| < 1$.}\\
    D'après le Chapitre~2, $|x^2 - a| < 1 \iff -1 < x^2 - a < 1 \iff \mathbf{a - 1 < x^2 < a + 1}$.
    \begin{itemize}
        \item \textbf{Si $a \le -1$ :} Alors $a + 1 \le 0$, or un carré réel $x^2 \ge 0$ ne peut jamais être strictement négatif $\implies \mathbf{S = \varnothing}$.
        \item \textbf{Si $-1 < a < 1$ :} Alors $a - 1 < 0$ (donc $x^2 > a - 1$ est toujours vrai !) et $a + 1 > 0$. Il reste uniquement $x^2 < a + 1 \iff \mathbf{S = \bigl]-\sqrt{a + 1},\, \sqrt{a + 1}\bigr[}$.
        \item \textbf{Si $a \ge 1$ :} Alors $0 \le a - 1 < a + 1$. En prenant la racine carrée des membres positifs : $\sqrt{a - 1} < |x| < \sqrt{a + 1}$, ce qui donne deux intervalles symétriques :
        \[
            \mathbf{S = \bigl]-\sqrt{a + 1},\, -\sqrt{a - 1}\bigr[ \,\cup\, \bigl]\sqrt{a - 1},\, \sqrt{a + 1}\bigr[}.
        \]
    \end{itemize}

    \item \textbf{Juillet 2023 -- Série 2, Q2 (Équation paramétrique avec condition sur les racines) :} \emph{Pour quelles valeurs du paramètre réel $a$ l'équation $\dfrac{x^2 + 2x - a}{x - a} = 0$ admet-elle deux racines réelles $x_1$ et $x_2$ telles que $x_1 x_2 > -1$ ?}\\
    Il faut réunir \textbf{trois conditions simultanées} :
    \begin{itemize}
        \item \emph{Condition 1 (Existence de deux racines réelles au numérateur $x^2 + 2x - a = 0$) :} $\Delta = 4 + 4a \ge 0 \iff \mathbf{a \ge -1}$ (ou $a > -1$ si les deux racines sont supposées distinctes), de racines $x_{1,2} = -1 \pm \sqrt{1 + a}$.
        \item \emph{Condition 2 (Produit des racines $x_1 x_2 > -1$) :} Par la relation Somme-Produit de Viète, le produit des racines du trinôme $x^2 + 2x - a$ vaut $x_1 x_2 = \frac{c}{1} = -a$. On impose donc $-a > -1 \iff \mathbf{a < 1}$.
        \item \emph{Condition 3 (Aucune des deux racines ne doit annuler le dénominateur $x - a \neq 0$) :} Il est interdit que $x = a$ soit racine du numérateur ! En remplaçant $x = a$ dans $x^2 + 2x - a \neq 0$ :
        \[
            a^2 + 2a - a \neq 0 \iff a^2 + a \neq 0 \iff a(a + 1) \neq 0 \iff \mathbf{a \neq 0 \text{ et } a \neq -1}.
        \]
    \end{itemize}
    L'intersection des trois conditions ($a \ge -1$, $a < 1$ et $a \notin \{-1, 0\}$) donne :
    \[
        \boxed{a \in \,]-1, 0[ \,\cup\, ]0, 1[}.
    \]
\end{enumerate}
\end{exemple}

% ----------------------------------------------------------------------------
\section{Inéquations logarithmiques paramétriques (dont base variable $a - 1$)}
% ----------------------------------------------------------------------------

Trois sessions récentes (\textbf{Septembre 2022 Q2}, \textbf{Septembre 2023 Q2} et \textbf{Juillet 2026 Q2}) ont posé à la Question~2 une inéquation logarithmique paramétrique.

\begin{theoreme}[Méthode de résolution des inéquations logarithmiques paramétriques]
\begin{enumerate}
    \item \textbf{Cas d'une base fixe $b > 1$ (comme $\ln$ ou $\log_2$) :}
    Puisque $\log_b$ est strictement croissante pour $b > 1$ et que la C.E. impose $U(x, a) > 0$, une inéquation de la forme $\log_b\bigl(U(x, a)\bigr) < k$ se traduit par une \textbf{double inégalité} (intersection de la C.E. et de l'inéquation) :
    \[
        \boxed{\log_b\bigl(U(x, a)\bigr) < k \iff 0 < U(x, a) < b^k}.
    \]
    En particulier, pour $k = 0$ : $\boxed{\ln\bigl(U(x,a)\bigr) < 0 \iff 0 < U(x, a) < 1}$.
    \item \textbf{Cas d'une base dépendante du paramètre $\log_{B(a)}\bigl(U(x, a)\bigr) > 0$ (ex. $B(a) = a - 1$) :}
    Passez au logarithme népérien par la \textbf{formule de changement de base} $\dfrac{\ln\bigl(U(x,a)\bigr)}{\ln\bigl(B(a)\bigr)} > 0$.
    \begin{itemize}
        \item \textbf{Si $B(a) \le 0$ ou $B(a) = 1$ :} La base du logarithme n'est pas valide $\implies \mathbf{S = \varnothing}$.
        \item \textbf{Si $B(a) > 1$ ($\ln B(a) > 0$, base croissante) :} $\log_{B(a)}\bigl(U(x,a)\bigr) > 0 \iff \boxed{U(x, a) > 1}$.
        \item \textbf{Si $0 < B(a) < 1$ ($\ln B(a) < 0$, base décroissante) :} Le sens s'inverse ! $\log_{B(a)}\bigl(U(x,a)\bigr) > 0 \iff \ln\bigl(U(x,a)\bigr) < 0 \iff \boxed{0 < U(x, a) < 1}$.
    \end{itemize}
\end{enumerate}
\end{theoreme}

\begin{exemple}[Les trois inéquations logarithmiques paramétriques des annales (2022, 2023, 2026)]
\begin{enumerate}
    \item \textbf{Septembre 2022, Q2 :} \emph{Résoudre dans $\R$ et discuter selon $a \in \R$ : $\ln(-x^2 + 2x + a) < 0$.}\\
    D'après le théorème, cette inéquation équivaut exactement à la double condition :
    \[
        \boxed{0 < -x^2 + 2x + a < 1} \iff \begin{cases} -x^2 + 2x + a > 0 & (1) \text{ [Condition d'existence]} \\ -x^2 + 2x + a - 1 < 0 & (2) \text{ [Inégalité } < \ln(1)] \end{cases}
    \]
    \begin{itemize}
        \item \textbf{Étude de $(1)$ : $-x^2 + 2x + a > 0$.} Le discriminant vaut $\Delta_1 = 4 + 4a = 4(1 + a)$. Comme le coefficient de $x^2$ est $-1 < 0$, ce trinôme n'est strictement positif qu'entre ses racines (à condition que $\Delta_1 > 0 \iff a > -1$) :
        \begin{itemize}
            \item Si $\mathbf{a \le -1}$ : $(1)$ n'est jamais vérifiée $\implies \mathbf{S = \varnothing}$.
            \item Si $\mathbf{a > -1}$ : $(1) \iff x \in \bigl]1 - \sqrt{1 + a},\, 1 + \sqrt{1 + a}\bigr[$.
        \end{itemize}
        \item \textbf{Étude de $(2)$ : $-x^2 + 2x + (a - 1) < 0 \iff x^2 - 2x + (1 - a) > 0$.} Son discriminant vaut $\Delta_2 = 4 - 4(1 - a) = 4a$ :
        \begin{itemize}
            \item Si $\mathbf{-1 < a < 0}$ ($\Delta_2 < 0$) : $(2)$ est \textbf{toujours vérifiée} sur $\R$ ! L'intersection de $(1)$ et $(2)$ redonne tout l'intervalle de $(1)$ :
            \[
                \mathbf{S = \bigl]1 - \sqrt{1 + a},\, 1 + \sqrt{1 + a}\bigr[}.
            \]
            \item Si $\mathbf{a \ge 0}$ ($\Delta_2 \ge 0$) : $(2)$ est vérifiée à l'extérieur des racines $1 - \sqrt{a}$ et $1 + \sqrt{a}$, soit $x < 1 - \sqrt{a}$ ou $x > 1 + \sqrt{a}$. Comme $0 \le a < 1 + a \implies \sqrt{a} < \sqrt{1 + a}$, on a l'emboîtement $1 - \sqrt{1+a} < 1 - \sqrt{a} \le 1 + \sqrt{a} < 1 + \sqrt{1+a}$. L'intersection de $(1)$ et $(2)$ retire le segment central $[1 - \sqrt{a}, 1 + \sqrt{a}]$ :
            \[
                \mathbf{S = \bigl]1 - \sqrt{1 + a},\, 1 - \sqrt{a}\bigr[ \,\cup\, \bigl]1 + \sqrt{a},\, 1 + \sqrt{1 + a}\bigr[}.
            \]
        \end{itemize}
    \end{itemize}

    \item \textbf{Septembre 2023, Q2 :} \emph{Résoudre dans $\R$ et discuter selon $a \in \R$ : $\log_2\left(\dfrac{x + a}{x - 1}\right) < 1$.}\\
    Puisque la base $2 > 1$ est strictement croissante et $2^1 = 2$, l'inéquation équivaut au système ($x \neq 1$) :
    \[
        0 < \frac{x + a}{x - 1} < 2 \iff \begin{cases} \dfrac{x + a}{x - 1} > 0 & (1) \\[1.8ex] \dfrac{x + a}{x - 1} - 2 < 0 \iff \dfrac{-x + a + 2}{x - 1} < 0 \iff \dfrac{x - (a + 2)}{x - 1} > 0 & (2) \end{cases}
    \]
    Chaque fraction $\frac{x - r}{x - 1}$ est strictement positive à l'\textbf{extérieur} de l'intervalle borné par $r$ et $1$. Comparons les racines $-a$ et $a + 2$ avec $1$ (elles se croisent toutes les deux en $1$ lorsque $\mathbf{a = -1}$ !) :
    \begin{itemize}
        \item \textbf{Si $a = -1$ :} Alors $\frac{x - 1}{x - 1} = 1$ pour tout $x \neq 1$, et $\log_2(1) = 0 < 1$ est toujours vrai ! (Attention : $x = 1$ annule le dénominateur $x - 1$, donc il est exclu) : $\mathbf{S = \R \setminus \{1\} = \,]-\infty, 1[ \,\cup\, ]1, +\infty[}$.
        \item \textbf{Si $a < -1$ :} Alors $a + 2 < 1 < -a$.
        $(1) \iff x \in \,]-\infty, 1[ \,\cup\, ]-a, +\infty[$ et $(2) \iff x \in \,]-\infty, a + 2[ \,\cup\, ]1, +\infty[$.
        L'intersection des deux ensembles donne : $\mathbf{S = \,]-\infty, a + 2[ \,\cup\, ]-a, +\infty[}$.
        \item \textbf{Si $a > -1$ :} Alors $-a < 1 < a + 2$.
        $(1) \iff x \in \,]-\infty, -a[ \,\cup\, ]1, +\infty[$ et $(2) \iff x \in \,]-\infty, 1[ \,\cup\, ]a + 2, +\infty[$.
        L'intersection des deux ensembles donne : $\mathbf{S = \,]-\infty, -a[ \,\cup\, ]a + 2, +\infty[}$.
    \end{itemize}

    \item \textbf{Juillet 2026, Q2 (Base paramétrique $a - 1$) :} \emph{Résoudre dans $\R$ et discuter selon $a \in \R$ : $\log_{a-1}\bigl(x^2 - 2ax + 1\bigr) > 0$.}\\
    Par changement de base : $\dfrac{\ln(x^2 - 2ax + 1)}{\ln(a - 1)} > 0$.
    \textbf{Conditions sur la base $a - 1$ :} Il faut $a - 1 > 0 \iff a > 1$ et $a - 1 \neq 1 \iff a \neq 2$.
    \begin{itemize}
        \item \textbf{Cas 0 : $a \le 1$ ou $a = 2$.} Le logarithme en base $a - 1$ n'est pas défini $\implies \mathbf{S = \varnothing}$.
        \item \textbf{Cas 1 : $a > 2$ (donc $a - 1 > 1$, base croissante, $\ln(a - 1) > 0$).}\\
        L'inéquation équivaut à $\ln(x^2 - 2ax + 1) > 0 \iff x^2 - 2ax + 1 > 1$ (ce qui garantit automatiquement la C.E. $> 0$). En simplifiant :
        \[
            x^2 - 2ax > 0 \iff x(x - 2a) > 0 \iff \mathbf{S = \,]-\infty, 0[ \,\cup\, ]2a, +\infty[} \quad (\text{car } 2a > 4 > 0).
        \]
        \item \textbf{Cas 2 : $1 < a < 2$ (donc $0 < a - 1 < 1$, base décroissante, $\ln(a - 1) < 0$).}\\
        En multipliant par $\ln(a - 1) < 0$, le sens s'inverse : $\ln(x^2 - 2ax + 1) < 0$, ce qui impose la double inégalité :
        \[
            0 < x^2 - 2ax + 1 < 1 \iff \begin{cases} x^2 - 2ax + 1 > 0 & (1) \\ x^2 - 2ax < 0 \iff x(x - 2a) < 0 & (2) \end{cases}
        \]
        Puisque $a \in \,]1, 2[$, $(2)$ donne l'intervalle $\mathbf{0 < x < 2a}$.
        Pour $(1)$, le trinôme $x^2 - 2ax + 1$ a pour discriminant $\Delta = 4(a^2 - 1) > 0$ (car $a > 1$), et est strictement positif à l'extérieur de ses racines $x_1 = a - \sqrt{a^2 - 1}$ et $x_2 = a + \sqrt{a^2 - 1}$.
        Or, comme $0 < \sqrt{a^2 - 1} < \sqrt{a^2} = a$, on a toujours l'emboîtement $\mathbf{0 < a - \sqrt{a^2 - 1} < a + \sqrt{a^2 - 1} < 2a}$ !
        L'intersection de $]0, 2a[$ avec $]-\infty, x_1[ \,\cup\, ]x_2, +\infty[$ donne donc :
        \[
            \mathbf{S = \bigl]0,\, a - \sqrt{a^2 - 1}\bigr[ \,\cup\, \bigl]a + \sqrt{a^2 - 1},\, 2a\bigr[}.
        \]
    \end{itemize}
\end{enumerate}
\end{exemple}

% ----------------------------------------------------------------------------
\section{Exercices d'entraînement du Chapitre 7}
% ----------------------------------------------------------------------------

\emph{Les corrigés détaillés de tous ces exercices se trouvent au Chapitre~10.}

\begin{enumerate}[label=\textbf{Exercice 7.\arabic* :}, leftmargin=*]
    \item \textbf{(UCLouvain -- Juillet 2021, Q2.2)} Soit $m$ un paramètre réel strictement négatif ($m < 0$). Résoudre et discuter dans $\R$ l'inéquation suivante :
    \[
        \frac{x^4 - 4x^2 + 4}{x^2 - mx + 4} < 0.
    \]
    \item \textbf{(Inéquation rationnelle paramétrique type Question 2)} Résoudre dans $\R$ et discuter en fonction du paramètre réel $m \in \R$ l'inéquation :
    \[
        \frac{x^2 - (m + 1)x + m}{x^2 - 1} \le 0.
    \]
    \item \textbf{(Inéquation logarithmique à base paramétrique type Question 2)} Résoudre dans $\R$ et discuter en fonction du paramètre réel $a \in \R$ l'inéquation :
    \[
        \log_a\bigl(x^2 - ax\bigr) < 1.
    \]
\end{enumerate}

% ============================================================================
% CHAPITRE 8 : MISE EN ÉQUATIONS ET SYSTÈMES ALGÉBRIQUES (PROBLÈMES TEXTUELS)
% ============================================================================
\chapter{Mise en Équations et Systèmes Algébriques (Problèmes Textuels)}

La \textbf{Question~3} de l'épreuve d'Algèbre à l'examen d'admission de l'UCLouvain est une grande question de \textbf{modélisation mathématique} : un récit concret d'une demi-page à une page décrit une situation mettant souvent en scène \emph{Charles, Laurent, Aude, Sophie} ou un projet de l'EPL. Il s'agit de traduire méthodiquement chaque phrase de l'énoncé en une équation algébrique, puis de résoudre le système obtenu (linéaire, ensembliste ou bilinéaire) sans calculatrice.

Ce chapitre expose la méthode de décodage en 5 étapes et analyse en détail \textbf{10 problèmes complets de Question~3 posés entre 2022 et 2026}.

% ----------------------------------------------------------------------------
\section{Méthodologie en 5 étapes pour décoder un problème textuel}
% ----------------------------------------------------------------------------

Face à un long texte, l'erreur la plus fréquente consiste à essayer de tout garder en tête ou à écrire des calculs numériques épars sans définir ses variables. Appliquez toujours ce protocole d'ingénieur :

\begin{theoreme}[Protocole en 5 étapes pour réussir la Question 3]
\begin{enumerate}
    \item \textbf{Étape 1 -- Définition explicite des inconnues (et de leurs unités !) :}
    Dès les premières lignes de votre copie, déclarez chaque variable par une phrase claire en précisant son \textbf{unité} (kg, €, m, $\text{m}^3$, nombre entier de sacs ou de sièges).
    \emph{Conseil :} Choisissez des lettres parlantes ($F, B, V$ pour les ballons de \textbf{F}oot, \textbf{B}asket, \textbf{V}olley ; $p, d, b$ pour \textbf{p}lat, \textbf{d}essert, \textbf{b}oisson) plutôt que $x_1, x_2, x_3, x_4$ : vous éviterez ainsi de confondre vos variables en fin de calcul !
    \item \textbf{Étape 2 -- Traduction phrase par phrase en équations numérotées $(1), (2), (3)\dots$ :}
    Surlignez chaque donnée chiffrée ou relation de proportionnalité dans le texte (« trois fois plus que », « un sixième de moins », « la moyenne de », « un nombre impair dont la majorité est $\frac{N+1}{2}$ ») et écrivez l'équation correspondante.
    \begin{itemize}
        \item \textbf{Homogénéité des unités :} Vérifiez que les deux membres de chaque équation sont exprimés dans la \textbf{même unité} (ne mélangez pas des grammes avec des kilogrammes, ou des centilitres avec des litres : convertissez $30\text{ kg} = 30\,000\text{ g}$ et $94 \times 25\text{ cL} = 2350\text{ cL}$ !).
    \end{itemize}
    \item \textbf{Étape 3 -- Simplification immédiate et choix de la variable pivot :}
    Avant de vous lancer dans de grandes substitutions, divisez chaque équation par son plus grand facteur commun (ex. $4p + 4d + 4b = 120 \implies p + d + b = 30$). Repérez si plusieurs inconnues s'expriment simplement en fonction d'une seule \textbf{variable pivot} (par exemple $E = 3C$, $D = 2B$, ou $B = 2V$ et $F = \frac{3}{2}V$).
    \item \textbf{Étape 4 -- Résolution par substitution en cascade ou combinaisons linéaires :}
    \begin{itemize}
        \item Pour un \textbf{système linéaire} ($ax + by + cz = d$), procédez par \textbf{combinaisons linéaires} (soustraire deux équations ayant des termes identiques comme $(1) - (2)$ pour éliminer $c_0$ ou $p + d$) ou par substitution de la variable pivot.
        \item Pour un \textbf{système bilinéaire} (comportant des produits d'inconnues comme $N \cdot p = 3240$ et $(N + 5)p = 3510$), développez $(N + 5)p = Np + 5p$ et remplacez le bloc $Np$ tout entier par sa valeur $3240$ !
    \end{itemize}
    \item \textbf{Étape 5 -- Contrôle final dans l'énoncé et report dans les cadres :}
    En 30 secondes, vérifiez que vos valeurs numériques redonnent bien les totaux de l'énoncé, puis inscrivez lisiblement vos réponses (avec unités) dans les cadres prévus.
\end{enumerate}
\end{theoreme}

% ----------------------------------------------------------------------------
\section{Systèmes linéaires classiques et variable pivot (2022, 2023, 2024, 2025, 2026)}
% ----------------------------------------------------------------------------

\subsection{Élimination rapide par soustraction d'équations}

\begin{exemple}[L'addition au restaurant à la Mer du Nord -- UCLouvain, Septembre 2022, Q3]
\textbf{Résumé de l'énoncé :} Au restaurant, tous les plats principaux sont au même prix $p$, tous les desserts au prix $d$ et toutes les boissons au prix $b$.
\begin{itemize}
    \item Une affiche annonce : « Formule familiale : plat + boisson + dessert pour $4$ personnes à $100$\,€ au lieu de $120$\,€ ! » (le prix normal pour $4$ formules complètes est donc de $120$\,€).
    \item Aude a pris $1$ plat et $2$ boissons. Laurent a pris $1$ plat, $1$ boisson et $1$ dessert. Charles a pris autant de boissons qu'Aude et Laurent ensemble (soit $2 + 1 = 3$ boissons) et $1$ dessert, sans plat.
    \item Charles a réglé l'addition totale de $80$\,€ pour les trois amis.
    \item Aude a remarqué : « Si on prend tous les trois un dessert, cela revient aussi cher que ton plat et le mien ! »
\end{itemize}
Combien chacun doit-il payer (et combien Aude et Laurent doivent-ils rembourser à Charles) ?

\textbf{Résolution structurée :}
\begin{enumerate}
    \item \textbf{Inconnues :} $p$ = prix d'un plat (€), $d$ = prix d'un dessert (€), $b$ = prix d'une boisson (€).\\
    Parts respectives : Laurent $= p + b + d$, Aude $= p + 2b$, Charles $= 3b + d$.
    \item \textbf{Mise en équations :}
    \begin{align}
        \text{Prix normal (4 formules) :}\quad 4p + 4d + 4b = 120 &\iff p + d + b = 30 \label{eq:rest1}\\
        \text{Facture totale (3 amis) :}\quad 2p + 2d + 6b = 80 &\iff p + d + 3b = 40 \label{eq:rest2}\\
        \text{Remarque d'Aude (3 desserts = 2 plats) :}\quad 3d &= 2p \label{eq:rest3}
    \end{align}
    \item \textbf{Résolution éclair :}
    \begin{itemize}
        \item Notons d'emblée que l'équation~\eqref{eq:rest1} $p + d + b = 30$ donne \textbf{directement la part de Laurent : $30$\,€} !
        \item En soustrayant \eqref{eq:rest2} $-$ \eqref{eq:rest1} : $(p + d + 3b) - (p + d + b) = 40 - 30 \iff 2b = 10 \iff \mathbf{b = 5\,\text{€}}$.
        \item On remplace $b = 5$ dans \eqref{eq:rest1} : $p + d = 25$. Avec \eqref{eq:rest3} ($p = \frac{3}{2}d$), il vient $\frac{5}{2}d = 25 \iff \mathbf{d = 10\,\text{€}}$ et $\mathbf{p = 15\,\text{€}}$.
    \end{itemize}
    \item \textbf{Conclusion :} Laurent doit payer $p + b + d = 15 + 5 + 10 = \boxed{30\,\text{€}}$, Aude doit payer $p + 2b = 15 + 10 = \boxed{25\,\text{€}}$, et Charles a consommé pour $3b + d = 15 + 10 = \boxed{25\,\text{€}}$ (ayant payé les $80$\,€, il doit recevoir $30$\,€ de Laurent et $25$\,€ d'Aude).
\end{enumerate}
\end{exemple}

\begin{exemple}[Logistique de brasserie : bouteilles, casiers, fûts et palettes -- UCLouvain, Juillet 2023 Série 1, Q3]
\textbf{Résumé de l'énoncé :} Chez \emph{Eaux, Pils et Liqueurs (EPL)}, la bière ($1\text{ L} = 1\text{ kg}$) est vendue en casiers de $24$ bouteilles de $25\text{ cL}$ (soit $24 \times 0,25 = 6\text{ L} = 6\text{ kg}$ de bière par casier plein) ou en fûts de $50\text{ L}$ (soit $50\text{ kg}$ de bière par fût plein).
\begin{itemize}
    \item Charles rapporte $6$ casiers de bouteilles vides et repart avec $8$ casiers pleins : sa voiture pèse $66\text{ kg}$ de plus à la sortie.
    \item Laurent se fait livrer $1$ palette de $8$ fûts pleins et $2$ palettes de $60$ casiers pleins chacune : le poids total (incluant les $3$ palettes) est de $2\,370\text{ kg}$.
    \item Le poids d'un fût plein équivaut à celui de $4$ casiers pleins, et le poids d'un fût vide équivaut à celui de $40$ bouteilles vides.
\end{itemize}
Déterminer le poids d'une bouteille vide ($b$), d'un casier sans bouteilles ($c$), d'un fût vide ($f$) et d'une palette seule ($p$).

\textbf{Résolution par variables agrégées (l'astuce qui simplifie tout !) :}
Plutôt que d'écrire tout de suite des équations compliquées en $b, c, f, p$, introduisons d'abord les poids intermédiaires d'un \textbf{casier de vidanges} $C_{\text{vide}} = c + 24b$, d'un \textbf{casier plein} $C_{\text{plein}} = C_{\text{vide}} + 6$ (car il contient $6\text{ kg}$ de bière) et d'un \textbf{fût plein} $F_{\text{plein}} = f + 50$ :
\begin{enumerate}
    \item \textbf{Voiture de Charles :} $8\,C_{\text{plein}} - 6\,C_{\text{vide}} = 66 \iff 8(C_{\text{vide}} + 6) - 6\,C_{\text{vide}} = 66 \iff 2\,C_{\text{vide}} + 48 = 66 \iff \mathbf{C_{\text{vide}} = 9\text{ kg}}$, et donc $\mathbf{C_{\text{plein}} = 9 + 6 = 15\text{ kg}}$.
    \item \textbf{Fût plein et fût vide :} $F_{\text{plein}} = 4\,C_{\text{plein}} = 4 \times 15 = \mathbf{60\text{ kg}}$. Comme il contient $50\text{ kg}$ de bière, le fût vide pèse $\mathbf{f = 60 - 50 = 10\text{ kg}}$.
    \item \textbf{Bouteille vide et casier vide :} $f = 40b \iff 40b = 10 \iff \mathbf{b = 0,25\text{ kg}}$. Puis de $C_{\text{vide}} = c + 24b = 9$, on tire $c = 9 - 24(0,25) = 9 - 6 = \mathbf{3\text{ kg}}$.
    \item \textbf{Camion de Laurent ($3$ palettes + $8$ fûts pleins + $120$ casiers pleins) :}
    \[
        3p + 8\,F_{\text{plein}} + 120\,C_{\text{plein}} = 2370 \iff 3p + 480 + 1800 = 2370 \iff 3p = 90 \iff \mathbf{p = 30\text{ kg}}.
    \]
\end{enumerate}
Réponses : $\boxed{b = 0,25\text{ kg},\quad c = 3\text{ kg},\quad f = 10\text{ kg},\quad p = 30\text{ kg}}$.
\end{exemple}

\subsection{Utilisation d'une variable pivot dans les systèmes à 4, 5 ou 6 inconnues}

\begin{exemple}[Les coalitions électorales à l'EPL -- UCLouvain, Juillet 2024, Q3]
\textbf{Résumé de l'énoncé :} Cinq partis ($A$ = Arch, $B$ = Bio, $C$ = Chimie, $D$ = Data, $E$ = EPL) se partagent un nombre total impair $N = A + B + C + D + E$ de sièges. Pour avoir la majorité absolue, il faut au moins $M = \frac{N + 1}{2}$ sièges (le plus petit entier strictement supérieur à $N/2$).
\begin{itemize}
    \item La coalition Data + EPL dispose de $6$ sièges de plus que la majorité : $D + E = M + 6$.
    \item La coalition Bio + Chimie + Data est tout juste à la majorité : $B + C + D = M$.
    \item EPL a trois fois plus de sièges que Chimie ($E = 3C$), Data en a deux fois plus que Bio ($D = 2B$), et Bio + Data ensemble représentent les $\frac{2}{5}$ du total des sièges ($B + D = \frac{2}{5}N$).
\end{itemize}
Combien de sièges chaque parti a-t-il obtenus ?

\textbf{Résolution par variable pivot (exprimer tout en fonction de $B$) :}
\begin{enumerate}
    \item \textbf{Exprimer $D$ et $N$ puis $M$ en fonction de $B$ :}
    Puisque $D = 2B$, on a $B + D = 3B$. Or $B + D = \frac{2}{5}N$, donc :
    \[
        3B = \frac{2}{5}N \iff \mathbf{N = \frac{15}{2}B} \qquad \text{et} \qquad \mathbf{M = \frac{N + 1}{2} = \frac{15B + 2}{4}}.
    \]
    \item \textbf{Trouver $C$ et $E$ en soustrayant les deux équations de coalition :}
    On sait que $D + E = M + 6$ et $B + C + D = M$. En soustrayant la seconde à la première ($D$ et $M$ s'éliminent !) :
    \[
        (D + E) - (B + C + D) = (M + 6) - M \iff E - B - C = 6.
    \]
    Comme $E = 3C$, cela donne $3C - B - C = 6 \iff 2C = B + 6 \iff \mathbf{C = \frac{B + 6}{2}}$.
    \item \textbf{Calculer $B$ via l'équation $B + C + D = M$ :}
    Remplaçons $D = 2B$, $C = \frac{B+6}{2}$ et $M = \frac{15B+2}{4}$ dans $B + D + C = M$ :
    \[
        3B + \frac{B + 6}{2} = \frac{15B + 2}{4} \xrightarrow{\times 4} 12B + 2B + 12 = 15B + 2 \iff \mathbf{B = 10}.
    \]
    \item \textbf{Déduction en cascade de toutes les autres inconnues :}
    $D = 2B = \mathbf{20}$, $C = \frac{10 + 6}{2} = \mathbf{8}$, $E = 3C = \mathbf{24}$, $N = \frac{15}{2}(10) = \mathbf{75}$ (et $M = \frac{75+1}{2} = 38$), puis enfin $A = N - (B + C + D + E) = 75 - (10 + 8 + 20 + 24) = 75 - 62 = \mathbf{13}$.
\end{enumerate}
Réponses : $\boxed{A = 13,\quad B = 10,\quad C = 8,\quad D = 20,\quad E = 24 \quad (\text{total } N = 75,\text{ majorité } M = 38)}$.
\end{exemple}

\begin{exemple}[Trois autres problèmes de systèmes linéaires et flux (Juillet 2025, Septembre 2025, Septembre 2026)]
\begin{enumerate}
    \item \textbf{Juillet 2025, Q3 (Le stand de jus d'oranges aux Portes Ouvertes) :}
    \emph{$94$ visiteurs boivent chacun un verre de $25\text{ cL}$ (volume total $= 94 \times 25 = 2\,350\text{ cL}$), pressé à partir de $185$ oranges pesant au total $30\text{ kg} = 30\,000\text{ g}$. Il y a trois variétés : Valencia ($x$ oranges à $150\text{ g}$ et $15\text{ cL}$/fruit), Navel ($y$ oranges à $200\text{ g}$ et $10\text{ cL}$/fruit) et Sanguines ($z$ oranges à $120\text{ g}$ et $10\text{ cL}$/fruit). Chaque sac d'oranges a le même poids $P$ (en kg), ne contient qu'une variété, et Laurent a utilisé $1$ sac de Valencia de plus que de Navel.}\\
    \textbf{Étape 1 (Système $3 \times 3$ pour le nombre d'oranges $x, y, z$) :}
    \[
        \begin{cases}
            x + y + z = 185 & (1) \text{ [Nombre d'oranges]} \\
            15x + 10y + 10z = 2350 \iff 3x + 2(y + z) = 470 & (2) \text{ [Volume de jus en cL]} \\
            150x + 200y + 120z = 30\,000 \iff 15x + 20y + 12z = 3000 & (3) \text{ [Masse totale en g]}
        \end{cases}
    \]
    De $(1)$, $y + z = 185 - x$. En injectant dans $(2)$ : $3x + 2(185 - x) = 470 \iff x + 370 = 470 \iff \mathbf{x = 100\text{ Valencia}}$ !
    Il reste $y + z = 85$ et, dans $(3)$, $20y + 12z = 3000 - 1500 = 1500 \iff 5y + 3z = 375$. Avec $z = 85 - y$ : $5y + 3(85 - y) = 375 \iff 2y = 120 \iff \mathbf{y = 60\text{ Navel}}$ et $\mathbf{z = 25\text{ Sanguines}}$.\\
    \textbf{Étape 2 (Poids d'un sac $P$ et nombres de sacs) :} Les masses totales par variété sont $M_{\text{Val}} = 100 \times 0,15\text{ kg} = 15\text{ kg}$, $M_{\text{Nav}} = 60 \times 0,20\text{ kg} = 12\text{ kg}$ et $M_{\text{Sang}} = 25 \times 0,12\text{ kg} = 3\text{ kg}$.
    La relation sur le nombre de sacs $n_{\text{Val}} = n_{\text{Nav}} + 1$ s'écrit $\dfrac{15}{P} = \dfrac{12}{P} + 1 \iff \boxed{P = 3\text{ kg}}$, d'où $\boxed{5\text{ sacs Valencia},\; 4\text{ sacs Navel et } 1\text{ sac Sanguines}}$.

    \item \textbf{Septembre 2025, Q3 (Les filets de ballons de sport) :}
    \emph{Soient $F, B, V$ les masses (en g) d'un ballon de foot, basket, volley, et $S$ celle d'un filet vide.
    Relations entre ballons : $F = \frac{B + V}{2}$ (moyenne) et $2F + 3V = 3B$.
    On remplit $10$ filets ($a$ filets de type A contenant $3F + 2B + 3V$, et $b$ filets de type B contenant $3B + 5V$, avec $a + b = 10$) de sorte qu'il y ait autant de ballons de volley que de foot et basket réunis. Un filet A pèse $150\text{ g}$ de plus qu'un filet B, et les $10$ filets remplis pèsent $34\,750\text{ g}$.}\\
    \textbf{Résolution :}
    \begin{itemize}
        \item En remplaçant $2F = B + V$ dans $2F + 3V = 3B$, on obtient $B + 4V = 3B \iff \mathbf{B = 2V}$, puis $\mathbf{F = \frac{3}{2}V}$ ($V$ est la variable pivot !).
        \item L'écart de poids entre un filet A et un filet B donne : $(3F + 2B + 3V + S) - (3B + 5V + S) = 150 \iff 3F - B - 2V = 150$. En remplaçant en fonction de $V$ : $\frac{9}{2}V - 2V - 2V = \frac{1}{2}V = 150 \implies \boxed{V = 300\text{ g}}$, $\boxed{B = 600\text{ g}}$ et $\boxed{F = 450\text{ g}}$.
        \item Nombre de filets $a$ et $b$ ($a + b = 10$) : la condition $V_{\text{tot}} = F_{\text{tot}} + B_{\text{tot}}$ s'écrit $3a + 5b = 3a + (2a + 3b) \iff 2b = 2a \iff \mathbf{a = b = 5}$.
        \item Masse d'un filet vide $S$ : un filet A (sans filet) pèse $3(450) + 2(600) + 3(300) = 3450\text{ g}$ et un filet B pèse $3(600) + 5(300) = 3300\text{ g}$. Pour $5$ filets A et $5$ filets B : $5(3450 + 3300) + 10S = 33\,750 + 10S = 34\,750 \implies 10S = 1000 \implies \boxed{S = 100\text{ g}}$.
    \end{itemize}

    \item \textbf{Septembre 2026, Q3 (L'incendie du bois de Lauzelle : débits et rotations) :}
    \emph{$1\,060\text{ m}^3$ d'eau ont été prélevés en $149$ passages par des camions de pompiers (capacité $C$), des citernes agricoles (capacité $A$) et l'hélicoptère Charlie (capacité $H$). Le jour, les camions ont fait $3$ fois plus de passages que les citernes ($3a$ contre $a$). La nuit, seuls les camions ont fait $16$ passages. Charlie a fait $5$ passages au total. Le débit d'arrosage d'une citerne est le double de celui d'un camion ($2q$ contre $q$), mais elle se vide $3$ fois plus vite ($t/3$ contre $t$). Enfin, les $5$ passages de Charlie ont emporté autant d'eau que $8$ passages de camions.}\\
    \textbf{Résolution en 2 temps :}
    \begin{itemize}
        \item \emph{1. Nombre de passages :} $a + 3a + 16 + 5 = 149 \iff 4a = 128 \iff \mathbf{a = 32}$. Les citernes ont donc effectué $\boxed{32\text{ passages}}$ et les camions $3(32) + 16 = \boxed{112\text{ passages}}$.
        \item \emph{2. Capacités (avec $C = q\cdot t$ comme variable pivot) :}
        La capacité d'une citerne vaut $A = \text{débit} \times \text{temps} = (2q)\cdot\frac{t}{3} = \frac{2}{3}qt = \mathbf{\frac{2}{3}C}$.
        Pour l'hélicoptère : $5H = 8C \iff \mathbf{H = \frac{8}{5}C}$.
        L'équation du volume total $112C + 32A + 5H = 1060$ devient :
        \[
            112C + \frac{64}{3}C + 8C = 1060 \iff \frac{424}{3}C = 1060 \implies C = \frac{3180}{424} = \boxed{7,5\text{ m}^3},
        \]
        d'où $A = \frac{2}{3}(7,5) = \boxed{5\text{ m}^3}$ et $H = \frac{8}{5}(7,5) = \boxed{12\text{ m}^3}$.
    \end{itemize}
\end{enumerate}
\end{exemple}

% ----------------------------------------------------------------------------
\section{Problèmes ensemblistes (Diagramme de Venn), bilinéaires et suites}
% ----------------------------------------------------------------------------

\subsection{Problème de dénombrement par diagramme de Venn à 3 ensembles}

Lorsqu'un problème répartit une population de $N$ individus selon leur appartenance à trois catégories qui se chevauchent (ex. les fans photographiés avec les trois membres d'un groupe de rock), \textbf{dessinez immédiatement un diagramme de Venn à 3 cercles} et inscrivez une variable dans chacune des \textbf{7 régions disjointes} !

\begin{exemple}[Les autographes du groupe de rock \emph{The EPLighters} -- UCLouvain, Juillet 2023 Série 2, Q3]
\textbf{Résumé de l'énoncé :} Chaque fan ($N$ fans au total) repart avec exactement une photo prise avec une ou plusieurs des trois stars du groupe : Charles ($C$), Aude ($A$) et Laurent ($L$).
\begin{itemize}
    \item Charles signe au total $200$ photos, soit deux fois plus qu'Aude ($100$ photos au total pour Aude).
    \item Si Laurent avait signé $6$ fois plus de photos au total, il lui en aurait encore manqué $10$ pour couvrir tous les fans : $N = 6\,\text{Tot}_L + 10$. De plus, Laurent apparaît \textbf{seul} sur $35$ clichés.
    \item Le nombre de photos avec Charles et Aude (sans Laurent) est le double du nombre de photos avec Charles et Laurent (sans Aude). Aude et Laurent n'apparaissent \textbf{jamais} ensemble sans Charles ($0$ photo). Enfin, $5$ photos réunissent le groupe au complet ($C, A$ et $L$).
\end{itemize}
Combien y avait-il de fans $N$ à cette soirée ?

\begin{center}
\begin{tikzpicture}[scale=0.95, >=Stealth]
    % Cercle Charles
    \draw[very thick, uclblue, fill=ucllightblue, fill opacity=0.35] (-1.2, 0.7) circle (1.8);
    % Cercle Aude
    \draw[very thick, defgreen, fill=deflightgreen, fill opacity=0.35] (1.2, 0.7) circle (1.8);
    % Cercle Laurent
    \draw[very thick, exorange, fill=exlightorange, fill opacity=0.35] (0, -1.2) circle (1.8);

    \node[uclblue, font=\bfseries] at (-2.7, 2.3) {Charles ($\text{Tot}_C = 200$)};
    \node[defgreen, font=\bfseries] at (2.7, 2.3) {Aude ($\text{Tot}_A = 100$)};
    \node[exorange, font=\bfseries] at (0, -3.35) {Laurent ($\text{Tot}_L = 40 + x$)};

    % Les 7 régions disjointes
    \node[font=\bfseries] at (-1.8, 1.1) {$c = 165$};
    \node[font=\bfseries] at (1.8, 1.1) {$a = 75$};
    \node[font=\bfseries] at (0, -1.8) {$35$};
    \node[uclblue, font=\bfseries] at (0, 1.35) {$2x = 20$};
    \node[exorange, font=\bfseries] at (-0.95, -0.35) {$x = 10$};
    \node[defgreen, font=\bfseries] at (0.95, -0.35) {$0$};
    \node[warnred, font=\bfseries] at (0, 0.25) {$5$};
\end{tikzpicture}
\end{center}

\textbf{Résolution avec le diagramme de Venn :}
Notons $x = CL$ le nombre de photos avec Charles et Laurent seuls (donc $2x$ avec Charles et Aude seuls), $c$ le nombre de photos avec Charles seul et $a$ avec Aude seule. Les $7$ régions disjointes contiennent donc $c,\; a,\; 35,\; 2x,\; x,\; 0,\; 5$ :
\begin{enumerate}
    \item \textbf{Total du cercle de Charles :} $c + 2x + x + 5 = 200 \iff \mathbf{c = 195 - 3x}$.
    \item \textbf{Total du cercle d'Aude :} $a + 2x + 0 + 5 = 100 \iff \mathbf{a = 95 - 2x}$.
    \item \textbf{Total du cercle de Laurent :} $\text{Tot}_L = 35 + x + 0 + 5 = \mathbf{40 + x}$.
    \item \textbf{Nombre total de fans $N$ (somme des 7 régions disjointes) :}
    \[
        N = c + a + 35 + 2x + x + 0 + 5 = (195 - 3x) + (95 - 2x) + 3x + 40 = \mathbf{330 - 2x}.
    \]
    \item \textbf{Équation sur le total de Laurent :} $N = 6\,\text{Tot}_L + 10 = 6(40 + x) + 10 = \mathbf{250 + 6x}$.
    \item \textbf{Égalisation des deux expressions de $N$ :}
    \[
        330 - 2x = 250 + 6x \iff 8x = 80 \iff \mathbf{x = 10} \implies N = 330 - 2(10) = \boxed{310\text{ fans}}.
    \]
\end{enumerate}
\end{exemple}

\subsection{Systèmes bilinéaires (Quantité $\times$ Rendement ou Longueur $\times$ Prix unitaire)}

\begin{exemple}[Deux problèmes bilinéaires avec astuce de bloc (Septembre 2024 \& Juillet 2026)]
\begin{enumerate}
    \item \textbf{Septembre 2024, Q3 (Charles l'électricien : voie rapide vs voie algébrique) :}
    \emph{Charles utilise des rouleaux de $50\text{ m}$ de fil (rouge, bleu, noir, tous au même prix $x$\,€/m). Chaque circuit utilise deux fils de couleurs différentes de même longueur.
    Chez Aude (facture $2\,000$\,€) : uniquement des circuits rouge-bleu (longueur $y$ de rouge, donc $y$ de bleu).
    Chez Laurent (facture $6\,000$\,€) : circuits rouge-bleu et noir-bleu, avec $2$ fois plus de rouge que chez Aude ($2y$) et autant de noir que de rouge chez Aude ($y$), donc la longueur de bleu vaut $2y + y = 3y$.
    Chez Sophie (facture $14\,000$\,€) : circuits rouge-noir et bleu-noir, avec autant de bleu que chez Laurent ($3y$, donc total de fil chez Sophie $= 2z + 6y$ où $z$ est la longueur de rouge chez Sophie).
    Au total, Charles a utilisé $44$ rouleaux toutes couleurs confondues, dont $14$ rouleaux de fil rouge.}\\
    \textbf{La « voie royale » en 3 lignes :}
    \begin{itemize}
        \item Les $44$ rouleaux ($44 \times 50 = 2\,200\text{ m}$ de câble au total) ont été facturés au total $2\,000 + 6\,000 + 14\,000 = 22\,000$\,€. Le prix d'un mètre de câble est donc immédiatement :
        \[
            x = \frac{22\,000\,\text{€}}{2\,200\text{ m}} = \boxed{10\,\text{€/m}} \quad (\text{soit } 500\,\text{€ le rouleau}).
        \]
        \item Chez Aude ($2\,000$\,€, soit $200\text{ m}$ de fil moitié rouge, moitié bleu), la longueur de fil rouge vaut donc $\boxed{y = 100\text{ m}}$ ($2$ rouleaux).
        \item Laurent a deux fois plus de rouge qu'Aude, soit $200\text{ m}$ ($4$ rouleaux). Comme $14$ rouleaux rouges ($700\text{ m}$) ont été utilisés au total, il reste pour Sophie $700 - 100 - 200 = \boxed{z = 400\text{ m}}$ ($8$ rouleaux rouges) !
    \end{itemize}

    \item \textbf{Juillet 2026, Q3 (Les ruches de Charles l'apiculteur : système bilinéaire en $N \cdot p$) :}
    \emph{Soient $N$ le nombre initial de ruches, $p$ le nombre de pots produits par une ruche saine, $v$ la perte par ruche due au varroa et $f$ la perte par ruche due aux frelons asiatiques.
    Année 1 : la moitié des ruches ($\frac{N}{2}$) a le varroa $\implies$ perte totale de $180$ pots.
    Année 2 : toutes les ruches subissent les frelons $\implies$ perte d'un sixième de la production normale $Np$.
    Année 3 : $\frac{1}{4}$ des ruches est éliminé (il reste $\frac{3}{4}N$ ruches), toutes touchées par le varroa ET les frelons (production $p - v - f$ par ruche) $\implies$ récolte de $1\,755$ pots.
    Année 4 : exploitation reconstituée avec $5$ ruches de plus ($N + 5$ ruches toutes saines) $\implies$ récolte doublée par rapport à l'année 3 ($2 \times 1755 = 3\,510$ pots).}\\
    \textbf{Mise en équations et résolution par le bloc $Np$ :}
    \begin{itemize}
        \item Année 1 : $\dfrac{N}{2} v = 180 \iff \mathbf{Nv = 360}$.
        \item Année 2 : $Nf = \dfrac{1}{6}Np \iff \mathbf{f = \frac{p}{6}}$ (une ruche touchée par les frelons perd $\frac{1}{6}$ de sa production, il lui reste donc $\frac{5}{6}p$).
        \item Année 3 : $\dfrac{3}{4}N(p - v - f) = 1755 \iff \dfrac{3}{4}N\left(\frac{5}{6}p - v\right) = 1755$. En développant le produit par $N$ pour faire apparaître les blocs connus $\mathbf{Np}$ et $\mathbf{Nv = 360}$ :
        \[
            \frac{5}{8}(Np) - \frac{3}{4}(360) = 1755 \iff \frac{5}{8}(Np) = 1755 + 270 = 2025 \implies \mathbf{Np = 3240}.
        \]
        \item Année 4 : $(N + 5)p = 3510 \iff Np + 5p = 3510$. En remplaçant le bloc $Np = 3240$ :
        \[
            3240 + 5p = 3510 \iff 5p = 270 \iff \boxed{p = 54\text{ pots/ruche saine}}.
        \]
        \item On déduit aussitôt : $N = \dfrac{3240}{54} = \boxed{60\text{ ruches}}$, puis $v = \dfrac{360}{60} = 6$ (ruche avec varroa : $p - v = 54 - 6 = \boxed{48\text{ pots}}$) et $f = \dfrac{54}{6} = 9$ (ruche avec frelons : $p - f = 54 - 9 = \boxed{45\text{ pots}}$).
    \end{itemize}
\end{enumerate}
\end{exemple}

\subsection{Problèmes combinant systèmes d'équations et suites (Juillet 2022 \& Septembre 2023)}

\begin{exemple}[Deux problèmes mêlant suites et systèmes (Juillet 2022 \& Septembre 2023)]
\begin{enumerate}
    \item \textbf{Juillet 2022, Q3 (Les panneaux photovoltaïques de Charles et Laurent) :}
    \emph{Un panneau neuf produit $P = 300\text{ kWh/an}$. Les panneaux polycristallins de Charles (coût unitaire $c_p$) perdent $5\,\%$ de production par an (raison géométrique $r = 0,95$). Charles installe le plus petit nombre entier $S$ de panneaux pour que la production de la $10$\ieme{} année ($P_{10} = S \cdot 300 \cdot 0,95^9$, avec $300 \times 0,95^9 = 189,0\text{ kWh}$) couvre sa consommation $W = 2600\text{ kWh}$. Sur $10$ ans, au prix $c_e$\,€/kWh et avec un coût fixe d'installation $c_0$, Charles réalise un bénéfice net $c_C = 1560$\,€ à l'achat, ou $c_R = 300$\,€ en leasing (loyer annuel de $12\,\%$ du prix des panneaux pendant $10$ ans, soit $10 \times 0,12 = 1,2$ fois le prix d'achat des panneaux). Laurent installe $S$ panneaux monocristallins (sans dégradation, $30\,\%$ plus chers : $c_m = 1,3\,c_p$) et réalise un bénéfice net $c_L = 3390$\,€ sur $10$ ans.}\\
    \textbf{Résolution :}
    \begin{itemize}
        \item \emph{Nombre de panneaux $S$ :} À la $10$\ieme{} année, $189 \cdot S \ge 2600 \iff S \ge \frac{2600}{189} \approx 13,75 \implies \boxed{S = 14\text{ panneaux}}$.
        \item \emph{Productions sur $10$ ans :} Pour Laurent (sans dégradation), $W_L = 10 \times 14 \times 300 = \mathbf{42\,000\text{ kWh}}$. Pour Charles (suite géométrique), $W_C = 14 \times 300 \times \frac{1 - 0,95^{10}}{1 - 0,95} = \mathbf{33\,706\text{ kWh}}$.
        \item \emph{Système linéaire pour $c_p, c_m, c_e, c_0$ :}
        \[
            \begin{cases}
                c_C = W_C c_e - c_0 - S c_p = 1560 & (1) \\
                c_R = W_C c_e - c_0 - 1,2\,S c_p = 300 & (2) \\
                c_L = W_L c_e - c_0 - 1,3\,S c_p = 3390 & (3)
            \end{cases}
        \]
        En soustrayant $(1) - (2)$, $W_C c_e$ et $c_0$ s'éliminent d'un coup : $0,2\,S c_p = 1560 - 300 = 1260 \implies c_p = \frac{1260}{2,8} = \boxed{450\,\text{€}}$, d'où $c_m = 1,3 \times 450 = \boxed{585\,\text{€}}$.
        Puis $(3) - (1)$ élimine $c_0$ : $(W_L - W_C)c_e - 0,3\,S c_p = 3390 - 1560 = 1830 \iff 8294\,c_e = 1830 + 1890 = 3720 \implies c_e = \boxed{0,45\,\text{€/kWh}}$, et enfin $(1)$ donne $c_0 = \boxed{7\,257\,\text{€}}$.
    \end{itemize}

    \item \textbf{Septembre 2023, Q3 (Rénovation du living et escalier suspendu) :}
    \emph{Pièce rectangulaire de dimensions $x$ et $y$ (hauteur sous plafond $h$) avec $2$ portes ($2\text{ m} \times 1\text{ m}$ chacune) et $1$ fenêtre ($1,6\text{ m}$ de haut sur toute la longueur $y$ de la pièce).
    Parquet au sol : surface $xy = 36\text{ m}^2$.
    Plinthes au bas des 4 murs (sauf les 2 portes de $1\text{ m}$ de large) : $2x + 2y - 2 = 23 \iff 2x + 2y = 25\text{ m}$.
    Escalier central de $16$ marches régulièrement espacées : les $15$ premières marches sont suspendues chacune par $2$ barres descendant du plafond (la $16$\ieme{} est au niveau du plafond, donc $0\text{ m}$ de barre). La longueur totale des $30$ barres est de $42\text{ m}$ (donc $21\text{ m}$ pour une rangée de $15$ barres).
    Quelle quantité de peinture (en kg) faut-il pour peindre les $4$ murs (hors portes et fenêtre) en $2$ couches avec un rendement de $4\text{ m}^2/\text{kg}$ ?}\\
    \textbf{Résolution :}
    \begin{itemize}
        \item \emph{1. Hauteur sous plafond $h$ (suite arithmétique des barres d'escalier) :} La hauteur $h$ est divisée en $16$ hauteurs de marche égales à $r = \frac{h}{16}$. Les longueurs d'une rangée de $15$ barres forment la suite arithmétique $u_1 = r,\; u_2 = 2r,\; \dots,\; u_{15} = 15r$, dont la somme vaut $\frac{42}{2} = 21\text{ m}$ :
        \[
            S_{15} = \frac{15 \times 16}{2}\,r = 120\,r = 21 \implies r = \frac{21}{120} = \frac{7}{40}\text{ m} \implies h = 16\,r = \frac{14}{5} = \mathbf{2,8\text{ m}}.
        \]
        \item \emph{2. Dimensions $x$ et $y$ de la pièce (Somme-Produit !) :} On a $x + y = \frac{25}{2} = 12,5$ et $xy = 36$. Les dimensions sont les racines de $Y^2 - \frac{25}{2}Y + 36 = 0 \iff 2Y^2 - 25Y + 72 = 0$ ($\Delta = 625 - 576 = 49 = 7^2$), ce qui donne $\mathbf{y = 8\text{ m}}$ (la longueur) et $\mathbf{x = 4,5\text{ m}}$ (la largeur).
        \emph{Astuce remarquable :} pour la surface des 4 murs $2h(x + y) = 2,8 \times 25 = 70\text{ m}^2$, on avait déjà $2(x + y) = 25\text{ m}$ sans même résoudre le second degré, mais il fallait connaître la longueur $y = 8\text{ m}$ pour calculer la surface de la fenêtre !
        \item \emph{3. Surface à peindre et masse de peinture :}
        $S_{\text{murs}} = 2h(x + y) - S_{\text{fenêtre}} - 2\,S_{\text{porte}} = 2,8 \times 25 - (1,6 \times 8) - 2(2 \times 1) = 70 - 12,8 - 4 = \mathbf{53,2\text{ m}^2} = \frac{266}{5}\text{ m}^2$.
        Pour $2$ couches (soit $2 \times 53,2 = 106,4\text{ m}^2$) à $4\text{ m}^2/\text{kg}$, il faut $\dfrac{106,4}{4} = \boxed{26,6\text{ kg de peinture}}$.
    \end{itemize}
\end{enumerate}
\end{exemple}

% ----------------------------------------------------------------------------
\section{Exercices d'entraînement du Chapitre 8}
% ----------------------------------------------------------------------------

\emph{Les corrigés détaillés de tous ces exercices se trouvent au Chapitre~10.}

\begin{enumerate}[label=\textbf{Exercice 8.\arabic* :}, leftmargin=*]
    \item \textbf{(UCLouvain -- Septembre 2021, Q3 : La collecte des déchets des trois voisins)} Trois voisins, Laurent ($L$), Nadine ($N$) et Bruno ($B$), paient ensemble une facture annuelle totale de $P = 251,60$\,€ pour la collecte de leurs déchets ménagers. Laurent et Nadine ont souscrit un petit forfait annuel à $50$\,€ (utilisant un conteneur de $5\text{ kg}$, qui couvre jusqu'à $50\text{ kg}$ de déchets par an et $12$ levées par an). Bruno a souscrit un grand forfait annuel à $80$\,€ (conteneur de $10\text{ kg}$, couvrant jusqu'à $70\text{ kg}$ de déchets par an et $12$ levées par an). Tout kilogramme au-delà du forfait est facturé $0,40$\,€/kg et toute levée au-delà des $12$ levées incluses est facturée $1,20$\,€/levée.
    Sur l'année, les trois voisins ont produit au total $290\text{ kg}$ de déchets et ont totalisé $43$ levées. Chacun ne sort son conteneur que lorsqu'il est intégralement rempli, et Laurent est le seul à ne pas avoir dépassé son forfait. Déterminer les quantités annuelles de déchets $Q_L, Q_N, Q_B$ (en kg) produites par chacun des trois voisins.
    \item \textbf{(Problème d'alliages et de proportions type Question 3)} Un atelier de métallurgie de l'EPL dispose de trois lingots d'alliages de cuivre, de zinc et d'étain :
    \begin{itemize}
        \item Lingot A : $60\,\%$ de cuivre, $30\,\%$ de zinc et $10\,\%$ d'étain ;
        \item Lingot B : $50\,\%$ de cuivre, $20\,\%$ de zinc et $30\,\%$ d'étain ;
        \item Lingot C : $40\,\%$ de cuivre, $50\,\%$ de zinc et $10\,\%$ d'étain.
    \end{itemize}
    Quelle masse (en kg) faut-il prélever et refondre de chacun des trois lingots pour obtenir $100\text{ kg}$ d'un nouvel alliage contenant exactement $52\,\%$ de cuivre, $31\,\%$ de zinc et $17\,\%$ d'étain ?
    \item \textbf{(Problème de diagramme de Venn type Question 3)} Parmi les $250$ étudiants inscrits à une journée de préparation à l'EPL, chacun suit au moins l'un des trois ateliers (Algèbre, Analyse, Trigonométrie) : $140$ suivent l'Algèbre, $130$ suivent l'Analyse et $110$ suivent la Trigonométrie. On sait que $30$ étudiants suivent les trois ateliers, que le nombre d'étudiants suivant uniquement l'Algèbre et l'Analyse est le double de ceux qui suivent uniquement l'Analyse et la Trigonométrie, et que $25$ étudiants suivent uniquement l'Algèbre et la Trigonométrie. Combien d'étudiants suivent \emph{uniquement} l'atelier d'Algèbre ?
\end{enumerate}

% ============================================================================
% CHAPITRE 9 : FORMULAIRE SYNTHÉTIQUE ET AIDE-MÉMOIRE (VERSION IMPRIMABLE 4 PAGES)
% ============================================================================
\cleardoublepage
\newgeometry{top=1.3cm,bottom=1.4cm,left=1.4cm,right=1.4cm}
\refstepcounter{chapter}
\addcontentsline{toc}{chapter}{\protect\numberline{\thechapter}Formulaire Synthétique et Aide-Mémoire (Imprimable -- 4 pages)}
\chaptermark{Formulaire Synthétique et Aide-Mémoire}

% ============================================================================
% PAGE 1/4 : POLYNÔMES, ÉQUATIONS, RADICAUX, EXP & LOG (1 PAGE A4 IMPRIMABLE)
% ============================================================================
\thispagestyle{plain}
\refstepcounter{section}
\addcontentsline{toc}{section}{\protect\numberline{\thesection}Polynômes, Équations, Radicaux, Exponentielles et Logarithmes}

\begin{center}
    {\LARGE\bfseries\color{uclblue} Formulaire d'Algèbre : Polynômes, Exp \& Log}\\[2pt]
    {\small\color{darkgray} Examen d'entrée Ingénieur Civil -- UCLouvain (Fiche récapitulative imprimable 1/4)}
    \vspace{2pt}
    \hrule height 1pt
\end{center}
\vspace{-6pt}

{\fontsize{9.5pt}{11.5pt}\selectfont
\setlength{\columnsep}{18pt}
\begin{multicols}{2}

\noindent{\large\bfseries\color{uclblue} 1. Identités Remarquables}
\begin{itemize}[leftmargin=1.2em, itemsep=1.5pt, topsep=2pt]
    \item $(a \pm b)^2 = a^2 \pm 2ab + b^2$
    \item $a^2 - b^2 = (a - b)(a + b)$
    \item $(a \pm b)^3 = a^3 \pm 3a^2b + 3ab^2 \pm b^3$
    \item $a^3 - b^3 = (a - b)(a^2 + ab + b^2)$
    \item $a^3 + b^3 = (a + b)(a^2 - ab + b^2)$
    \item $a^4 + 4b^4 = (a^2 - 2ab + 2b^2)(a^2 + 2ab + 2b^2)$
\end{itemize}

\vspace{3pt}
\noindent{\large\bfseries\color{uclblue} 2. Second Degré \& Viète ($ax^2 + bx + c = 0$)}
\begin{itemize}[leftmargin=1.2em, itemsep=1.5pt, topsep=2pt]
    \item \textbf{Discriminant :} $\Delta = b^2 - 4ac$
    \item \textbf{Racines ($\Delta \ge 0$) :} $x_{1,2} = \dfrac{-b \pm \sqrt{\Delta}}{2a}$
    \item \textbf{Factorisation :} $a(x - x_1)(x - x_2)$
    \item \textbf{Signe :} Du signe de $a$ à l'\textbf{extérieur} des racines ($x < x_1$ ou $x > x_2$), du signe de $-a$ entre les racines. Si $\Delta < 0$, toujours du signe de $a$ sur $\R$.
    \item \textbf{Somme et Produit (Viète) :}
    \[
        S = x_1 + x_2 = -\frac{b}{a}, \qquad P = x_1 x_2 = \frac{c}{a}
    \]
    Trinôme de somme $S$ et produit $P$ : $a(x^2 - Sx + P) = 0$.
\end{itemize}

\vspace{3pt}
\noindent{\large\bfseries\color{uclblue} 3. Division Euclidienne de Polynômes}
\begin{itemize}[leftmargin=1.2em, itemsep=1.5pt, topsep=2pt]
    \item \textbf{Écriture unique :} $P(x) = D(x)\cdot Q(x) + R(x)$ avec $\deg(R) < \deg(D)$.
    \item \textbf{Divisibilité :} $P$ est divisible par $D \iff R(x) = 0$.
    \item \textbf{Loi du reste (division par $x - \alpha$) :} Le reste vaut la constante $R = P(\alpha)$. En particulier :
    \[
        P(x) \text{ divisible par } (x - \alpha) \iff P(\alpha) = 0.
    \]
    \item \textbf{Divisibilité de $x^4 + ax + b$ par $x^2 + px + q$ :} Poser $x^4 + ax + b = (x^2 + px + q)(x^2 + ux + v)$, développer et identifier les coefficients de $x^3, x^2, x^1, x^0$.
\end{itemize}

\vspace{3pt}
\noindent{\large\bfseries\color{uclblue} 4. Valeurs Absolues \& Équations Irrationnelles}
\begin{itemize}[leftmargin=1.2em, itemsep=1.5pt, topsep=2pt]
    \item $\sqrt{x^2} = |x|$ ($= x$ si $x \ge 0$, $= -x$ si $x < 0$).
    \item $|A(x)| < k \;(k > 0) \iff -k < A(x) < k$ (intersection).
    \item $|A(x)| > k \;(k \ge 0) \iff A(x) < -k \text{ ou } A(x) > k$.
    \item \textbf{Équation avec racine carrée :}
    \[
        \boxed{\sqrt{A(x)} = B(x) \iff \begin{cases} B(x) \ge 0 \\ A(x) = \bigl(B(x)\bigr)^2 \end{cases}}
    \]
    \emph{Attention aux solutions parasites où $B(x) < 0$ !}
\end{itemize}

\columnbreak

\noindent{\large\bfseries\color{uclblue} 5. Conditions d'Existence (Domaines)}
\begin{itemize}[leftmargin=1.2em, itemsep=1.5pt, topsep=2pt]
    \item Fraction $\dfrac{N(x)}{D(x)}$ : $D(x) \neq 0$.
    \item Racine $\sqrt{A(x)}$ : $A(x) \ge 0$ (si en bas $\frac{1}{\sqrt{A(x)}}$ : $A(x) > 0$).
    \item Logarithme $\log_a\bigl(U(x)\bigr)$ : $U(x) > 0$ (et base $a > 0, a \neq 1$).
    \item \textbf{Cas composés fréquents à l'examen :}
    \begin{align*}
        \dom\sqrt{\ln\bigl(U(x)\bigr)} &: \ln\bigl(U(x)\bigr) \ge 0 \iff \boxed{U(x) \ge 1} \\
        \dom\frac{1}{\sqrt{\ln\bigl(U(x)\bigr)}} &: \ln\bigl(U(x)\bigr) > 0 \iff \boxed{U(x) > 1}
    \end{align*}
\end{itemize}

\vspace{3pt}
\noindent{\large\bfseries\color{uclblue} 6. Exponentielles \& Logarithmes}
\begin{itemize}[leftmargin=1.2em, itemsep=1.5pt, topsep=2pt]
    \item \textbf{Puissances ($a,b > 0$) :} $a^0 = 1$, $a^{-x} = \frac{1}{a^x}$, $a^{x+y} = a^x a^y$, $a^{x-y} = \frac{a^x}{a^y}$, $(a^x)^y = a^{xy}$, $a^{p/n} = \sqrt[n]{a^p}$.
    \item \textbf{Réciprocité ($a > 0, a \neq 1, y > 0$) :}
    \[
        y = a^x \iff x = \log_a(y), \quad \log_a(1) = 0, \quad \log_a(a) = 1
    \]
    \[
        \log_a(a^x) = x \qquad \text{et} \qquad a^{\log_a(y)} = y.
    \]
    \item \textbf{Lois logarithmiques ($u > 0, v > 0$) :}
    \begin{align*}
        \log_a(u \cdot v) &= \log_a(u) + \log_a(v) \\
        \log_a\left(\frac{u}{v}\right) &= \log_a(u) - \log_a(v) \\
        \log_a(u^p) &= p\cdot\log_a(u)
    \end{align*}
    \item \textbf{Changement de base :}
    \[
        \boxed{\log_a(u) = \frac{\ln(u)}{\ln(a)} = \frac{\log_b(u)}{\log_b(a)}}, \quad \log_a(b)\log_b(a) = 1.
    \]
    \item \textbf{Monotonie et inéquations :}
    \begin{itemize}[leftmargin=1em, itemsep=1pt]
        \item Si $\mathbf{a > 1}$ (\textbf{croissante}, ordre conservé) :\\
        $\log_a(u) > 0 \iff u > 1$ ; \; $\log_a(u) < 0 \iff 0 < u < 1$.
        \item Si $\mathbf{0 < a < 1}$ (\textbf{décroissante}, ordre inversé !) :\\
        $\log_a(u) > 0 \iff 0 < u < 1$ ; \; $\log_a(u) < 0 \iff u > 1$.
    \end{itemize}
\end{itemize}

\end{multicols}
}

\newpage
% ============================================================================
% PAGE 2/4 : NOMBRES COMPLEXES, COMBINATOIRE, BINÔME & PROBABILITÉS (1 PAGE A4)
% ============================================================================
\thispagestyle{plain}
\refstepcounter{section}
\addcontentsline{toc}{section}{\protect\numberline{\thesection}Nombres Complexes, Combinatoire, Binôme de Newton et Probabilités}

\begin{center}
    {\LARGE\bfseries\color{uclblue} Formulaire : Complexes, Combinatoire \& Probabilités}\\[2pt]
    {\small\color{darkgray} Examen d'entrée Ingénieur Civil -- UCLouvain (Fiche récapitulative imprimable 2/4)}
    \vspace{2pt}
    \hrule height 1pt
\end{center}
\vspace{-6pt}

{\fontsize{9.5pt}{11.5pt}\selectfont
\setlength{\columnsep}{18pt}
\begin{multicols}{2}

\noindent{\large\bfseries\color{uclblue} 1. Nombres Complexes dans $\C$ ($i^2 = -1$)}
\begin{itemize}[leftmargin=1.2em, itemsep=1.5pt, topsep=2pt]
    \item \textbf{Forme cartésienne :} $z = a + bi$ ($a = \Reop(z), b = \Imop(z) \in \R$).
    \item \textbf{Conjugué :} $z^* = a - bi$.
    \[
        z + z^* = 2a, \quad z - z^* = 2bi, \quad \boxed{z \cdot z^* = a^2 + b^2 = |z|^2}.
    \]
    \item \textbf{Division :} $\dfrac{a + bi}{c + di} = \dfrac{(a + bi)(c - di)}{c^2 + d^2}$.
    \item \textbf{Astuce des carrés :} $(1 + i)^2 = 2i$ et $(1 - i)^2 = -2i$.
    \item \textbf{Module et Argument (forme polaire $z = \rho e^{i\theta}$) :}
    \[
        \rho = |z| = \sqrt{a^2 + b^2}, \qquad \begin{cases} \cos\theta = a/\rho \\ \sin\theta = b/\rho \end{cases}
    \]
    \item \textbf{Produit, Quotient et De Moivre :}
    \[
        (\rho_1 e^{i\theta_1})(\rho_2 e^{i\theta_2}) = \rho_1\rho_2 e^{i(\theta_1+\theta_2)}, \quad \frac{\rho_1 e^{i\theta_1}}{\rho_2 e^{i\theta_2}} = \frac{\rho_1}{\rho_2}e^{i(\theta_1-\theta_2)}
    \]
    \[
        \boxed{z^n = \rho^n e^{in\theta} = \rho^n\bigl(\cos(n\theta) + i\sin(n\theta)\bigr)}.
    \]
    \item \textbf{Racines $n$-ièmes ($z^n = w = R e^{i\alpha}$, $n$ solutions) :}
    \[
        \boxed{z_k = \sqrt[n]{R} \cdot e^{i\left(\frac{\alpha}{n} + \frac{2k\pi}{n}\right)}, \quad k \in \{0, 1, \dots, n-1\}}.
    \]
\end{itemize}

\vspace{3pt}
\noindent{\large\bfseries\color{uclblue} 2. Analyse Combinatoire (Dénombrement)}
\begin{itemize}[leftmargin=1.2em, itemsep=1.5pt, topsep=2pt]
    \item \textbf{Permutations ($n$ objets distincts) :} $P_n = n!$ ($0! = 1$).
    \item \textbf{Arrangements avec répétition :} $n^k$ (listes de longueur $k$).
    \item \textbf{Arrangements sans répétition (ordre important) :}
    \[
        A_n^k = n(n-1)\dots(n-k+1) = \frac{n!}{(n-k)!}.
    \]
    \item \textbf{Combinaisons (ordre SANS importance) :}
    \[
        \boxed{C_n^k = \binom{n}{k} = \frac{n!}{k!\,(n-k)!} = C_n^{n-k}}.
    \]
    \item \textbf{Relation de Pascal :} $C_n^k = C_{n-1}^{k-1} + C_{n-1}^k$.
\end{itemize}

\columnbreak

\noindent{\large\bfseries\color{uclblue} 3. Formule du Binôme de Newton}
\begin{itemize}[leftmargin=1.2em, itemsep=1.5pt, topsep=2pt]
    \item \textbf{Développement :}
    \[
        \boxed{(a + b)^n = \sum_{k=0}^n \binom{n}{k} a^{n-k} b^k}.
    \]
    \item \textbf{Terme général de $\left(\alpha x^r + \dfrac{\beta}{x}\right)^n$ :}
    \[
        T_k = \binom{n}{k}\alpha^{n-k}\beta^k \cdot x^{r(n-k) - k}.
    \]
    \emph{Terme d'ordre nul (constant en $x^0$) :} résoudre $r(n-k) - k = 0$.
    \item \textbf{Somme et différence :}
    \begin{align*}
        (x+1)^n + (x-1)^n &= 2\sum_{k \text{ pair}} \binom{n}{k} x^{n-k} \\
        (x+1)^n - (x-1)^n &= 2\sum_{k \text{ impair}} \binom{n}{k} x^{n-k}
    \end{align*}
\end{itemize}

\vspace{3pt}
\noindent{\large\bfseries\color{uclblue} 4. Probabilités \& Statistiques}
\begin{itemize}[leftmargin=1.2em, itemsep=1.5pt, topsep=2pt]
    \item \textbf{Loi de Laplace (issues équiprobables) :}
    \[
        P(A) = \frac{\text{Nombre de cas favorables}}{\text{Nombre de cas possibles}}.
    \]
    \item \textbf{Complémentaire (« au moins un ») :} $P(A) = 1 - P(\bar{A})$.
    \item \textbf{Lancers de $n$ dés à $6$ faces :} $6^n$ issues ordonnées équiprobables :
    \begin{itemize}[leftmargin=1em, itemsep=1pt]
        \item Paire exacte sur $3$ dés : $\dfrac{\binom{3}{2}\cdot 6 \cdot 5}{6^3} = \dfrac{5}{12}$.
        \item Full ($3+2$) sur $5$ dés : $\dfrac{\binom{5}{2}\cdot 6 \cdot 5}{6^5} = \dfrac{25}{648}$.
    \end{itemize}
    \item \textbf{Tirage simultané sans remise (urne / cartes) :} Utiliser les combinaisons $\binom{N}{k}$ en haut et en bas.
    \item \textbf{Moyenne et Écart-type ($N = \sum n_i$) :}
    \[
        \bar{x} = \frac{1}{N}\sum n_i x_i, \qquad \sigma = \sqrt{\frac{1}{N}\sum n_i(x_i - \bar{x})^2} = \sqrt{\overline{x^2} - \bar{x}^2}.
    \]
\end{itemize}

\end{multicols}
}

\newpage
% ============================================================================
% PAGE 3/4 : SUITES ARITHMÉTIQUES & GÉOMÉTRIQUES ET DISCUSSION PARAMÉTRIQUE (1 PAGE A4)
% ============================================================================
\thispagestyle{plain}
\refstepcounter{section}
\addcontentsline{toc}{section}{\protect\numberline{\thesection}Suites Arithmétiques et Géométriques, Discussion Paramétrique}

\begin{center}
    {\LARGE\bfseries\color{uclblue} Formulaire : Suites \& Inéquations Paramétriques}\\[2pt]
    {\small\color{darkgray} Examen d'entrée Ingénieur Civil -- UCLouvain (Fiche récapitulative imprimable 3/4)}
    \vspace{2pt}
    \hrule height 1pt
\end{center}
\vspace{-6pt}

{\fontsize{9.5pt}{11.5pt}\selectfont
\setlength{\columnsep}{18pt}
\begin{multicols}{2}

\noindent{\large\bfseries\color{uclblue} 1. Suites Arithmétiques (pas $+r$)}
\begin{itemize}[leftmargin=1.2em, itemsep=1.5pt, topsep=2pt]
    \item \textbf{Récurrence :} $u_{n+1} = u_n + r$.
    \item \textbf{Terme général :} $u_n = u_1 + (n - 1)r = u_p + (n - p)r$.
    \item \textbf{Moyenne arithmétique :} $u_n = \dfrac{u_{n-1} + u_{n+1}}{2}$.
    \item \textbf{Somme des $n$ premiers termes ($S_n = u_1 + \dots + u_n$) :}
    \[
        \boxed{S_n = n \cdot \frac{u_1 + u_n}{2} = \frac{n}{2}\bigl(2u_1 + (n - 1)r\bigr)}.
    \]
    \item \textbf{Sommes remarquables :}
    \begin{align*}
        1 + 2 + 3 + \dots + n &= \frac{n(n+1)}{2} \\
        2 + 4 + 6 + \dots + 2n &= n(n+1) \\
        1 + 3 + 5 + \dots + (2n-1) &= n^2
    \end{align*}
\end{itemize}

\vspace{3pt}
\noindent{\large\bfseries\color{uclblue} 2. Suites Géométriques (facteur $\times q$)}
\begin{itemize}[leftmargin=1.2em, itemsep=1.5pt, topsep=2pt]
    \item \textbf{Récurrence :} $u_{n+1} = q \cdot u_n$ ($q \neq 0$).
    \item \textbf{Terme général :} $u_n = u_1 \cdot q^{n-1} = u_p \cdot q^{n-p}$.
    \item \textbf{Moyenne géométrique ($u_k > 0$) :} $u_n = \sqrt{u_{n-1}\cdot u_{n+1}}$.
    \item \textbf{Somme des $n$ premiers termes ($q \neq 1$) :}
    \[
        \boxed{S_n = u_1 \cdot \frac{1 - q^n}{1 - q} = u_1 \cdot \frac{q^n - 1}{q - 1}}.
    \]
    \item \textbf{Doublement ($q = 2$) :} $S_n = u_1(2^n - 1)$ ($2^7 - 1 = 127$, $2^{10} - 1 = 1023$).
    \item \textbf{Hausse / baisse de $p\,\%$ :} $q = 1 \pm \dfrac{p}{100}$.
\end{itemize}

\columnbreak

\noindent{\large\bfseries\color{uclblue} 3. Méthodologie de la Question 2 (Discussion Paramétrique)}
\begin{enumerate}[leftmargin=1.4em, itemsep=2pt, topsep=2pt]
    \item \textbf{Mise sous forme standard :} Tout passer à gauche et réduire au même dénominateur :
    \[
        \frac{N(x, m)}{D(x, m)} \ge 0 \quad (\text{ou } > 0, \le 0, < 0).
    \]
    \emph{Ne jamais multiplier par un dénominateur de signe inconnu !}
    \item \textbf{Valeurs interdites (C.E.) :} Résoudre $D(x, m) = 0$ (doubles barres $\|$ dans le tableau, toujours exclues de $S$).
    \item \textbf{Factorisation et discriminants $\Delta(m)$ :} Trouver les plages de $m$ où chaque trinôme a $0$, $1$ ou $2$ racines réelles.
    \item \textbf{Ordre des valeurs critiques :} Résoudre $x_{\text{num}}(m) = x_{\text{dén}}(m)$ pour trouver les valeurs charnières de $m$ où deux racines se croisent.
    \item \textbf{Tableaux de signes :} Dresser un tableau par plage de $m$ (attention au signe du coefficient dominant s'il dépend de $m$ !).
\end{enumerate}

\vspace{3pt}
\noindent{\large\bfseries\color{uclblue} 4. Inéquations Logarithmiques Paramétriques}
\begin{itemize}[leftmargin=1.2em, itemsep=1.5pt, topsep=2pt]
    \item \textbf{Base fixe $b > 1$ :}
    \[
        \log_b\bigl(U(x,a)\bigr) < k \iff \boxed{0 < U(x,a) < b^k}.
    \]
    \item \textbf{Base variable $B(a)$ (ex. $B(a) = a - 1$) :}
    Écrire $\dfrac{\ln U(x,a)}{\ln B(a)} > 0$.
    \begin{itemize}[leftmargin=1em, itemsep=1pt]
        \item Si $B(a) \le 0$ ou $B(a) = 1$ : $S = \varnothing$.
        \item Si $B(a) > 1$ : $\boxed{U(x,a) > 1}$.
        \item Si $0 < B(a) < 1$ : $\boxed{0 < U(x,a) < 1}$.
    \end{itemize}
\end{itemize}

\end{multicols}
}

\newpage
% ============================================================================
% PAGE 4/4 : MODÉLISATION DE PROBLÈMES TEXTUELS (QUESTION 3) (1 PAGE A4)
% ============================================================================
\thispagestyle{plain}
\refstepcounter{section}
\addcontentsline{toc}{section}{\protect\numberline{\thesection}Modélisation et Systèmes Algébriques (Question 3)}

\begin{center}
    {\LARGE\bfseries\color{uclblue} Formulaire : Modélisation \& Systèmes (Question 3)}\\[2pt]
    {\small\color{darkgray} Examen d'entrée Ingénieur Civil -- UCLouvain (Fiche récapitulative imprimable 4/4)}
    \vspace{2pt}
    \hrule height 1pt
\end{center}
\vspace{-6pt}

{\fontsize{9.5pt}{11.5pt}\selectfont
\setlength{\columnsep}{18pt}
\begin{multicols}{2}

\noindent{\large\bfseries\color{uclblue} 1. Dictionnaire de Traduction Algébrique}
\begin{itemize}[leftmargin=1.2em, itemsep=1.5pt, topsep=2pt]
    \item « $A$ a trois fois plus que $B$ » : $A = 3B$.
    \item « $A$ dépasse $B$ de $150\text{ g}$ » : $A - B = 150$.
    \item « Perte d'un sixième de la normale $P$ » : $\text{Perte} = \frac{1}{6}P$, il reste $\frac{5}{6}P$.
    \item « $F$ est la moyenne de $B$ et $V$ » : $F = \dfrac{B + V}{2} \iff 2F = B + V$.
    \item « Autant de $V$ que de $F$ et $B$ réunis » : $V_{\text{tot}} = F_{\text{tot}} + B_{\text{tot}}$.
    \item « Nombre impair $N$ de sièges, majorité absolue $M$ » :
    \[
        M = \frac{N + 1}{2} \iff N = 2M - 1.
    \]
    \item « Capacité = Débit $\times$ Temps de vidange » ($C = q \cdot t$) :
    Si le débit double ($2q$) et qu'elle se vide $3$ fois plus vite ($t/3$), alors $A = (2q)\frac{t}{3} = \frac{2}{3}C$.
\end{itemize}

\vspace{3pt}
\noindent{\large\bfseries\color{uclblue} 2. Astuces de Résolution des Systèmes}
\begin{itemize}[leftmargin=1.2em, itemsep=1.5pt, topsep=2pt]
    \item \textbf{Méthode de la variable pivot :} Exprimer toutes les inconnues en fonction d'une seule (ex. $B = 2V$, $F = \frac{3}{2}V$, ou $A = \frac{2}{3}C$, $H = \frac{8}{5}C$) avant d'injecter dans l'équation du total.
    \item \textbf{Soustraction membre à membre $(1) - (2)$ :} Idéal pour éliminer un coût fixe $c_0$ ou le poids d'un filet vide $S$ qui apparaît avec le même coefficient dans deux situations.
    \item \textbf{Substitution par bloc bilinéaire ($N\cdot p$) :}
    Si l'on connaît $Nv = 360$ et que l'on a $(N + 5)p = 3510$, développer en $Np + 5p = 3510$ et isoler d'abord le produit global $Np$ !
\end{itemize}

\columnbreak

\noindent{\large\bfseries\color{uclblue} 3. Diagramme de Venn à 3 Ensembles ($C, A, L$)}
Lorsque chaque individu appartient à au moins un des $3$ ensembles :
\begin{itemize}[leftmargin=1.2em, itemsep=1.5pt, topsep=2pt]
    \item Poser une inconnue dans chacune des \textbf{$7$ régions disjointes} :
    \begin{itemize}[leftmargin=1em, itemsep=1pt]
        \item $c, a, \ell$ : appartenant à un seul ensemble ($C$ seul, $A$ seul, $L$ seul) ;
        \item $x_{CA}, x_{CL}, x_{AL}$ : appartenant à exactement deux ensembles ;
        \item $x_{CAL}$ : appartenant aux trois ensembles au centre.
    \end{itemize}
    \item \textbf{Équations des cercles complets :}
    \begin{align*}
        \text{Tot}_C &= c + x_{CA} + x_{CL} + x_{CAL} \\
        \text{Tot}_A &= a + x_{CA} + x_{AL} + x_{CAL} \\
        \text{Tot}_L &= \ell + x_{CL} + x_{AL} + x_{CAL}
    \end{align*}
    \item \textbf{Population totale $N$ (somme des 7 régions disjointes) :}
    \[
        N = c + a + \ell + x_{CA} + x_{CL} + x_{AL} + x_{CAL}.
    \]
    \item \textbf{Formule du crible de Poincaré (Inclusion-Exclusion) :}
    \begin{align*}
        \text{card}(C \cup A \cup L) &= \text{Tot}_C + \text{Tot}_A + \text{Tot}_L \\
        &\quad - \bigl(\text{card}(C\cap A) + \text{card}(C\cap L) + \text{card}(A\cap L)\bigr) \\
        &\quad + \text{card}(C \cap A \cap L).
    \end{align*}
\end{itemize}

\vspace{3pt}
\noindent{\large\bfseries\color{uclblue} 4. Check-list de Vérification Finale}
\begin{enumerate}[leftmargin=1.4em, itemsep=1pt, topsep=2pt]
    \item Toutes les inconnues sont-elles définies avec leurs unités ?
    \item Les unités sont-elles homogènes dans chaque équation (g vs kg, cL vs L) ?
    \item Les solutions entières (nombre de sacs, de panneaux, de ruches, de sièges) sont-elles bien des entiers naturels positifs ?
    \item Les réponses finales sont-elles toutes reportées dans les cadres ?
\end{enumerate}

\end{multicols}
}

\restoregeometry

% ============================================================================
% ANNEXE A : LOGIQUE, ENSEMBLES ET MÉTHODES DE RAISONNEMENT EN ALGÈBRE
% ============================================================================
\chapter*{Annexe A : Raisonnement Logique, Ensembles et Pièges Classiques en Algèbre}
\addcontentsline{toc}{chapter}{Annexe A : Raisonnement Logique, Ensembles et Pièges Classiques en Algèbre}
\markboth{Annexe A : Logique et Raisonnement}{Annexe A : Logique et Raisonnement}

\emph{Bien que la logique formelle ne fasse pas l'objet d'une question théorique isolée à l'examen d'admission, la maîtrise des connecteurs logiques (« et », « ou », $\implies$, $\iff$), des opérations sur les intervalles ($\cap$, $\cup$, complémentaire) et des conditions d'équivalence lors de la résolution d'équations ou d'inéquations paramétriques est indispensable pour éviter de perdre des solutions ou d'introduire des solutions parasites.}

% ----------------------------------------------------------------------------
\section*{A.1 Connecteurs logiques, ensembles et lois de De Morgan}
% ----------------------------------------------------------------------------

À chaque propriété algébrique $P(x)$ portant sur une variable réelle $x \in \R$ correspond l'\textbf{ensemble des solutions} $E_P = \{x \in \R \mid P(x) \text{ est vraie}\}$. Les opérations logiques sur les propositions se traduisent directement par des opérations ensemblistes sur les intervalles de $\R$ :

\begin{theoreme}[Correspondance Logique $\longleftrightarrow$ Ensembles et Lois de De Morgan]
Soient $P(x)$ et $Q(x)$ deux conditions portant sur $x \in \R$, d'ensembles de solutions respectifs $A$ et $B$ :
\begin{enumerate}
    \item \textbf{Négation ($\text{non } P$) $\longleftrightarrow$ Complémentaire ($\R \setminus A$) :} Ensemble des $x$ qui ne vérifient pas $P(x)$.
    \item \textbf{Conjonction (« $P(x) \text{ \textbf{et} } Q(x)$ », notée par une accolade système $\begin{cases} P(x) \\ Q(x) \end{cases}$) $\longleftrightarrow$ Intersection ($A \cap B$) :}
    Ensemble des réels $x$ appartenant \textbf{simultanément} à $A$ et à $B$.
    \item \textbf{Disjonction (« $P(x) \text{ \textbf{ou} } Q(x)$ ») $\longleftrightarrow$ Union ($A \cup B$) :}
    Ensemble des réels $x$ appartenant à \textbf{au moins l'un} des deux ensembles $A$ ou $B$.
    \item \textbf{Lois de De Morgan :}
    \begin{align*}
        \boxed{\text{non}(P \text{ ou } Q) \iff (\text{non } P) \text{ et } (\text{non } Q)} \\
        \boxed{\text{non}(P \text{ et } Q) \iff (\text{non } P) \text{ ou } (\text{non } Q)}
    \end{align*}
\end{enumerate}
\end{theoreme}

\begin{attention}[Ne confondez jamais l'accolade « ET » ($\cap$) avec le « OU » ($\cup$) !]
\begin{itemize}
    \item Pour un \textbf{produit nul} : $A(x)\cdot B(x) = 0 \iff A(x) = 0 \;\text{ \textbf{ou} }\; B(x) = 0$ (on fait l'\textbf{union} $\cup$ des racines).
    \item Pour un \textbf{dénominateur non nul} (C.E. de $\frac{1}{A(x)B(x)}$) : par la loi de De Morgan,
    \[
        A(x)\cdot B(x) \neq 0 \iff A(x) \neq 0 \;\text{ \textbf{et} }\; B(x) \neq 0 \iff \begin{cases} A(x) \neq 0 \\ B(x) \neq 0 \end{cases}.
    \]
    \item Pour une \textbf{inéquation logarithmique} $\ln\bigl(U(x)\bigr) < 0$ : il faut la condition d'existence $U(x) > 0$ \textbf{ET} l'inégalité $U(x) < 1$, ce qui correspond à l'\textbf{intersection} $\{U(x) > 0\} \cap \{U(x) < 1\} = \{0 < U(x) < 1\}$.
\end{itemize}
\end{attention}

% ----------------------------------------------------------------------------
\section*{A.2 Implication ($\implies$) vs Équivalence ($\iff$) dans les équations et inéquations}
% ----------------------------------------------------------------------------

Résoudre une équation ou une inéquation $E(x)$ dans $\R$, c'est transformer $E(x)$ par une chaîne d'\textbf{équivalences logiques} ($\iff$) jusqu'à isoler $x$. Or, trois opérations algébriques très courantes ne sont \textbf{pas} des équivalences si l'on n'ajoute pas une condition de signe ou de non-nullité :

\begin{theoreme}[Les 3 règles d'équivalence qui sauvent vos points à l'examen]
\begin{enumerate}
    \item \textbf{Multiplication ou simplification par une expression $C(x, m)$ :}
    \begin{itemize}
        \item Dans une \textbf{équation} : on ne peut diviser ou simplifier par $C(x,m)$ que si $\boxed{C(x,m) \neq 0}$. Si $C(x,m)$ peut s'annuler, il faut \textbf{factoriser} ($A\cdot C = B\cdot C \iff C\cdot(A - B) = 0 \iff C = 0 \text{ ou } A = B$) ! Simplifier sauvagement par $z$ dans $(z-2)z^* = -iz$ ferait perdre la solution $z = 0$ !
        \item Dans une \textbf{inéquation} : on ne multiplie \textbf{jamais} par $C(x,m)$ si son signe dépend de $x$. Si $C(m)$ ne dépend que du paramètre $m$ (ex. $\ln(a - 1)$ en Juillet~2026 Q2), il faut \textbf{distinguer deux cas} : le cas $C(m) > 0$ (sens conservé) et le cas $C(m) < 0$ (sens inversé !).
    \end{itemize}
    \item \textbf{Élévation au carré des deux membres d'une équation :}
    L'implication $A(x) = B(x) \implies \bigl(A(x)\bigr)^2 = \bigl(B(x)\bigr)^2$ n'est une équivalence que si $A(x)$ et $B(x)$ ont le \textbf{même signe} :
    \[
        \boxed{\sqrt{A(x)} = B(x) \iff \begin{cases} B(x) \ge 0 \\ A(x) = \bigl(B(x)\bigr)^2 \end{cases}}.
    \]
    \item \textbf{Regroupement de logarithmes $\log_a\bigl(U(x)\bigr) + \log_a\bigl(V(x)\bigr) = \log_a\bigl(U(x)V(x)\bigr)$ :}
    Cette transformation élargit le domaine (car $U(x)V(x) > 0$ autorise $U(x) < 0$ et $V(x) < 0$ simultanément, ce qui est interdit dans l'équation initiale !). L'équivalence exacte est :
    \[
        \boxed{\log_a\bigl(U(x)\bigr) + \log_a\bigl(V(x)\bigr) = k \iff \begin{cases} U(x) > 0 \text{ et } V(x) > 0 \\ U(x)\cdot V(x) = a^k \end{cases}}.
    \]
\end{enumerate}
\end{theoreme}

% ============================================================================
% CHAPITRE 10 : CORRIGÉS DÉTAILLÉS DES EXERCICES DE FIN DE CHAPITRE
% ============================================================================
\chapter{Corrigés Détaillés des Exercices}

Ce chapitre rassemble les solutions complètes, méthodiques et commentées de tous les exercices d'entraînement proposés à la fin des Chapitres~1 à~8.

% ----------------------------------------------------------------------------
\section{Corrigés du Chapitre 1 (Polynômes, Factorisation et Division Euclidienne)}
% ----------------------------------------------------------------------------

\paragraph{Corrigé de l'Exercice 1.1 (UCLouvain -- Septembre 2021, Q1.6) :}
Par l'identité remarquable de la différence de deux carrés appliquée à $x^4 - 1 = (x^2)^2 - 1^2$ :
\[
    x^4 - 1 = (x^2 + 1)(x^2 - 1) + 0.
\]
Par unicité de la division euclidienne de $x^4 - 1$ par $x^2 + 1$, le quotient et le reste valent directement :
\[
    \boxed{Q(x) = x^2 - 1 \quad \text{et} \quad R(x) = 0}.
\]

\paragraph{Corrigé de l'Exercice 1.2 (UCLouvain -- Juillet 2023, Série 1, Q1.6) :}
\textbf{Méthode 1 (Division euclidienne posée) :} En divisant $x^4 + 0x^3 + 0x^2 + 0x + 4$ par $x^2 + 2x + 2$ :
\begin{itemize}
    \item $x^4 / x^2 = x^2$, et $(x^4 + 4) - x^2(x^2 + 2x + 2) = -2x^3 - 2x^2 + 4$.
    \item $-2x^3 / x^2 = -2x$, et $(-2x^3 - 2x^2 + 4) - (-2x)(x^2 + 2x + 2) = 2x^2 + 4x + 4$.
    \item $2x^2 / x^2 = 2$, et $(2x^2 + 4x + 4) - 2(x^2 + 2x + 2) = 0$.
\end{itemize}
\textbf{Méthode 2 (Identité de Sophie Germain) :} $x^4 + 4 = (x^4 + 4x^2 + 4) - 4x^2 = (x^2 + 2)^2 - (2x)^2 = (x^2 + 2x + 2)(x^2 - 2x + 2)$.
Le quotient vaut donc $\boxed{Q(x) = x^2 - 2x + 2}$ (et le reste est nul).

\paragraph{Corrigé de l'Exercice 1.3 (Division par Horner) :}
On divise $P(x) = 2x^4 - 5x^3 + 4x^2 - 3x + 7$ par $x - 2$ ($\alpha = 2$) :
\[
\begin{array}{c|cccc|c}
    & x^4 & x^3 & x^2 & x^1 & x^0 \\
    & 2 & -5 & 4 & -3 & 7 \\
    \alpha = 2 & \downarrow & 4 & -2 & 4 & 2 \\
    \hline
    & 2 & -1 & 2 & 1 & \mathbf{9 = R}
\end{array}
\]
On obtient : $\boxed{Q(x) = 2x^3 - x^2 + 2x + 1 \quad \text{et} \quad R(x) = 9}$.

\paragraph{Corrigé de l'Exercice 1.4 (Factorisation d'un polynôme de degré 3) :}
Puisque $x_1 = 3$ est racine de $P(x) = 2x^3 - 3x^2 - 11x + 6$, divisons $P(x)$ par $x - 3$ via le schéma de Horner :
\[
\begin{array}{c|ccc|c}
    & 2 & -3 & -11 & 6 \\
    \alpha = 3 & \downarrow & 6 & 9 & -6 \\
    \hline
    & 2 & 3 & -2 & \mathbf{0 = R}
\end{array}
\]
Le quotient est $Q(x) = 2x^2 + 3x - 2$. Ses racines ($\Delta = 9 + 16 = 25 = 5^2$) valent $x_2 = \frac{-3 - 5}{4} = -2$ et $x_3 = \frac{-3 + 5}{4} = \frac{1}{2}$.
Le deux autres racines sont donc $\boxed{-2 \text{ et } \dfrac{1}{2}}$, et la forme factorisée s'écrit :
\[
    \boxed{P(x) = 2(x - 3)(x + 2)\left(x - \frac{1}{2}\right) = (x - 3)(x + 2)(2x - 1)}.
\]

\paragraph{Corrigé de l'Exercice 1.5 (Divisibilité par $(x - 1)(x + 2)$) :}
Comme $1$ et $-2$ sont distincts, $P(x) = x^4 - 2x^3 + ax^2 + bx + 4$ est divisible par $(x - 1)(x + 2)$ si et seulement s'il est divisible par $(x - 1)$ \textbf{et} par $(x + 2)$, soit $P(1) = 0$ et $P(-2) = 0$ :
\[
    \begin{cases}
        P(1) = 1 - 2 + a + b + 4 = 0 \iff a + b = -3 \\
        P(-2) = 16 - 2(-8) + 4a - 2b + 4 = 0 \iff 4a - 2b = -36 \iff 2a - b = -18
    \end{cases}
\]
En additionnant les deux équations : $3a = -21 \implies \boxed{a = -7}$, puis $b = -3 - (-7) = \boxed{4}$.

\paragraph{Corrigé de l'Exercice 1.6 (Reste de la division par un produit $(x - 1)(x + 2)$) :}
Le diviseur $D(x) = (x - 1)(x + 2)$ étant de degré~$2$, le reste est au plus de degré~$1$ : posons $R(x) = ux + v$, de sorte que :
\[
    P(x) = (x - 1)(x + 2)Q(x) + ux + v.
\]
En évaluant cette égalité en $x = 1$ (où $P(1) = 5$) et en $x = -2$ (où $P(-2) = -4$), le terme contenant $Q(x)$ s'annule :
\[
    \begin{cases}
        P(1) = u + v = 5 \\
        P(-2) = -2u + v = -4
    \end{cases}
    \implies (u + v) - (-2u + v) = 5 - (-4) \iff 3u = 9 \iff u = 3 \;\text{ et }\; v = 2.
\]
Le reste cherché est donc $\boxed{R(x) = 3x + 2}$.

\paragraph{Corrigé de l'Exercice 1.7 (Factorisation de degré 6 par groupement) :}
Groupons les termes deux par deux dans $P(x) = x^6 - 9x^4 - x^2 + 9$ :
\[
    P(x) = x^4(x^2 - 9) - (x^2 - 9) = (x^2 - 9)(x^4 - 1) = \boxed{(x - 3)(x + 3)(x - 1)(x + 1)(x^2 + 1)}
\]
(le facteur $x^2 + 1$ étant irréductible dans $\R$).

% ----------------------------------------------------------------------------
\section{Corrigés du Chapitre 2 (Équations, Valeurs Absolues, Radicaux et Domaines)}
% ----------------------------------------------------------------------------

\paragraph{Corrigé de l'Exercice 2.1 (Équation $|2x - 3| = |x + 4|$) :}
Deux valeurs absolues sont égales ssi les expressions sont égales ou opposées :
\[
    |2x - 3| = |x + 4| \iff 2x - 3 = x + 4 \;\text{ ou }\; 2x - 3 = -(x + 4) \iff \boxed{S = \left\{-\frac{1}{3},\; 7\right\}}.
\]

\paragraph{Corrigé de l'Exercice 2.2 (Inéquation $|x^2 - 5x + 4| < 2$) :}
Puisque $2 > 0$, l'inéquation équivaut à la double inégalité $-2 < x^2 - 5x + 4 < 2$, soit au système :
\[
    \begin{cases}
        x^2 - 5x + 4 > -2 \iff x^2 - 5x + 6 > 0 \iff (x - 2)(x - 3) > 0 & (1) \\
        x^2 - 5x + 4 < 2 \iff x^2 - 5x + 2 < 0 & (2)
    \end{cases}
\]
\begin{itemize}
    \item La condition $(1)$ donne $x \in \,]-\infty, 2[ \,\cup\, ]3, +\infty[$.
    \item Le trinôme de $(2)$ a pour discriminant $\Delta = 25 - 8 = 17$, et est strictement négatif entre ses racines $\frac{5 - \sqrt{17}}{2} \approx 0,44$ et $\frac{5 + \sqrt{17}}{2} \approx 4,56$.
\end{itemize}
Comme $4 < \sqrt{17} < 5$, on a $\frac{5-\sqrt{17}}{2} < 2 < 3 < \frac{5+\sqrt{17}}{2}$. L'intersection de $(1)$ et $(2)$ donne :
\[
    \boxed{S = \left]\frac{5 - \sqrt{17}}{2},\, 2\right[ \cup \left]3,\, \frac{5 + \sqrt{17}}{2}\right[}.
\]

\paragraph{Corrigé de l'Exercice 2.3 (Équations irrationnelles) :}
\begin{enumerate}[label=(\alph*)]
    \item $\sqrt{2x + 3} = x \iff \begin{cases} x \ge 0 \\ 2x + 3 = x^2 \end{cases} \iff \begin{cases} x \ge 0 \\ x^2 - 2x - 3 = (x - 3)(x + 1) = 0 \end{cases}$. On rejette $x = -1 < 0$ et on garde $\boxed{x = 3}$.
    \item On isole d'abord le radical : $\sqrt{x + 4} = 2 - x \iff \begin{cases} 2 - x \ge 0 \iff x \le 2 \\ x + 4 = (2 - x)^2 = x^2 - 4x + 4 \end{cases} \iff \begin{cases} x \le 2 \\ x^2 - 5x = x(x - 5) = 0 \end{cases}$. La condition $x \le 2$ rejette la solution parasite $x = 5$ et conserve $\boxed{x = 0}$.
    \item C.E. des radicaux : $3x + 1 \ge 0$ et $x - 1 \ge 0 \iff \mathbf{x \ge 1}$. Isolons le plus gros radical avant d'élever au carré (les deux membres étant positifs pour $x \ge 1$) :
    \[
        \sqrt{3x + 1} = 2 + \sqrt{x - 1} \iff 3x + 1 = 4 + 4\sqrt{x - 1} + (x - 1) \iff x - 1 = 2\sqrt{x - 1}.
    \]
    En posant $u = \sqrt{x - 1} \ge 0$, on a $u^2 = 2u \iff u(u - 2) = 0 \iff u = 0 \text{ ou } u = 2$, soit $x - 1 = 0$ ou $x - 1 = 4$, d'où les deux solutions valides : $\boxed{x = 1 \quad \text{et} \quad x = 5}$.
\end{enumerate}

\paragraph{Corrigé de l'Exercice 2.4 (Domaines de définition) :}
\begin{enumerate}[label=(\alph*)]
    \item Pour $f_1(x) = \dfrac{1}{\sqrt{\ln(x^2 - 4x + 5)}}$, il faut $\ln(x^2 - 4x + 5) > 0 \iff x^2 - 4x + 5 > 1 \iff x^2 - 4x + 4 > 0 \iff (x - 2)^2 > 0$. Or un carré parfait est strictement positif partout sauf là où il s'annule ($x = 2$) ! D'où :
    \[
        \boxed{\dom f_1 = \R \setminus \{2\} = \,]-\infty, 2[ \,\cup\, ]2, +\infty[}.
    \]
    \item Pour $f_2(x) = \sqrt{\ln\left(\dfrac{2x + 1}{x - 1}\right)}$, la racine au numérateur impose $\ln\left(\dfrac{2x + 1}{x - 1}\right) \ge 0 \iff \dfrac{2x + 1}{x - 1} \ge 1$ (avec $x \neq 1$). En passant $1$ à gauche :
    \[
        \frac{(2x + 1) - (x - 1)}{x - 1} \ge 0 \iff \frac{x + 2}{x - 1} \ge 0 \iff \boxed{\dom f_2 = \,]-\infty, -2] \,\cup\, ]1, +\infty[}.
    \]
    \item Pour $f_3(x) = \dfrac{\sqrt{9 - x^2}}{\ln(x + 2)}$, il faut trois conditions simultanées :
    \begin{itemize}
        \item Radicande positif ou nul : $9 - x^2 \ge 0 \iff x \in [-3, 3]$.
        \item Argument du logarithme strictement positif : $x + 2 > 0 \iff x > -2$.
        \item Dénominateur non nul : $\ln(x + 2) \neq 0 \iff x + 2 \neq 1 \iff x \neq -1$.
    \end{itemize}
    L'intersection $[-3, 3] \cap \,]-2, +\infty[ \setminus \{-1\}$ donne $\boxed{\dom f_3 = \,]-2, -1[ \,\cup\, ]-1, 3]}$.
\end{enumerate}

% ----------------------------------------------------------------------------
\section{Corrigés du Chapitre 3 (Exponentielles et Logarithmes)}
% ----------------------------------------------------------------------------

\paragraph{Corrigé de l'Exercice 3.1 (Calculs logarithmiques sans calculatrice) :}
\begin{enumerate}[label=(\alph*)]
    \item Remarquons que $\log_4(9) = \log_{2^2}(3^2) = \frac{2}{2}\log_2(3) = \log_2(3)$, qui s'annule avec $-\log_2(3)$ ! Il reste :
    \[
        A = \log_2(24) - \log_2(3^{1/2}) = \log_2(8 \times 3) - \frac{1}{2}\log_2(3) = \dots
    \]
    \emph{Plus précisément :} en regroupant tout en base~$2$ : $\log_2(24) - \log_2(3) = \log_2(24/3) = \log_2(8) = 3$, et $\log_4(9) - \log_2(\sqrt{3}) = \log_2(3) - \frac{1}{2}\log_2(3) = \frac{1}{2}\log_2(3)$, soit $A = \boxed{3 + \frac{1}{2}\log_2(3)}$.
    \item Par changement de base en logarithme népérien (produit télescopique !) :
    \[
        B = \frac{\ln 5}{\ln 3} \cdot \frac{\ln 7}{\ln 5} \cdot \frac{\ln(3^4)}{\ln 7} = \frac{4\ln 3}{\ln 3} = \boxed{4}.
    \]
    \item Comme $9 = 3^2$ et $8 = 2^3$ : $C = (3^2)^{\log_3(5)} - (2^3)^{\log_2(3)} = \bigl(3^{\log_3(5)}\bigr)^2 - \bigl(2^{\log_2(3)}\bigr)^3 = 5^2 - 3^3 = 25 - 27 = \boxed{-2}$.
\end{enumerate}

\paragraph{Corrigé de l'Exercice 3.2 (Équations exponentielles) :}
\begin{enumerate}[label=(\alph*)]
    \item Écrivons $4^{x+1} = 4\cdot 4^x = 4\cdot(2^x)^2$ et posons $t = 2^x > 0$ :
    \[
        4t^2 - 9t + 2 = 0 \iff (4t - 1)(t - 2) = 0 \iff t = \frac{1}{4} = 2^{-2} \;\text{ ou }\; t = 2 = 2^1 \iff \boxed{x = -2 \;\text{ ou }\; x = 1}.
    \]
    \item Avec $3^{2-x} = \dfrac{9}{3^x}$ et $t = 3^x > 0$ : $t + \dfrac{9}{t} = 10 \iff t^2 - 10t + 9 = (t - 1)(t - 9) = 0 \iff t = 1 \text{ ou } t = 9 \iff \boxed{x = 0 \;\text{ ou }\; x = 2}$.
    \item Détachons les constantes : $2^2 \cdot 2^x \cdot \frac{3^x}{3} = 144 \iff \frac{4}{3} \cdot 6^x = 144 \iff 6^x = \frac{144 \times 3}{4} = 36 \times 3 = 108 \implies \boxed{x = \log_6(108) = 2 + \log_6(3)}$. (Et si le second membre vaut $48$, $6^x = 36 \implies x = 2$).
\end{enumerate}

\paragraph{Corrigé de l'Exercice 3.3 (Équations logarithmiques) :}
\begin{enumerate}[label=(\alph*)]
    \item \textbf{C.E. :} $x + 2 > 0$ et $x - 4 > 0 \iff \mathbf{x > 4}$. En regroupant :
    \[
        \log_3\bigl((x + 2)(x - 4)\bigr) = 3 \iff x^2 - 2x - 8 = 3^3 = 27 \iff x^2 - 2x - 35 = (x - 7)(x + 5) = 0.
    \]
    La condition $x > 4$ élimine $x = -5$ et conserve l'unique solution $\boxed{x = 7}$.
    \item \textbf{C.E. :} $x > 0$. Comme $\ln(x^4) = 4\ln x$, posons $t = \ln x \in \R$ : $t^2 - 4t + 3 = (t - 1)(t - 3) = 0 \iff \ln x = 1 \text{ ou } \ln x = 3 \iff \boxed{x = e \;\text{ ou }\; x = e^3}$.
    \item \textbf{C.E. :} $x > 0$ et $x \neq 1$. Par la formule d'inversion de la base $\log_x(2) = \dfrac{1}{\log_2(x)}$, en posant $t = \log_2(x) \in \R_0$ :
    \[
        t + \frac{1}{t} = \frac{5}{2} \iff 2t^2 - 5t + 2 = (2t - 1)(t - 2) = 0 \iff \boxed{x = \sqrt{2} \;\text{ ou }\; x = 4}.
    \]
\end{enumerate}

% ----------------------------------------------------------------------------
\section{Corrigés du Chapitre 4 (Nombres Complexes dans $\C$)}
% ----------------------------------------------------------------------------

\paragraph{Corrigé de l'Exercice 4.1 (Formes cartésienne et polaire) :}
\begin{enumerate}[label=(\alph*)]
    \item Au numérateur : $(1 + i)^8 = \bigl((1 + i)^2\bigr)^4 = (2i)^4 = 16 i^4 = 16 = 16 e^{i\cdot 0}$.\\
    Au dénominateur : $w = 1 - i\sqrt{3}$ a pour module $|w| = 2$ et argument $\arg(w) = -\frac{\pi}{3}$, donc $w^3 = 2^3 e^{i\cdot 3(-\pi/3)} = 8 e^{-i\pi} = -8$.\\
    En divisant : $z_1 = \dfrac{16}{-8} = \boxed{-2 = 2 e^{i\pi}}$ ($\text{module} = 2$, $\text{argument} = \pi$, forme cartésienne $-2 + 0i$).
    \item Au numérateur : $2 - 2i = 2\sqrt{2}\,e^{-i\frac{\pi}{4}}$. Au dénominateur : $1 + i\sqrt{3} = 2\,e^{i\frac{\pi}{3}}$. Par quotient :
    \[
        z_2 = \frac{2\sqrt{2}}{2}\,e^{i\left(-\frac{\pi}{4} - \frac{\pi}{3}\right)} = \boxed{\sqrt{2}\,e^{-i\frac{7\pi}{12}}} \quad \left(\text{module } \rho = \sqrt{2},\; \text{argument } \theta = -\frac{7\pi}{12}\right).
    \]
    En forme cartésienne : $z_2 = \dfrac{(2 - 2i)(1 - i\sqrt{3})}{1 + 3} = \dfrac{(2 - 2\sqrt{3}) - (2 + 2\sqrt{3})i}{4} = \boxed{\frac{1 - \sqrt{3}}{2} - \frac{1 + \sqrt{3}}{2}\,i}$.
\end{enumerate}

\paragraph{Corrigé de l'Exercice 4.2 (Équations avec conjugué $z^*$) :}
\begin{enumerate}[label=(\alph*)]
    \item Posons $z = a + bi$ ($a, b \in \R$) dans $2z + iz^* = 5 - 4i$ :
    \[
        2(a + bi) + i(a - bi) = 5 - 4i \iff (2a + b) + (a + 2b)i = 5 - 4i \iff \begin{cases} 2a + b = 5 \\ a + 2b = -4 \end{cases}
    \]
    En résolvant ce système linéaire $2 \times 2$ ($b = 5 - 2a \implies a + 2(5 - 2a) = -4 \iff -3a = -14$... ou par soustraction $2L_1 - L_2$ : $3a = 14 \implies a = \frac{14}{3}$ et $b = -\frac{13}{3}$), on trouve l'unique solution :
    \[
        \boxed{z = \frac{14}{3} - \frac{13}{3}i}.
    \]
    \item Posons $z = a + bi$ dans $z^2 + 2z^* + 1 = 0$ :
    \[
        (a^2 - b^2 + 2a + 1) + (2ab - 2b)i = 0 \iff \begin{cases} (a + 1)^2 - b^2 = 0 & (1) \\ 2b(a - 1) = 0 & (2) \end{cases}
    \]
    L'équation $(2)$ donne deux cas très simples :
    \begin{itemize}
        \item Si $b = 0$, alors $(1)$ donne $(a + 1)^2 = 0 \iff a = -1 \implies z_1 = -1$.
        \item Si $a = 1$, alors $(1)$ donne $(1 + 1)^2 - b^2 = 0 \iff b^2 = 4 \iff b = \pm 2 \implies z_{2,3} = 1 \pm 2i$.
    \end{itemize}
    D'où les trois solutions : $\boxed{S = \{-1,\; 1 + 2i,\; 1 - 2i\}}$.
\end{enumerate}

\paragraph{Corrigé de l'Exercice 4.3 (Racines $n$-ièmes dans $\C$) :}
\begin{enumerate}[label=(\alph*)]
    \item $z^3 = 27 = 3^3 e^{i\cdot 0} \iff z_k = 3 e^{i\frac{2k\pi}{3}}$ pour $k \in \{0, 1, 2\}$, ce qui donne en forme cartésienne :
    \[
        \boxed{S = \left\{3,\;\; -\frac{3}{2} + i\frac{3\sqrt{3}}{2},\;\; -\frac{3}{2} - i\frac{3\sqrt{3}}{2}\right\}}.
    \]
    \item $z^4 = -16 = 2^4 e^{i\pi} \iff z_k = 2 e^{i\left(\frac{\pi}{4} + k\frac{\pi}{2}\right)}$ pour $k \in \{0, 1, 2, 3\}$. Comme $\cos\frac{\pi}{4} = \sin\frac{\pi}{4} = \frac{\sqrt{2}}{2}$ :
    \[
        \boxed{S = \bigl\{\sqrt{2} + i\sqrt{2},\;\; -\sqrt{2} + i\sqrt{2},\;\; -\sqrt{2} - i\sqrt{2},\;\; \sqrt{2} - i\sqrt{2}\bigr\}}.
    \]
    \item Cherchons $z = a + bi$ tel que $z^2 = 3 + 4i$. En identifiant parties réelle, imaginaire et module ($|z|^2 = |3+4i| = \sqrt{9+16} = 5$) :
    \[
        \begin{cases} a^2 - b^2 = 3 \\ a^2 + b^2 = 5 \\ 2ab = 4 > 0 \end{cases} \implies \begin{cases} 2a^2 = 8 \iff a^2 = 4 \iff a = \pm 2 \\ 2b^2 = 2 \iff b^2 = 1 \iff b = \pm 1 \\ a \text{ et } b \text{ de même signe (car } ab = 2 > 0\text{)} \end{cases}
    \]
    D'où les deux racines carrées opposées : $\boxed{S = \{2 + i,\; -2 - i\}}$.
\end{enumerate}

% ----------------------------------------------------------------------------
\section{Corrigés du Chapitre 5 (Analyse Combinatoire, Binôme et Probabilités)}
% ----------------------------------------------------------------------------

\paragraph{Corrigé de l'Exercice 5.1 (UCLouvain -- Septembre 2021, Q1.5) :}
Le nombre total de façons de distribuer les $52$ cartes en $4$ paquets ordonnés de $13$ cartes est $\binom{52}{13}\binom{39}{13}\binom{26}{13}\binom{13}{13}$.
Pour qu'un paquet donné (parmi les $\binom{4}{1} = 4$ paquets) ne contienne que des cartes rouges, on choisit ses $13$ cartes parmi les $26$ cartes rouges ($\binom{26}{13}$ façons), puis on répartit les $39$ cartes restantes dans les $3$ autres paquets ($\binom{39}{13}\binom{26}{13}\binom{13}{13}$ façons).
\textbf{Attention au principe d'inclusion-exclusion :} Il est possible que \emph{deux} paquets simultanément soient entièrement rouges (l'un prend $13$ rouges, le second prend les $13$ rouges restantes, et les deux derniers se partagent les $26$ cartes noires : $\binom{4}{2}\binom{26}{13}\binom{13}{13}\binom{26}{13}\binom{13}{13}$ cas comptés deux fois !). En soustrayant cette intersection :
\[
    \boxed{P = \frac{\binom{4}{1}\binom{26}{13}\binom{39}{13}\binom{26}{13} - \binom{4}{2}\binom{26}{13}\binom{26}{13}}{\binom{52}{13}\binom{39}{13}\binom{26}{13}} = \frac{4\binom{26}{13}\binom{39}{13} - 6\binom{26}{13}}{\binom{52}{13}\binom{39}{13}}}.
\]

\paragraph{Corrigé de l'Exercice 5.2 (Binôme de Newton) :}
\begin{enumerate}[label=(\alph*)]
    \item Dans $(2x - 1)^8$, le terme général est $\binom{8}{k}(2x)^{8-k}(-1)^k = \binom{8}{k}2^{8-k}(-1)^k x^{8-k}$. Pour $x^5$, il faut $8 - k = 5 \iff k = 3$ :
    \[
        \text{Coeff} = \binom{8}{3}\cdot 2^5 \cdot (-1)^3 = 56 \times 32 \times (-1) = \boxed{-1792}.
    \]
    \item Dans $\left(3x^2 - \frac{1}{3x}\right)^9$, le terme général est $\binom{9}{k}(3x^2)^{9-k}\left(-\frac{1}{3}x^{-1}\right)^k = \binom{9}{k}3^{9-2k}(-1)^k x^{18-3k}$.
    Pour le terme constant ($x^0$), on résout $18 - 3k = 0 \iff k = 6$, d'où :
    \[
        T_6 = \binom{9}{6}\cdot 3^{9-12}\cdot (-1)^6 = \binom{9}{3}\cdot \frac{1}{3^3} = \frac{9 \times 8 \times 7}{3 \times 2 \times 1} \cdot \frac{1}{27} = \frac{84}{27} = \boxed{\frac{28}{9}}.
    \]
\end{enumerate}

\paragraph{Corrigé de l'Exercice 5.3 (Probabilités aux dés et aux urnes) :}
\begin{enumerate}[label=(\alph*)]
    \item Sur $4$ dés ($6^4 = 1296$ issues équiprobables) :
    \begin{itemize}
        \item Il y a $6$ carrés possibles (1-1-1-1 à 6-6-6-6), donc $P(\text{carré}) = \dfrac{6}{6^4} = \dfrac{1}{6^3} = \boxed{\dfrac{1}{216}}$.
        \item Il y a $A_6^4 = 6 \times 5 \times 4 \times 3 = 360$ tirages à $4$ chiffres tous distincts, donc $P(\text{tous différents}) = \dfrac{6 \times 5 \times 4 \times 3}{6^4} = \dfrac{20}{36} = \boxed{\dfrac{5}{9}}$.
    \end{itemize}
    \item L'urne contient $5 + 3 + 2 = 10$ boules, donc $\binom{10}{3} = 120$ tirages de $3$ boules équiprobables :
    \begin{itemize}
        \item $3$ boules de même couleur (soit $3$ blanches, soit $3$ noires, car il n'y a que $2$ rouges) : $P = \dfrac{\binom{5}{3} + \binom{3}{3}}{\binom{10}{3}} = \dfrac{10 + 1}{120} = \boxed{\dfrac{11}{120}}$.
        \item $1$ boule de chaque couleur ($1$ blanche, $1$ noire et $1$ rouge) : $P = \dfrac{\binom{5}{1}\binom{3}{1}\binom{2}{1}}{\binom{10}{3}} = \dfrac{30}{120} = \boxed{\dfrac{1}{4}}$.
    \end{itemize}
\end{enumerate}

% ----------------------------------------------------------------------------
\section{Corrigés du Chapitre 6 (Suites Arithmétiques et Géométriques)}
% ----------------------------------------------------------------------------

\paragraph{Corrigé de l'Exercice 6.1 (Suite arithmétique) :}
On a $u_{11} - u_4 = (11 - 4)r = 7r = 35 - 14 = 21 \implies \boxed{r = 3}$.
Puis $u_1 = u_4 - 3r = 14 - 9 = \boxed{5}$. Enfin, la somme des $20$ premiers termes vaut :
\[
    S_{20} = \frac{20}{2}\bigl(2(5) + 19(3)\bigr) = 10(10 + 57) = \boxed{670}.
\]

\paragraph{Corrigé de l'Exercice 6.2 (Suite géométrique et moyenne géométrique) :}
Notons les trois termes en progression géométrique $\frac{b}{q},\; b,\; bq$ (avec $q > 1$ car la suite est strictement croissante).
Leur produit donne immédiatement le terme central : $\left(\frac{b}{q}\right)\cdot b \cdot (bq) = b^3 = 216 = 6^3 \implies \mathbf{b = 6}$ !
La somme devient $\frac{6}{q} + 6 + 6q = 26 \iff 6q^2 - 20q + 6 = 0 \iff 3q^2 - 10q + 3 = (3q - 1)(q - 3) = 0$.
Comme $q > 1$, on retient $q = 3$, et les trois termes sont $\boxed{2,\;\; 6,\;\; 18}$.

\paragraph{Corrigé de l'Exercice 6.3 (Le puits de forage) :}
Le coût du $n$-ième mètre suit une suite arithmétique de premier terme $u_1 = 100$\,€ et de raison $r = 15$\,€. Le coût total pour forer $n$ mètres est :
\[
    S_n = \frac{n}{2}\bigl(200 + 15(n - 1)\bigr) \le 5200 \iff 15n^2 + 185n - 10\,400 \le 0 \iff 3n^2 + 37n - 2080 \le 0.
\]
Cherchons la racine positive : $\Delta = 37^2 + 4(3)(2080) = 1369 + 24\,960 = 26\,329$. Plus simplement, testons $n = 20$ : $3(400) + 37(20) = 1200 + 740 = 1940 \le 2080$, puis $n = 21$ : $3(441) + 37(21) = 1323 + 777 = 2100 > 2080$ !
On peut donc forer au maximum $\boxed{20\text{ mètres entiers}}$ (ou exactement $20,86\text{ m}$ si une fraction de mètre est facturée au prorata).

\paragraph{Corrigé de l'Exercice 6.4 (UCLouvain -- Juillet 2021, Q3 : Le robot-tondeuse) :}
\begin{enumerate}[label=(\alph*)]
    \item Au cycle $n$, l'énergie disponible pour la tonte est $C_b(n) - C_{b,m} = C_{b,0}(1 - d_b)^{n-1} - C_{b,m}$. Avec une puissance consommée $p_t$, la durée de tonte vaut :
    \[
        \boxed{t_t(n) = \frac{C_{b,0}(1 - d_b)^{n-1} - C_{b,m}}{p_t}}.
    \]
    Sur un cycle complet de durée totale $T(n) = t_t(n) + t_c + t_a(n)$, l'équilibre entre le volume d'herbe tondu et le volume qui a poussé s'écrit $v_t \cdot t_t(n) = v_p \cdot T(n)$. Au cycle limite $n_\ell$, le temps d'attente tombe à zéro ($t_a(n_\ell) = 0$, donc $T(n_\ell) = t_t(n_\ell) + t_c$), d'où :
    \[
        (v_t - v_p)\,t_t(n_\ell) = v_p\,t_c \iff (1 - d_b)^{n_\ell - 1} = \frac{\frac{v_p}{v_t - v_p}\,p_t\,t_c + C_{b,m}}{C_{b,0}}.
    \]
    En prenant le logarithme, on obtient la formule littérale exacte :
    \[
        \boxed{n_\ell = 1 + \frac{\ln\left(\dfrac{\frac{v_p}{v_t - v_p}\,p_t\,t_c + C_{b,m}}{C_{b,0}}\right)}{\ln(1 - d_b)}} \quad \left(\text{numériquement : } 1 + \frac{\ln(5,555/100)}{\ln(0,995)} \approx 449,6 \text{ cycles}\right).
    \]
    \item Le temps total sur les $n_\ell$ cycles est $T_{\text{tot}} = \sum_{n=1}^{n_\ell} T(n) = \frac{v_t}{v_p}\sum_{n=1}^{n_\ell} t_t(n)$. Par la formule de la somme d'une suite géométrique de raison $q = 1 - d_b$ (avec $1 - q = d_b$) :
    \[
        \boxed{T_{\text{tot}} = \frac{v_t}{v_p \cdot p_t}\left[C_{b,0}\,\frac{1 - (1 - d_b)^{n_\ell}}{d_b} - n_\ell\,C_{b,m}\right]}.
    \]
\end{enumerate}

% ----------------------------------------------------------------------------
\section{Corrigés du Chapitre 7 (Discussion d'Inéquations Paramétriques)}
% ----------------------------------------------------------------------------

\paragraph{Corrigé de l'Exercice 7.1 (UCLouvain -- Juillet 2021, Q2.2) :}
\begin{enumerate}
    \item \textbf{Étude du numérateur :} $N(x) = x^4 - 4x^2 + 4 = (x^2 - 2)^2 = \bigl((x - \sqrt{2})(x + \sqrt{2})\bigr)^2 \ge 0$.
    C'est un \textbf{carré parfait} ! Il est strictement positif pour tout $x \in \R \setminus \{-\sqrt{2}, \sqrt{2}\}$ et s'annule en $x = \pm\sqrt{2}$.
    Puisque l'inégalité demandée est \textbf{stricte} ($\frac{N(x)}{D(x)} < 0$), les valeurs $x = -\sqrt{2}$ et $x = \sqrt{2}$ (qui annulent le numérateur) ne peuvent \textbf{jamais} être solutions, et pour $x \neq \pm\sqrt{2}$ le quotient est strictement négatif si et seulement si le dénominateur est \textbf{strictement négatif} : $D(x) = x^2 - mx + 4 < 0$.
    \item \textbf{Étude du dénominateur $D(x) = x^2 - mx + 4$ (avec $m < 0$) :}
    Son discriminant vaut $\Delta = m^2 - 16 = (m - 4)(m + 4)$. Comme $m < 0$ par hypothèse :
    \begin{itemize}
        \item \textbf{Si $-4 \le m < 0$ ($\Delta \le 0$) :} $D(x) \ge 0$ sur tout $\R$ (jamais strictement négatif) $\implies \mathbf{S = \varnothing}$.
        \item \textbf{Si $m < -4$ ($\Delta > 0$) :} $D(x) < 0$ strictement entre ses deux racines $x_1 = \frac{m - \sqrt{m^2 - 16}}{2}$ et $x_2 = \frac{m + \sqrt{m^2 - 16}}{2}$.
        Comme $x_1 + x_2 = m < 0$ et $x_1 x_2 = 4 > 0$, les deux racines $x_1$ et $x_2$ sont \textbf{strictement négatives} ! La racine positive $\sqrt{2}$ du numérateur n'est donc jamais dans $]x_1, x_2[$. En revanche, la racine négative $-\sqrt{2}$ peut-elle tomber dans $]x_1, x_2[$ ?
        Calculons $D(-\sqrt{2}) = (-\sqrt{2})^2 - m(-\sqrt{2}) + 4 = 6 + m\sqrt{2}$.
        On a $D(-\sqrt{2}) < 0 \iff -\sqrt{2} \in \,]x_1, x_2[ \iff 6 + m\sqrt{2} < 0 \iff m < -\frac{6}{\sqrt{2}} = -3\sqrt{2} \approx -4,24$ (et pour $m = -3\sqrt{2}$, $x_2 = -\sqrt{2}$ est déjà exclu de l'intervalle ouvert $]x_1, x_2[$).
    \end{itemize}
\end{enumerate}
\textbf{Conclusion :}
\[
    \boxed{\begin{cases}
        \text{Si } -4 \le m < 0 : & S = \varnothing \\
        \text{Si } -3\sqrt{2} \le m < -4 : & S = \left]\dfrac{m - \sqrt{m^2 - 16}}{2},\, \dfrac{m + \sqrt{m^2 - 16}}{2}\right[ \\[2ex]
        \text{Si } m < -3\sqrt{2} : & S = \left]\dfrac{m - \sqrt{m^2 - 16}}{2},\, \dfrac{m + \sqrt{m^2 - 16}}{2}\right[ \setminus \{-\sqrt{2}\}
    \end{cases}}
\]

\paragraph{Corrigé de l'Exercice 7.2 (Inéquation rationnelle $\frac{x^2 - (m+1)x + m}{x^2 - 1} \le 0$) :}
\textbf{C.E. :} $x^2 - 1 \neq 0 \iff x \notin \{-1, 1\}$.
Factorisons le numérateur et le dénominateur : $x^2 - (m + 1)x + m = (x - 1)(x - m)$ (somme $m + 1$, produit $m$ !) et $x^2 - 1 = (x - 1)(x + 1)$.
Pour tout $x \notin \{-1, 1\}$, on peut simplifier par $(x - 1) \neq 0$ (tout en gardant la restriction $\mathbf{x \neq 1}$) :
\[
    \frac{(x - 1)(x - m)}{(x - 1)(x + 1)} \le 0 \iff \begin{cases} \dfrac{x - m}{x + 1} \le 0 \\ x \neq 1 \text{ et } x \neq -1 \end{cases}
\]
Le quotient $\frac{x - m}{x + 1}$ est négatif ou nul entre ses valeurs critiques $-1$ (exclue) et $m$ (incluse si $m \neq \pm 1$). On compare donc $m$ avec $-1$ et $1$ :
\[
    \boxed{\begin{cases}
        \text{Si } m < -1 : & S = [m, -1[ \\
        \text{Si } m = -1 : & S = \varnothing \quad (\text{car pour } x \neq \pm 1,\; \frac{x+1}{x+1} = 1 > 0) \\
        \text{Si } -1 < m < 1 : & S = \,]-1, m] \\
        \text{Si } m = 1 : & S = \,]-1, 1[ \quad (\text{car } x = 1 \text{ est une valeur interdite !}) \\
        \text{Si } m > 1 : & S = \,]-1, 1[ \,\cup\, ]1, m]
    \end{cases}}
\]

\paragraph{Corrigé de l'Exercice 7.3 (Inéquation logarithmique $\log_a(x^2 - ax) < 1$) :}
\textbf{Condition sur la base $a$ :} $a > 0$ et $a \neq 1$ (donc si $a \le 0$ ou $a = 1$, $\mathbf{S = \varnothing}$).
\begin{itemize}
    \item \textbf{Cas 1 : $a > 1$ (base croissante, ordre conservé).}
    $\log_a(x^2 - ax) < 1 = \log_a(a) \iff 0 < x^2 - ax < a \iff \begin{cases} x(x - a) > 0 & (1) \\ x^2 - ax - a < 0 & (2) \end{cases}$.
    La condition $(1)$ donne $x \in \,]-\infty, 0[ \,\cup\, ]a, +\infty[$.
    Le trinôme de $(2)$ a pour discriminant $\Delta = a^2 + 4a > 0$ (car $a > 1$), et pour racines $x_1 = \frac{a - \sqrt{a^2+4a}}{2} < 0$ et $x_2 = \frac{a + \sqrt{a^2+4a}}{2} > a$. L'intersection de $]x_1, x_2[$ avec $(1)$ donne :
    \[
        S = \left]\frac{a - \sqrt{a^2 + 4a}}{2},\, 0\right[ \cup \left]a,\, \frac{a + \sqrt{a^2 + 4a}}{2}\right[.
    \]
    \item \textbf{Cas 2 : $0 < a < 1$ (base décroissante, ordre inversé !).}
    $\log_a(x^2 - ax) < 1 \iff x^2 - ax > a \iff x^2 - ax - a > 0$ (ce qui garantit automatiquement la C.E. $x^2 - ax > a > 0$). Le trinôme est strictement positif à l'extérieur de ses racines $x_1 < x_2$ :
    \[
        S = \left]-\infty,\, \frac{a - \sqrt{a^2 + 4a}}{2}\right[ \cup \left]\frac{a + \sqrt{a^2 + 4a}}{2},\, +\infty\right[.
    \]
\end{itemize}

% ----------------------------------------------------------------------------
\section{Corrigés du Chapitre 8 (Mise en Équations et Systèmes Algébriques)}
% ----------------------------------------------------------------------------

\paragraph{Corrigé de l'Exercice 8.1 (UCLouvain -- Septembre 2021, Q3 : La collecte des déchets) :}
\begin{enumerate}
    \item \textbf{Relation entre la masse de déchets $Q_i$ et le nombre de levées $N_i$ :}
    Puisque les conteneurs ne sont sortis que lorsqu'ils sont intégralement remplis (capacité $5\text{ kg}$ pour Laurent et Nadine, $10\text{ kg}$ pour Bruno), on a :
    \[
        N_L = \frac{Q_L}{5}, \qquad N_N = \frac{Q_N}{5}, \qquad N_B = \frac{Q_B}{10}.
    \]
    \item \textbf{Mise en équations en fonction des seules inconnues $Q_L, Q_N, Q_B$ :}
    \begin{itemize}
        \item \emph{Masse totale de déchets :} $Q_L + Q_N + Q_B = 290$ \quad $(1)$.
        \item \emph{Nombre total de levées :} $N_L + N_N + N_B = \dfrac{Q_L}{5} + \dfrac{Q_N}{5} + \dfrac{Q_B}{10} = 43 \xrightarrow{\times 10} 2Q_L + 2Q_N + Q_B = 430$ \quad $(2)$.
        \item \emph{Facture totale ($251,60$\,€) :} Laurent ne dépasse pas son forfait (il paie uniquement $50$\,€). Nadine (forfait $50$\,€, franchises $50\text{ kg}$ et $12$ levées) et Bruno (forfait $80$\,€, franchises $70\text{ kg}$ et $12$ levées) paient les suppléments :
        \[
            50 + 50 + 80 + 0,4(Q_N - 50) + 1,2\left(\frac{Q_N}{5} - 12\right) + 0,4(Q_B - 70) + 1,2\left(\frac{Q_B}{10} - 12\right) = 251,6.
        \]
    \end{itemize}
    \item \textbf{Résolution :}
    De $(1)$, $Q_L + Q_N = 290 - Q_B$. En remplaçant dans $(2)$ : $2(290 - Q_B) + Q_B = 430 \iff 580 - Q_B = 430 \implies \boxed{Q_B = 150\text{ kg}}$ (et donc $Q_L + Q_N = 140\text{ kg}$).
    Dans l'équation de la facture, remplaçons $Q_B = 150$ (les suppléments de Bruno valent $0,4(80) + 1,2(15 - 12) = 32 + 3,6 = 35,60$\,€) :
    \[
        180 + 0,4(Q_N - 50) + 0,24\,Q_N - 14,4 + 35,6 = 251,6 \iff 0,64\,Q_N + 181,2 = 251,6,
    \]
    d'où $0,64\,Q_N = 70,4 \implies \boxed{Q_N = \frac{7040}{64} = 110\text{ kg}}$,
    et enfin $Q_L = 140 - 110 = \boxed{30\text{ kg}}$ (qui est bien $\le 50\text{ kg}$ et correspond à $6 \le 12$ levées : Laurent n'a effectivement pas dépassé son forfait !).
\end{enumerate}

\paragraph{Corrigé de l'Exercice 8.2 (Problème d'alliages métallurgiques) :}
Soient $x, y, z$ les masses (en kg) prélevées respectivement des lingots A, B et C. Les équations de conservation de la masse de cuivre, de zinc et d'étain (ou de la masse totale $x + y + z = 100$) s'écrivent :
\[
    \begin{cases}
        x + y + z = 100 & (1) \text{ [Masse totale]} \\
        0,6x + 0,5y + 0,4z = 52 \iff 6x + 5y + 4z = 520 & (2) \text{ [Cuivre]} \\
        0,1x + 0,3y + 0,1z = 17 \iff x + 3y + z = 170 & (3) \text{ [Étain]}
    \end{cases}
\]
\textbf{Astuce :} Regardez $(3) - (1)$ (l'étain a le même pourcentage $10\,\%$ dans A et C !) :
\[
    (x + 3y + z) - (x + y + z) = 170 - 100 \iff 2y = 70 \implies \boxed{y = 35\text{ kg du lingot B}} !
\]
On remplace $y = 35$ dans $(1)$ et $(2)$ : $x + z = 65$ et $6x + 4z = 520 - 175 = 345 \iff 6x + 4(65 - x) = 345 \iff 2x + 260 = 345 \iff 2x = 85 \implies \boxed{x = 42,5\text{ kg du lingot A}}$ et $\boxed{z = 22,5\text{ kg du lingot C}}$.

\paragraph{Corrigé de l'Exercice 8.3 (Diagramme de Venn à 3 ateliers) :}
Notons $x$ le nombre d'étudiants suivant uniquement Analyse et Trigonométrie ($x_{An,Tr} = x$). D'après l'énoncé, $x_{Al,An} = 2x$, $x_{Al,Tr} = 25$ et le centre des trois ateliers vaut $x_{Al,An,Tr} = 30$.
Par la formule du crible de Poincaré (ou en sommant les 7 régions du diagramme de Venn), sachant que les intersections deux à deux complètes valent $\text{card}(Al \cap An) = 2x + 30$, $\text{card}(An \cap Tr) = x + 30$ et $\text{card}(Al \cap Tr) = 25 + 30 = 55$ :
\[
    250 = (140 + 130 + 110) - \bigl((2x + 30) + (x + 30) + 55\bigr) + 30 \iff 250 = 380 - (3x + 115) + 30 = 295 - 3x,
\]
d'où $3x = 45 \implies \mathbf{x = 15}$.
Le nombre d'étudiants inscrits \emph{uniquement} à l'atelier d'Algèbre est alors le total du cercle d'Algèbre ($140$) diminué des trois régions d'intersection de ce cercle ($2x = 30$, $25$ et $30$) :
\[
    n_{\text{Algèbre seul}} = 140 - (30 + 25 + 30) = 140 - 85 = \boxed{55\text{ étudiants}}.
\]


\end{document}
